\documentclass[12pt,a4paper,openany]{book}

% ------------------------------------------------------------
% Encodage, langue et typographie
% ------------------------------------------------------------
\usepackage[utf8]{inputenc}
\usepackage[T1]{fontenc}
\usepackage[french]{babel}
\usepackage{newpxtext}
\usepackage[scaled=0.94]{sourcesanspro}
\usepackage{microtype}
\usepackage{setspace}
\setstretch{1.18}
\usepackage{geometry}
\geometry{a4paper,inner=31mm,outer=25mm,top=27mm,bottom=29mm,bindingoffset=4mm}
\setlength{\parindent}{0pt}
\setlength{\parskip}{0.45em}
\raggedbottom
\clubpenalty=10000
\widowpenalty=10000
\displaywidowpenalty=10000
\emergencystretch=2em
\setcounter{tocdepth}{1}
\setcounter{secnumdepth}{1}

% ------------------------------------------------------------
% Mathématiques et tableaux
% ------------------------------------------------------------
\usepackage{amsmath,amssymb,amsthm,mathtools}
\usepackage{newpxmath}
\usepackage{booktabs,tabularx,array,multicol,graphicx}
\usepackage{caption}
\captionsetup{font={small,sf},labelfont={bf,color=navy},textfont={color=slate},skip=6pt}
\usepackage{enumitem}
\setlist[itemize]{leftmargin=1.35em,itemsep=0.2em,topsep=0.25em}
\setlist[enumerate]{leftmargin=1.55em,itemsep=0.25em,topsep=0.25em}
\newcolumntype{Y}{>{\raggedright\arraybackslash}X}

% ------------------------------------------------------------
% Graphiques : palette moderne et discrète
% ------------------------------------------------------------
\usepackage{xcolor}
\definecolor{ink}{HTML}{202B36}
\definecolor{navy}{HTML}{203A5F}
\definecolor{teal}{HTML}{2E6F78}
\definecolor{sage}{HTML}{55777A}
\definecolor{gold}{HTML}{A98245}
\definecolor{terracotta}{HTML}{A45C50}
\definecolor{slate}{HTML}{5E6B76}
\definecolor{mistblue}{HTML}{F5F7F9}
\definecolor{mistteal}{HTML}{F3F7F6}
\definecolor{mistsage}{HTML}{F4F7F6}
\definecolor{mistgold}{HTML}{FAF8F2}
\definecolor{mistred}{HTML}{FBF5F3}
\definecolor{softgray}{HTML}{F7F7F6}
% Compatibilité avec les figures déjà écrites.
\colorlet{uqamblue}{navy}
\colorlet{deepblue}{ink}
\colorlet{softblue}{mistblue}
\colorlet{softgreen}{mistsage}
\colorlet{softyellow}{mistgold}
\colorlet{softred}{mistred}
\usepackage{tikz}
\usetikzlibrary{arrows.meta,calc,positioning,patterns,angles,quotes,decorations.pathreplacing,3d}
\usepackage{pgfplots}
\pgfplotsset{compat=1.18}

% ------------------------------------------------------------
% Repères pédagogiques : surfaces légères, accents cohérents
% ------------------------------------------------------------
\usepackage[most]{tcolorbox}
\tcbset{
  enhanced jigsaw,
  breakable,
  frame hidden,
  boxrule=0pt,
  arc=0.8pt,
  outer arc=0.8pt,
  left=10pt,
  right=8pt,
  top=6pt,
  bottom=6pt,
  before skip=0.8em,
  after skip=0.8em,
  fonttitle=\sffamily\bfseries\small,
  separator sign none,
  before title={\vphantom{É}},
  after title={\par\smallskip}
}

\newtcolorbox{questionbox}[1][]{
  title={Question directrice},
  colback=mistteal,
  coltitle=teal,
  borderline west={2.2pt}{0pt}{teal},
  fontupper=\itshape,
  #1
}
\newtcolorbox{intuitionbox}[1][]{
  title={Idée intuitive},
  colback=mistsage,
  coltitle=sage,
  borderline west={1.8pt}{0pt}{sage},
  #1
}
\newtcolorbox{visualbox}[1][]{
  title={Lecture visuelle},
  colback=white,
  coltitle=teal,
  borderline west={2.1pt}{0pt}{teal},
  borderline north={0.35pt}{0pt}{teal!30},
  borderline south={0.35pt}{0pt}{teal!30},
  #1
}
\newtcolorbox{connectionbox}[1][]{
  title={Lien entre les représentations},
  colback=mistgold,
  coltitle=gold!75!black,
  borderline west={2pt}{0pt}{gold},
  #1
}
\newtcolorbox{methodbox}[1][]{
  title={Méthode},
  colback=mistblue,
  coltitle=navy,
  borderline west={2.2pt}{0pt}{navy},
  #1
}
\newtcolorbox{warningbox}[1][]{
  title={Point de vigilance},
  colback=mistred,
  coltitle=terracotta,
  borderline west={2.4pt}{0pt}{terracotta},
  #1
}
\newtcolorbox{summarybox}[1][]{
  title={Synthèse},
  colback=mistgold,
  coltitle=gold!75!black,
  borderline west={2.2pt}{0pt}{gold},
  #1
}
\newtcolorbox{examplebox}[1][]{
  title={Exemple},
  colback=softgray,
  coltitle=slate,
  borderline west={1.5pt}{0pt}{slate},
  #1
}
\newtcolorbox{counterbox}[1][]{
  title={Contre-exemple},
  colback=mistred,
  coltitle=terracotta,
  borderline west={2pt}{0pt}{terracotta},
  #1
}

\newtcbtheorem[number within=chapter]{definition}{Définition}{
  colback=mistteal,
  coltitle=teal,
  borderline west={2pt}{0pt}{teal},
  fonttitle=\sffamily\bfseries
}{def}
\newtcbtheorem[use counter from=definition]{theorem}{Théorème}{
  colback=mistblue,
  coltitle=navy,
  borderline west={2.4pt}{0pt}{navy},
  fonttitle=\sffamily\bfseries
}{thm}
\newtcbtheorem[use counter from=definition]{proposition}{Proposition}{
  colback=mistgold,
  coltitle=gold!75!black,
  borderline west={2pt}{0pt}{gold},
  fonttitle=\sffamily\bfseries
}{prop}

% ------------------------------------------------------------
% Titres, en-têtes et hyperliens
% ------------------------------------------------------------
\usepackage{titlesec}
\titleformat{\part}[display]
  {\centering\sffamily\bfseries\color{navy}}
  {\color{gold}\fontsize{44}{46}\selectfont\thepart}
  {0.6ex}
  {\Huge}
  [\vspace{1.1ex}{\color{teal}\titlerule}]
\titlespacing*{\part}{0pt}{1.5cm}{2.2cm}
\titleformat{\chapter}[display]
  {\sffamily\bfseries\color{navy}}
  {\raggedleft\color{teal}\large Chapitre \thechapter}
  {0.45ex}
  {\color{navy}\huge}
  [\vspace{0.7ex}{\color{gold}\titlerule}]
\titlespacing*{\chapter}{0pt}{-10pt}{25pt}
\titleformat{\section}
  {\Large\sffamily\bfseries\color{navy}}
  {\color{teal}\thesection}{0.7em}{}
\titleformat{\subsection}
  {\large\sffamily\bfseries\color{ink}}
  {\color{teal}\thesubsection}{0.7em}{}

\usepackage{fancyhdr}
\pagestyle{fancy}
\fancyhf{}
\fancyhead[LE]{\sffamily\small\color{slate}\thepage\quad\nouppercase{\leftmark}}
\fancyhead[RO]{\sffamily\small\color{slate}\nouppercase{\rightmark}\quad\thepage}
\renewcommand{\headrulewidth}{0.45pt}
\renewcommand{\headrule}{\hbox to\headwidth{\color{teal}\leaders\hrule height \headrulewidth\hfill}}
\setlength{\headheight}{26pt}

\usepackage[hidelinks,bookmarksopen=true,pdfstartview=FitH]{hyperref}
\usepackage{bookmark}
\hypersetup{
  pdftitle={MAT0339 - Mathématiques, manuel de cours},
  pdfauthor={Xia Xiao},
  pdfsubject={Manuel structuré et illustré du cours MAT0339},
  pdfkeywords={MAT0339, mathématiques, algèbre, fonctions, probabilités, vecteurs, matrices, trigonométrie}
}

% ------------------------------------------------------------
% Raccourcis
% ------------------------------------------------------------
\newcommand{\R}{\mathbb{R}}
\newcommand{\N}{\mathbb{N}}
\newcommand{\Z}{\mathbb{Z}}
\newcommand{\Q}{\mathbb{Q}}
\newcommand{\C}{\mathbb{C}}
\newcommand{\vect}[1]{\overrightarrow{#1}}
\newcommand{\norm}[1]{\left\lVert #1\right\rVert}
\newcommand{\abs}[1]{\left\lvert #1\right\rvert}
\newcommand{\dom}{\operatorname{Dom}}
\newcommand{\im}{\operatorname{Im}}
\newcommand{\e}{\mathrm{e}}
\newcommand{\dd}{\,\mathrm{d}}

\begin{document}
\color{ink}
\frontmatter

\begin{titlepage}
\centering
\vspace*{2.2cm}
{\Huge\sffamily\bfseries\color{navy} MAT0339\par}
\vspace{0.35cm}
{\LARGE\sffamily\color{ink} Mathématiques\par}
\vspace{0.8cm}
{\large\sffamily\color{teal} Manuel de cours structuré et illustré\par}
\vspace{0.35cm}
{\normalsize Algèbre \textbullet\ Fonctions \textbullet\ Probabilités \textbullet\ Géométrie analytique \textbullet\ Trigonométrie\par}
\vspace{1.5cm}
\begin{tikzpicture}[scale=1.0]
  \draw[thick,slate,-{Stealth[length=3mm]}] (-4,0)--(4,0);
  \draw[thick,slate,-{Stealth[length=3mm]}] (0,-1.7)--(0,2.5);
  \draw[navy,very thick,domain=-3.4:3.4,samples=100] plot(\x,{0.18*(\x)^2-0.8});
  \draw[teal,very thick,dashed,domain=-3.5:3.5,samples=100] plot(\x,{1.0+0.55*sin(70*\x)});
  \fill[gold] (1.6,0.8) circle (2.2pt);
  \draw[terracotta,very thick,-{Stealth[length=3mm]}] (-2.8,-1.1)--(-1.1,1.6);
\end{tikzpicture}
\vfill
{\Large\sffamily\bfseries\color{navy} Xia Xiao\par}
\vspace{0.25cm}
{\large\sffamily\color{slate} Département de mathématiques, UQAM\par}
\vspace{0.8cm}
{\large\sffamily\color{teal} Édition révisée -- août 2026\par}
\end{titlepage}

\thispagestyle{empty}
\vspace*{\fill}
\begin{center}
Ce manuel est mis gracieusement à la disposition des étudiantes et étudiants.\\
Toute vente ou utilisation commerciale est strictement interdite.
\end{center}
\vspace{0.6cm}
\textbf{Avis de non-responsabilité.} Bien que tout ait été mis en œuvre pour assurer l'exactitude du contenu, l'auteur ne peut être tenu responsable d'erreurs ou d'omissions éventuelles.
\vspace{0.6cm}
\begin{center}\textit{Dernière mise à jour : \today}\end{center}
\vspace*{\fill}
\clearpage

\chapter*{Préface}
\addcontentsline{toc}{chapter}{Préface}
Bienvenue dans le cours MAT0339.

Ce manuel a été reconstruit autour d'une idée simple : une notion mathématique devient plus claire lorsqu'on peut la voir sous plusieurs formes. Une droite est à la fois un graphique, une équation et un modèle de variation constante. Une matrice est à la fois un tableau de nombres, une transformation et un outil de résolution. Une fonction sinus est à la fois une hauteur sur un cercle et un modèle de phénomène périodique.

Le texte suit donc un mouvement régulier : partir d'une question concrète, construire une intuition, introduire la notation, puis appliquer la notion dans un exemple complet. Les exercices ont été retirés afin que le manuel serve de référence de cours, de lecture préparatoire et de soutien aux vidéos. Les exemples ne sont pas des exercices cachés : ils font partie intégrante de l'explication.

À l'ère de l'intelligence artificielle, il est facile d'obtenir un calcul. Il reste toutefois essentiel de savoir formuler le problème, choisir le bon modèle, vérifier les hypothèses et interpréter le résultat. C'est cette autonomie de raisonnement que le cours cherche à développer.

\begin{summarybox}
Le manuel couvre et approfondit l'ensemble du programme : nombres réels, modèles linéaires et quadratiques, fonctions polynomiales, rationnelles, exponentielles et logarithmiques, dénombrement et probabilités, vecteurs, droites et plans, matrices et systèmes linéaires, programmation linéaire, trigonométrie.
\end{summarybox}

\chapter*{Comment utiliser ce manuel}
\addcontentsline{toc}{chapter}{Comment utiliser ce manuel}
Chaque chapitre commence par une \textbf{question directrice}. Elle indique ce que l'on cherche à comprendre avant de présenter les formules. Des repères typographiques sobres distinguent les questions directrices, les idées intuitives, les méthodes, les avertissements et les synthèses. La couleur est utilisée avec parcimonie pour signaler la fonction d'un passage; la hiérarchie reste entièrement lisible en impression monochrome.

Les chapitres sont organisés en cinq parties :
\begin{enumerate}
  \item nombres et algèbre;
  \item fonctions et modèles;
  \item dénombrement et probabilités;
  \item vecteurs, géométrie analytique, matrices et optimisation;
  \item trigonométrie et phénomènes périodiques.
\end{enumerate}

Les notions sont cumulatives. Il est donc préférable de lire les chapitres dans l'ordre, mais chaque chapitre peut aussi servir de référence autonome.

Le manuel ne cherche pas à multiplier les résultats. Lorsqu'une formule importante apparaît, on privilégie quatre questions : \emph{que représente-t-elle?}, \emph{pourquoi a-t-elle cette forme?}, \emph{comment la voir sur une figure?} et \emph{comment vérifier qu'on l'utilise correctement?} Les thèmes secondaires sont condensés lorsqu'ils n'améliorent pas la compréhension; l'espace ainsi libéré est consacré aux idées structurantes, aux schémas et aux changements de représentation.

\tableofcontents
\mainmatter

% ============================================================
\part{Nombres réels et algèbre}
% ============================================================

\chapter{Les nombres réels et leurs opérations}

\begin{questionbox}
Comment organiser les différents nombres et calculer avec eux sans perdre le sens des opérations?
\end{questionbox}

\section{Les grandes familles de nombres}

\begin{intuitionbox}
Les ensembles de nombres se sont agrandis chaque fois qu'une opération devenait impossible. Les naturels permettent de compter; les entiers permettent de soustraire; les rationnels permettent de diviser; les réels remplissent toute la droite numérique.
\end{intuitionbox}

\begin{definition}{Ensembles usuels de nombres}{}
\begin{align*}
\N&=\{0,1,2,3,\ldots\},\\
\Z&=\{\ldots,-2,-1,0,1,2,\ldots\},\\
\Q&=\left\{\frac ab:a,b\in\Z,\ b\neq0\right\},\\
\R&=\Q\cup\{\text{nombres irrationnels}\}.
\end{align*}
On a les inclusions
\[
\N\subset\Z\subset\Q\subset\R.
\]
\end{definition}

\begin{center}
\begin{tikzpicture}
\draw[rounded corners,fill=softblue,draw=uqamblue,thick] (-4,-2.2) rectangle (4,2.2);
\draw[rounded corners,fill=white,draw=teal,thick] (-3,-1.55) rectangle (3,1.55);
\draw[rounded corners,fill=softyellow,draw=gold,thick] (-2.2,-0.95) rectangle (2.2,0.95);
\draw[rounded corners,fill=white,draw=deepblue,thick] (-1.3,-0.45) rectangle (1.3,0.45);
\node[anchor=north west] at (-3.8,2.0) {$\R$};
\node[anchor=north west] at (-2.8,1.35) {$\Q$};
\node[anchor=north west] at (-2.0,0.75) {$\Z$};
\node at (0,0) {$\N$};
\node at (-3.2,-1.75) {$\sqrt2,\ \pi$};
\node at (-2.2,-1.15) {$\frac23,\ -0{,}4$};
\node at (-1.5,-0.65) {$-7$};
\node at (0.7,0) {$0,1,2,\ldots$};
\end{tikzpicture}
\end{center}

\begin{examplebox}
Le nombre $-5$ appartient à $\Z$, donc aussi à $\Q$ et à $\R$. Le nombre $0{,}125=\frac18$ est rationnel. Le nombre $\sqrt2$ est réel mais n'est pas rationnel.
\end{examplebox}

\subsection{Intervalles et valeur absolue}

Un intervalle décrit une partie continue de la droite réelle. Les crochets indiquent si les extrémités sont incluses :
\[
[a,b]=\{x\in\R:a\le x\le b\},\qquad ]a,b[=\{x\in\R:a<x<b\}.
\]
L'infini n'est jamais inclus : on écrit $[2,+\infty[$ et non $[2,+\infty]$.

\begin{definition}{Valeur absolue}{}
La valeur absolue $\abs{x}$ est la distance entre $x$ et $0$ sur la droite réelle :
\[
\abs{x}=\begin{cases}x,&x\ge0,\\-x,&x<0.\end{cases}
\]
Plus généralement, $\abs{x-a}$ est la distance entre $x$ et $a$.
\end{definition}

\begin{examplebox}
L'inégalité $\abs{x-3}\le2$ signifie que la distance entre $x$ et $3$ est au plus $2$. Elle équivaut à
\[
1\le x\le5,
\]
donc l'ensemble solution est $[1,5]$.
\end{examplebox}

\section{Priorité des opérations et signes}

\begin{methodbox}
Pour évaluer une expression :
\begin{enumerate}
\item traiter les parenthèses;
\item calculer les puissances et les racines;
\item effectuer multiplications et divisions de gauche à droite;
\item effectuer additions et soustractions de gauche à droite.
\end{enumerate}
\end{methodbox}

\begin{examplebox}
\[
3-2(5-8)^2=3-2(-3)^2=3-2\cdot9=-15.
\]
La puissance s'applique à toute la parenthèse. En revanche,
\[
-3^2=-(3^2)=-9,\qquad (-3)^2=9.
\]
\end{examplebox}

\begin{warningbox}
Le signe moins placé devant une puissance n'est pas automatiquement dans la base. Les parenthèses déterminent la base : $(-a)^2=a^2$, mais $-a^2=-(a^2)$.
\end{warningbox}

\section{Fractions et nombres décimaux}

\begin{intuitionbox}
Une fraction représente une division. Le dénominateur indique la taille des parts et le numérateur indique combien de parts sont prises.
\end{intuitionbox}

Pour $b\neq0$ et $d\neq0$ :
\begin{align*}
\frac ab+\frac cd&=\frac{ad+bc}{bd},\\
\frac ab\cdot\frac cd&=\frac{ac}{bd},\\
\frac ab\div\frac cd&=\frac ab\cdot\frac dc\qquad(c\neq0).
\end{align*}

\begin{examplebox}
\[
\frac34-\frac25=\frac{15}{20}-\frac8{20}=\frac7{20}.
\]
Pour la division,
\[
\frac56\div\frac{10}{9}=\frac56\cdot\frac9{10}=\frac34.
\]
\end{examplebox}

Un nombre rationnel possède une écriture décimale finie ou périodique. Par exemple,
\[
0{,}\overline{27}=\frac{27}{99}=\frac3{11}.
\]
La méthode générale consiste à multiplier par une puissance de $10$ puis à soustraire afin d'éliminer la partie périodique.

\section{Puissances, notation scientifique et racines}

\begin{definition}{Puissance entière}{}
Pour $n\in\N$,
\[
a^n=\underbrace{a\cdot a\cdots a}_{n\text{ facteurs}},\qquad a^0=1\quad(a\neq0).
\]
Pour $n>0$, on définit $a^{-n}=1/a^n$.
\end{definition}

Les lois fondamentales sont
\begin{align*}
a^m a^n&=a^{m+n},&\frac{a^m}{a^n}&=a^{m-n},\\
(a^m)^n&=a^{mn},&(ab)^n&=a^n b^n.
\end{align*}
Elles sont valides lorsque les expressions sont définies.

\begin{examplebox}
\[
\frac{2^7\cdot2^{-3}}{2^2}=2^{7-3-2}=2^2=4.
\]
La notation scientifique écrit un nombre sous la forme $a\times10^n$ avec $1\le\abs a<10$. Ainsi,
\[
0{,}000047=4{,}7\times10^{-5}.
\]
\end{examplebox}

\begin{definition}{Racine carrée}{}
Pour $a\ge0$, $\sqrt a$ désigne l'unique nombre non négatif dont le carré vaut $a$.
\end{definition}

\begin{warningbox}
On a toujours $\sqrt{x^2}=\abs{x}$, et non simplement $x$. Par exemple, si $x=-4$, alors $\sqrt{x^2}=4$.
\end{warningbox}

\begin{examplebox}
\[
\sqrt{72}=\sqrt{36\cdot2}=6\sqrt2,
\qquad
\sqrt{12}\sqrt3=\sqrt{36}=6.
\]
Pour $a,b\ge0$, $\sqrt{ab}=\sqrt a\sqrt b$. Cette formule n'est pas une règle générale sur les nombres négatifs dans $\R$.
\end{examplebox}


\section{La droite réelle comme modèle unificateur}

\begin{visualbox}
La droite réelle transforme trois idées abstraites en gestes visibles : \textbf{comparer} revient à regarder qui est le plus à droite, \textbf{soustraire} revient à mesurer un déplacement orienté, et \textbf{prendre une valeur absolue} revient à oublier le sens pour ne garder que la distance.
\end{visualbox}

\begin{center}
\begin{tikzpicture}[x=1.15cm,y=1cm]
  \draw[-{Stealth[length=3mm]},thick,slate] (-4.5,0)--(4.8,0) node[right] {$x$};
  \foreach \x in {-4,-3,...,4}{\draw[slate] (\x,0.09)--(\x,-0.09) node[below=3pt] {\small $\x$};}
  \fill[terracotta] (-2,0) circle (2.4pt) node[above=5pt] {$a=-2$};
  \fill[teal] (3,0) circle (2.4pt) node[above=5pt] {$b=3$};
  \draw[gold,very thick,{Stealth[length=2.5mm]}-{Stealth[length=2.5mm]}] (-2,0.55)--(3,0.55);
  \node[gold!75!black,fill=white,inner sep=2pt] at (0.5,0.55) {$|b-a|=5$};
  \draw[navy,thick,decorate,decoration={brace,amplitude=5pt,mirror}] (0,-0.55)--(3,-0.55)
    node[midway,below=7pt] {$|3|=3$};
\end{tikzpicture}
\end{center}

\begin{connectionbox}
Les intervalles sont des morceaux de cette droite. Une inégalité comme $-2\le x<3$ et l'intervalle $[-2,3[$ décrivent exactement le même ensemble de points. Le passage entre langage, symbole et figure est plus important que la mémorisation des crochets.
\end{connectionbox}

\subsection{Estimer avant de calculer}
Avant tout calcul numérique, une estimation donne un ordre de grandeur. Par exemple,
\[
\frac{19{,}8\times 4{,}9}{0{,}51}\approx\frac{20\times5}{0{,}5}=200.
\]
Si la calculatrice affiche $19{,}02$, l'erreur est immédiatement visible. Une valeur exacte, comme $3\sqrt2$, conserve toute l'information; une valeur approchée, comme $4{,}24$, sert à interpréter ou mesurer.

\subsection{Les unités comme contrôle logique}
Une équation n'est cohérente que si les deux membres ont les mêmes unités. On peut additionner des mètres à des mètres, mais pas des mètres à des secondes. Cette règle simple détecte de nombreuses erreurs de modèle avant même le calcul.


\section{Cohérence des règles de calcul}

Les règles de calcul ne sont pas une liste arbitraire. Elles prolongent des idées simples déjà visibles avec les nombres positifs.

\subsection{Le signe d'un produit}
Multiplier par un entier positif signifie additionner plusieurs fois. Ainsi, $3\times(-2)$ représente trois déplacements de deux unités vers la gauche, donc $-6$. Pour comprendre $(-3)\times(-2)$, on demande que la distributivité reste vraie :
\[
0=(-3)\times0=(-3)\times(2+(-2))=(-3)\times2+(-3)\times(-2).
\]
Comme $(-3)\times2=-6$, le second terme doit valoir $6$. La règle « moins par moins donne plus » est donc imposée par la cohérence de l'algèbre.

\begin{center}
\begin{tikzpicture}[x=0.8cm,y=1cm]
 \draw[-{Stealth},thick,slate] (-7,0)--(7.5,0);
 \foreach \x in {-6,-4,-2,0,2,4,6}{\draw (\x,0.1)--(\x,-0.1) node[below] {\small $\x$};}
 \draw[navy,very thick,-{Stealth[length=3mm]}] (0,0.55)--(-2,0.55);
 \draw[navy,very thick,-{Stealth[length=3mm]}] (-2,0.55)--(-4,0.55);
 \draw[navy,very thick,-{Stealth[length=3mm]}] (-4,0.55)--(-6,0.55);
 \node[navy,above] at (-3,0.75) {$3\times(-2)=-6$};
 \draw[teal,very thick,-{Stealth[length=3mm]}] (-6,-0.65)--(0,-0.65);
 \node[teal,below] at (-3,-0.85) {changer le signe de $3$ inverse le déplacement};
\end{tikzpicture}
\end{center}

\subsection{Les puissances et les racines sont des opérations inverses}
La notation $a^n$ condense une multiplication répétée. La racine $\sqrt[n]{a}$ cherche au contraire le nombre dont la puissance $n$ vaut $a$. Cette réciprocité explique les exposants fractionnaires :
\[
a^{1/n}=\sqrt[n]{a},
\qquad
(a^{1/n})^n=a.
\]
Il faut toutefois distinguer une opération inverse d'une égalité sans condition. Sur les réels, $\sqrt{x^2}=|x|$, car la racine carrée désigne toujours la valeur non négative.

\begin{visualbox}
Les trois niveaux de contrôle d'un calcul numérique sont :
\begin{enumerate}
\item le \textbf{signe} attendu;
\item l'\textbf{ordre de grandeur};
\item l'\textbf{unité} du résultat.
\end{enumerate}
Un résultat qui échoue à l'un de ces contrôles doit être réexaminé avant toute interprétation.
\end{visualbox}


\section{Mesurer, estimer et contrôler un résultat}

Les règles de calcul deviennent plus faciles à retenir lorsqu'on comprend ce qu'elles protègent. Elles empêchent qu'une même expression reçoive deux valeurs différentes et permettent de comparer des résultats obtenus par des chemins distincts. Cette section reprend les notions du chapitre sous trois points de vue complémentaires : la droite réelle, la mesure et le contrôle numérique.

\subsection{La droite réelle comme modèle de comparaison}

Tout nombre réel correspond à un point de la droite. L'ordre $a<b$ signifie que le point représentant $a$ est situé à gauche de celui qui représente $b$. Cette image rend immédiates plusieurs propriétés :
\begin{itemize}
\item ajouter le même nombre aux deux membres déplace les deux points de la même distance et conserve l'ordre;
\item multiplier par un nombre positif étire ou contracte la droite sans changer son orientation;
\item multiplier par un nombre négatif ajoute une réflexion et renverse donc l'ordre.
\end{itemize}

Ainsi,
\[
-4<1
\]
devient, après multiplication par $-3$,
\[
12>-3.
\]
Le changement du sens de l'inégalité n'est pas une règle arbitraire : il traduit la réflexion de la droite réelle.

\begin{examplebox}
Comparons $\sqrt{50}$ et $7$. Comme les deux nombres sont positifs, on peut comparer leurs carrés :
\[
(\sqrt{50})^2=50,
\qquad
7^2=49.
\]
Donc $\sqrt{50}>7$. Cette méthode évite d'utiliser une approximation décimale prématurée.
\end{examplebox}

\subsection{Valeur absolue, distance et voisinage}

La valeur absolue mesure une distance :
\[
\abs{x-a}
\]
est la distance entre $x$ et $a$. Cette interprétation transforme des inégalités en descriptions géométriques. Par exemple,
\[
\abs{x-3}\leq2
\]
signifie que $x$ est à une distance d'au plus $2$ de $3$. On obtient donc
\[
1\leq x\leq5.
\]
De même,
\[
\abs{x+1}>4
\]
signifie que $x$ est à plus de $4$ unités de $-1$ :
\[
x<-5\quad\text{ou}\quad x>3.
\]

Cette lecture est plus robuste que la mémorisation de deux modèles de résolution, car elle reste valide dans les problèmes de tolérance, d'erreur de mesure et d'intervalle acceptable.

\subsection{Valeur exacte et valeur approchée}

Une valeur exacte conserve toute l'information mathématique. Les écritures
\[
\frac{7}{12},\qquad \sqrt{13},\qquad \frac{3\pi}{5}
\]
sont exactes. Une écriture décimale obtenue à la calculatrice est souvent une approximation :
\[
\sqrt{13}\approx3{,}606.
\]
Le symbole $=$ ne doit pas remplacer $\approx$. Une approximation est utile pour interpréter un résultat, mais elle ne doit pas être introduite trop tôt dans une chaîne de calcul, car les erreurs d'arrondi s'accumulent.

\begin{methodbox}
Pour un calcul en plusieurs étapes :
\begin{enumerate}
\item conserver les fractions, les racines et $\pi$ aussi longtemps que possible;
\item effectuer les simplifications exactes;
\item convertir en décimal seulement à la fin;
\item annoncer la précision et l'unité.
\end{enumerate}
\end{methodbox}

\subsection{Rapports, proportions et pourcentages}

Un pourcentage est un rapport à $100$. Une hausse de $r\%$ correspond à une multiplication par
\[
1+\frac r{100},
\]
et une baisse de $r\%$ à une multiplication par
\[
1-\frac r{100}.
\]
Deux variations successives ne s'additionnent généralement pas. Une hausse de $20\%$ suivie d'une baisse de $20\%$ transforme une valeur $V$ en
\[
V\times1{,}20\times0{,}80=0{,}96V.
\]
La valeur finale est donc inférieure de $4\%$ à la valeur initiale. La raison est que les deux pourcentages ne sont pas calculés sur la même base.

\begin{examplebox}
Un prix de $240\,$\$ reçoit une réduction de $15\%$, puis une taxe de $14{,}975\%$. Le prix final est
\[
240(0{,}85)(1{,}14975)\approx234{,}55\ \$.
\]
Soustraire $15$ puis ajouter $14{,}975$ serait incorrect : les pourcentages sont des facteurs multiplicatifs, non des montants fixes.
\end{examplebox}

\subsection{Unités et cohérence dimensionnelle}

Les unités font partie du raisonnement. Si une distance est donnée en kilomètres et une vitesse en kilomètres par heure, le temps obtenu par
\[
t=\frac d v
\]
est mesuré en heures. Une égalité comme
\[
3\ \mathrm{m}+4\ \mathrm{s}=7
\]
n'a pas de sens, même si l'addition des nombres $3$ et $4$ est possible.

La cohérence des unités est aussi un outil de détection d'erreur. Une aire doit être exprimée en unités carrées, un volume en unités cubiques et une vitesse comme un rapport distance/temps.

\subsection{Contrôles numériques essentiels}

Avant d'accepter un résultat, quatre contrôles rapides sont utiles :
\begin{enumerate}
\item \textbf{signe :} le résultat devrait-il être positif, négatif ou nul?
\item \textbf{ordre de grandeur :} $0{,}49\times200$ doit être voisin de $100$, non de $1000$;
\item \textbf{unité :} l'unité finale correspond-elle à la grandeur cherchée?
\item \textbf{encadrement :} si $9<11<16$, alors $3<\sqrt{11}<4$.
\end{enumerate}

\begin{summarybox}
Les nombres réels servent à comparer et à mesurer. La valeur absolue exprime une distance, les pourcentages sont multiplicatifs, les unités contrôlent la cohérence et les approximations doivent être annoncées comme telles.
\end{summarybox}


\chapter{Expressions algébriques, équations et inéquations}

\begin{questionbox}
Comment transformer une expression sans changer sa valeur, et comment isoler une inconnue sans perdre ni créer de solutions?
\end{questionbox}

\section{Variables, monômes et polynômes}

\begin{definition}{Polynôme}{}
Un polynôme en $x$ est une somme finie de termes $a_kx^k$ où les exposants $k$ sont des entiers naturels :
\[
P(x)=a_nx^n+\cdots+a_1x+a_0.
\]
Le plus grand exposant associé à un coefficient non nul est le degré du polynôme.
\end{definition}

\begin{examplebox}
Pour $P(x)=3x^3-2x+5$, le degré est $3$ et
\[
P(-2)=3(-2)^3-2(-2)+5=-15.
\]
\end{examplebox}

On additionne les termes semblables et on développe avec la distributivité :
\[
a(b+c)=ab+ac.
\]
Les identités remarquables les plus utiles sont
\begin{align*}
(a+b)^2&=a^2+2ab+b^2,\\
(a-b)^2&=a^2-2ab+b^2,\\
a^2-b^2&=(a-b)(a+b).
\end{align*}

\section{Factorisation et fractions algébriques}

\begin{intuitionbox}
Développer transforme un produit en somme. Factoriser fait le mouvement inverse : on cherche la structure multiplicative cachée dans une somme.
\end{intuitionbox}

\begin{methodbox}
Pour factoriser, chercher dans cet ordre :
\begin{enumerate}
\item un facteur commun;
\item une identité remarquable;
\item un trinôme factorisable;
\item un regroupement de termes.
\end{enumerate}
\end{methodbox}

\begin{examplebox}
\[
6x^3-9x^2=3x^2(2x-3),
\qquad
x^2-9=(x-3)(x+3).
\]
Pour $x^2+5x+6$, on cherche deux nombres dont la somme vaut $5$ et le produit vaut $6$ :
\[
x^2+5x+6=(x+2)(x+3).
\]
\end{examplebox}

Une fraction algébrique est un quotient de polynômes. Les valeurs annulant le dénominateur sont interdites, même si un facteur se simplifie.

\begin{examplebox}
\[
\frac{x^2-1}{x-1}=\frac{(x-1)(x+1)}{x-1}=x+1,
\qquad x\neq1.
\]
La formule simplifiée est $x+1$, mais la fonction initiale possède un trou en $x=1$.
\end{examplebox}

\begin{counterbox}
On ne peut pas simplifier des termes séparés dans une somme :
\[
\frac{x+2}{x}\neq 1+2.
\]
La simplification porte sur des \emph{facteurs}, pas sur des termes reliés par une addition.
\end{counterbox}

\section{Équations du premier degré}

\begin{definition}{Équation équivalente}{}
Deux équations sont équivalentes si elles ont le même ensemble de solutions. Ajouter le même nombre aux deux membres, soustraire le même nombre ou multiplier les deux membres par un nombre non nul conserve l'équivalence.
\end{definition}

\begin{methodbox}
Pour résoudre une équation linéaire :
\begin{enumerate}
\item développer et réduire;
\item regrouper les termes contenant l'inconnue;
\item regrouper les constantes de l'autre côté;
\item diviser par le coefficient non nul de l'inconnue;
\item vérifier dans l'équation initiale.
\end{enumerate}
\end{methodbox}

\begin{examplebox}
\begin{align*}
3(2x-1)-4&=2(x+5),\\
6x-7&=2x+10,\\
4x&=17,\\
x&=\frac{17}{4}.
\end{align*}
\end{examplebox}

Une équation peut aussi avoir aucune solution ou une infinité de solutions. Par exemple,
\[
2(x+1)=2x+5
\]
conduit à $2=5$ : aucune solution. En revanche,
\[
2(x+1)=2x+2
\]
conduit à une identité : tout réel est solution.

\section{Inéquations et tableaux de signes}

\begin{warningbox}
Lorsqu'on multiplie ou divise les deux membres d'une inégalité par un nombre négatif, le sens de l'inégalité s'inverse.
\end{warningbox}

\begin{examplebox}
\[
-3x+2\le11
\quad\Longrightarrow\quad
-3x\le9
\quad\Longrightarrow\quad
x\ge-3.
\]
\end{examplebox}

Pour une inéquation factorisée, un tableau de signes indique le signe de chaque facteur entre ses zéros.

\begin{examplebox}
Résolvons $(x-2)(x+1)\le0$. Les changements de signe peuvent se produire en $x=-1$ et $x=2$.
\[
\begin{array}{c|ccccc}
x&(-\infty,-1)&-1&(-1,2)&2&(2,+\infty)\\ \hline
x+1&-&0&+&+&+\\
x-2&-&-&-&0&+\\ \hline
(x-2)(x+1)&+&0&-&0&+
\end{array}
\]
La solution est $[-1,2]$.
\end{examplebox}

\section{Équations avec racines ou valeurs absolues}

Pour $b\ge0$,
\[
\abs{x-a}=b\quad\Longleftrightarrow\quad x-a=b\ \text{ou}\ x-a=-b.
\]

\begin{examplebox}
\[
\abs{2x-3}=5
\]
donne $2x-3=5$ ou $2x-3=-5$, donc $x=4$ ou $x=-1$.
\end{examplebox}

Élever au carré peut créer des solutions parasites, car l'implication
\[
u=v\Longrightarrow u^2=v^2
\]
n'est pas réversible en général.

\begin{examplebox}
Résolvons $\sqrt{x+1}=x-1$. Le domaine impose $x\ge1$. En élevant au carré,
\[
x+1=(x-1)^2=x^2-2x+1,
\]
donc $x(x-3)=0$. Les candidats sont $0$ et $3$, mais seul $3$ appartient au domaine et vérifie l'équation initiale.
\end{examplebox}


\section{Équivalences et transformations à contrôler}

\begin{visualbox}
Résoudre une équation ressemble à maintenir une balance en équilibre. Ajouter, soustraire, multiplier ou diviser par une même quantité non nulle des deux côtés conserve l'équilibre. En revanche, élever au carré ou multiplier par une expression pouvant être nulle peut créer de nouvelles possibilités : ces étapes ne sont pas automatiquement réversibles.
\end{visualbox}

\begin{center}
\begin{tikzpicture}[scale=0.9]
  \draw[navy,very thick] (-3.2,0)--(3.2,0);
  \draw[navy,very thick] (0,0)--(0,-1.2);
  \draw[navy,very thick] (-0.7,-1.2)--(0.7,-1.2);
  \draw[slate,thick] (-2.3,0)--(-2.3,-0.65) --(-3.1,-1.05)--(-1.5,-1.05)--cycle;
  \draw[slate,thick] (2.3,0)--(2.3,-0.65) --(1.5,-1.05)--(3.1,-1.05)--cycle;
  \node at (-2.3,-0.82) {$3x+2$};
  \node at (2.3,-0.82) {$11$};
  \node[teal] at (0,0.5) {même opération des deux côtés};
\end{tikzpicture}
\end{center}

\begin{center}
\begin{tikzpicture}[node distance=1.25cm and 1.4cm,every node/.style={font=\small}]
  \node[draw=teal,rounded corners,fill=mistteal,inner sep=6pt] (a) {$x=3$};
  \node[draw=teal,rounded corners,fill=mistteal,inner sep=6pt,right=of a] (b) {$2x=6$};
  \draw[teal,thick,<->] (a)--node[above]{équivalence}(b);
  \node[draw=terracotta,rounded corners,fill=mistred,inner sep=6pt,below=of a] (c) {$x=3$};
  \node[draw=terracotta,rounded corners,fill=mistred,inner sep=6pt,right=of c] (d) {$x^2=9$};
  \draw[terracotta,thick,->] (c)--node[above]{implication}(d);
  \draw[terracotta,thick,dashed,->] (d) to[bend left=20] node[below]{faux sans contrôle} (c);
\end{tikzpicture}
\end{center}

\begin{connectionbox}
Le \textbf{domaine initial} appartient au problème. Simplifier
\[
\frac{(x-2)(x+1)}{x-2}=x+1
\]
ne rend pas $x=2$ admissible : l'expression de départ n'y était pas définie. De même, après avoir élevé une équation au carré, toute solution obtenue doit être substituée dans l'équation initiale.
\end{connectionbox}

\subsection{Une stratégie plutôt qu'une collection de recettes}
La factorisation, la mise au même dénominateur, l'isolement d'une racine et le tableau de signes ne sont pas des méthodes concurrentes. Chacune révèle une structure différente : un produit, un quotient, une distance ou un changement de signe. Le bon premier geste est celui qui rend cette structure visible.


\section{Équivalence, implication et ensemble des solutions}

Une équation ne demande pas seulement de « trouver $x$ ». Elle définit un ensemble de valeurs qui rendent une affirmation vraie. Chaque transformation doit donc être jugée selon son effet sur cet ensemble.

\subsection{Équivalence et perte d'information}
Considérons
\[
x-2=0.
\]
Multiplier par $x+1$ donne
\[
(x-2)(x+1)=0.
\]
Toute solution de la première équation satisfait la seconde, mais la seconde possède aussi $x=-1$. La transformation a ajouté une possibilité parce que le facteur $x+1$ peut s'annuler. À l'inverse, diviser par une expression sans vérifier qu'elle est non nulle peut supprimer une solution.

\begin{center}
\begin{tikzpicture}[node distance=1.1cm,every node/.style={font=\small,rounded corners,inner sep=7pt,align=center}]
 \node[draw=navy,fill=mistblue] (a) {ensemble initial\\$\{2\}$};
 \node[draw=terracotta,fill=mistred,right=of a] (b) {après multiplication\\$\{-1,2\}$};
 \draw[terracotta,very thick,->] (a)--node[above]{une solution ajoutée}(b);
\end{tikzpicture}
\end{center}

\subsection{Le tableau de signes est une carte}
Pour résoudre $(x-1)(x+3)\le0$, on ne teste pas au hasard une multitude de valeurs. Les racines $-3$ et $1$ découpent la droite en trois régions où le signe de chaque facteur reste constant.

\begin{center}
\begin{tikzpicture}[x=1.25cm,y=0.8cm]
 \draw[-{Stealth},thick,slate] (-4.5,0)--(3.5,0);
 \fill[teal!20] (-3,0.15) rectangle (1,0.55);
 \draw[teal,very thick] (-3,0.35)--(1,0.35);
 \fill[navy] (-3,0) circle (2.5pt) node[below] {$-3$};
 \fill[navy] (1,0) circle (2.5pt) node[below] {$1$};
 \node at (-3.75,0.85) {$+$}; \node at (-1,0.85) {$-$}; \node at (2,0.85) {$+$};
 \node[teal] at (-1,1.35) {solution : produit négatif ou nul};
\end{tikzpicture}
\end{center}

La carte des signes relie directement la forme factorisée à une inéquation. Elle est plus informative que le seul résultat final, car elle montre où et pourquoi le signe change.

\begin{connectionbox}
Une bonne résolution conserve trois traces : le domaine initial, le type logique de chaque étape (équivalence ou implication) et une vérification finale. Ces traces rendent le raisonnement lisible et empêchent les solutions parasites de se glisser dans la réponse.
\end{connectionbox}


\section{Résoudre sans perdre d'information}

Résoudre une équation ne consiste pas seulement à déplacer des termes. Il faut savoir quelles transformations conservent exactement l'ensemble des solutions, quelles transformations peuvent en créer de nouvelles et quelles restrictions existaient avant le calcul.

\subsection{Identité, équation et équivalence}

Une identité est vraie pour toutes les valeurs où les deux membres sont définis :
\[
(x+3)^2=x^2+6x+9.
\]
Une équation demande pour quelles valeurs une égalité devient vraie :
\[
(x+3)^2=25.
\]
Dans une résolution, le symbole $\Longleftrightarrow$ annonce que les deux lignes ont exactement les mêmes solutions. Le symbole $\Longrightarrow$ annonce seulement que toute solution de la première ligne est solution de la suivante.

Les opérations suivantes produisent des équations équivalentes :
\begin{itemize}
\item ajouter ou soustraire la même expression aux deux membres;
\item multiplier ou diviser les deux membres par une quantité connue non nulle;
\item développer ou factoriser à l'aide d'une identité.
\end{itemize}

En revanche, élever au carré ou multiplier par une expression qui peut être nulle peut introduire des solutions parasites.

\subsection{Une résolution linéaire commentée}

Résolvons
\[
\frac{2x-1}{3}-\frac{x+4}{5}=2.
\]
Le plus petit dénominateur commun est $15$. Multiplier les deux membres par $15$ donne
\[
5(2x-1)-3(x+4)=30.
\]
On développe seulement après avoir supprimé les fractions :
\[
10x-5-3x-12=30,
\]
\[
7x=47,
\qquad
x=\frac{47}{7}.
\]
La substitution dans l'équation initiale constitue le contrôle le plus direct.

\subsection{Fractions algébriques : le domaine vient avant la simplification}

Étudions l'équation
\[
\frac{x^2-9}{x-3}=8.
\]
Le membre de gauche n'est pas défini pour $x=3$. Pour $x\neq3$,
\[
\frac{(x-3)(x+3)}{x-3}=x+3.
\]
On résout alors
\[
x+3=8,
\qquad x=5.
\]
La valeur $5$ est admissible. La simplification n'a pas transformé la fonction initiale en une fonction définie en $3$; elle a seulement donné une formule équivalente sur le domaine autorisé.

\begin{warningbox}
On peut annuler des \emph{facteurs} communs, jamais des termes séparés par une addition. Ainsi
\[
\frac{x(x+2)}{x}=x+2\quad(x\neq0),
\]
mais
\[
\frac{x+2}{x}\neq2.
\]
\end{warningbox}

\subsection{Équations avec une racine}

Résolvons
\[
\sqrt{2x+3}=x.
\]
Le membre de gauche est non négatif, donc toute solution doit satisfaire $x\geq0$. En élevant au carré,
\[
2x+3=x^2,
\]
\[
x^2-2x-3=0,
\]
\[
(x-3)(x+1)=0.
\]
Les candidats sont $x=3$ et $x=-1$. La condition $x\geq0$ élimine déjà $-1$, et la vérification confirme que $x=3$ est la seule solution.

L'étape d'élévation au carré est une implication, non une équivalence automatique. Le contrôle final est donc indispensable.

\subsection{Équations avec une valeur absolue}

La valeur absolue conduit naturellement à deux cas. Pour $c\geq0$,
\[
\abs{u}=c
\quad\Longleftrightarrow\quad
u=c\quad\text{ou}\quad u=-c.
\]
Par exemple,
\[
\abs{3x-2}=7
\]
donne
\[
3x-2=7\quad\text{ou}\quad3x-2=-7,
\]
d'où
\[
x=3\quad\text{ou}\quad x=-\frac53.
\]
Si le membre de droite est négatif, aucune solution n'est possible.

\subsection{Inéquations et tableau de signes}

Pour résoudre
\[
\frac{(x-2)(x+1)}{x-4}\leq0,
\]
on place les valeurs critiques $-1$, $2$ et $4$ sur une droite. Les deux premières annulent le numérateur; la dernière est interdite. Le signe de chaque facteur est constant entre deux valeurs critiques. Le tableau de signes montre alors que la solution est
\[
]-\infty,-1]\cup[2,4[.
\]
La valeur $4$ est toujours exclue, même si l'inégalité est large.

\subsection{Choisir une stratégie de factorisation}

Avant d'appliquer une formule, chercher la structure :
\begin{itemize}
\item facteur commun : $6x^3-9x^2=3x^2(2x-3)$;
\item différence de carrés : $x^2-16=(x-4)(x+4)$;
\item trinôme simple : $x^2+5x+6=(x+2)(x+3)$;
\item regroupement : $ax+ay+bx+by=(a+b)(x+y)$.
\end{itemize}
La factorisation rend visibles les zéros et les signes; le développement rend visibles les coefficients. Le bon choix dépend de la question.

\subsection{Procédure de vérification}

Après une résolution :
\begin{enumerate}
\item reprendre les restrictions du problème initial;
\item substituer chaque candidat dans l'équation initiale;
\item vérifier que les dénominateurs sont non nuls et les racines définies;
\item contrôler le signe et l'ordre de grandeur;
\item écrire l'ensemble solution clairement.
\end{enumerate}

\begin{summarybox}
Une résolution fiable distingue équivalence et implication, note le domaine avant toute simplification et vérifie les candidats dans l'énoncé initial. La structure de l'expression guide la méthode.
\end{summarybox}


\part{Fonctions et modèles}
% ============================================================

\chapter[Fonctions : définition, graphe et opérations]{Fonctions : définition, graphe\\et opérations}

\begin{questionbox}
Comment une même règle peut-elle être comprise comme une machine, un tableau, une formule et un graphe?
\end{questionbox}

\section{La notion de fonction}

\begin{intuitionbox}
Une fonction reçoit une entrée et produit une seule sortie. L'idée importante n'est pas la présence d'une formule, mais l'unicité de la sortie associée à chaque entrée autorisée.
\end{intuitionbox}

\begin{definition}{Fonction, domaine et image}{}
Une fonction $f:A\to B$ associe à chaque $x\in A$ un unique élément $f(x)\in B$. L'ensemble $A$ est le domaine. L'image est
\[
\im(f)=\{f(x):x\in A\}\subseteq B.
\]
\end{definition}

\begin{center}
\begin{tikzpicture}[node distance=1.5cm]
\node[draw,rounded corners,fill=softblue,minimum width=1.6cm,minimum height=0.9cm] (x) {$x$};
\node[draw,rounded corners,fill=softgreen,minimum width=3.2cm,minimum height=1.2cm,right=of x] (f) {règle $f$};
\node[draw,rounded corners,fill=softyellow,minimum width=1.8cm,minimum height=0.9cm,right=of f] (y) {$f(x)$};
\draw[-{Stealth},thick,uqamblue] (x)--(f);
\draw[-{Stealth},thick,uqamblue] (f)--(y);
\end{tikzpicture}
\end{center}

\begin{connectionbox}[title={Une fonction ne se réduit pas à sa formule}]
Les deux expressions
\[
f(x)=x^2\quad (x\in\R),
\qquad
g(x)=x^2\quad (x\ge0)
\]
utilisent la même règle de calcul, mais elles ne définissent pas la même fonction : leurs domaines sont différents. La seconde est injective et possède une fonction réciproque; la première ne l'est pas sur tout $\R$.
\end{connectionbox}

\begin{center}
\begin{tikzpicture}[scale=0.72]
  \draw[-{Stealth},thick,slate] (-3.1,0)--(3.3,0);
  \draw[-{Stealth},thick,slate] (0,-0.6)--(0,3.5);
  \draw[navy,very thick,domain=-2.2:2.2,samples=90] plot(\x,{0.65*\x*\x});
  \node[navy] at (-1.65,2.9) {$f:\R\to[0,+\infty[$};
\end{tikzpicture}
\qquad
\begin{tikzpicture}[scale=0.72]
  \draw[-{Stealth},thick,slate] (-0.5,0)--(3.5,0);
  \draw[-{Stealth},thick,slate] (0,-0.6)--(0,3.5);
  \draw[teal,very thick,domain=0:2.2,samples=70] plot(\x,{0.65*\x*\x});
  \node[teal] at (1.75,2.9) {$g:[0,+\infty[\to[0,+\infty[$};
\end{tikzpicture}
\end{center}

\begin{examplebox}
Pour $f(x)=x^2-1$, on a
\[
f(3)=8,
\qquad f(-3)=8.
\]
Deux entrées différentes peuvent produire la même sortie. Cela ne viole pas la définition d'une fonction.
\end{examplebox}

Le domaine dépend de la formule :
\begin{itemize}
\item un polynôme est défini sur $\R$;
\item un dénominateur ne peut pas être nul;
\item un radicande de racine carrée doit être non négatif;
\item l'argument d'un logarithme doit être strictement positif.
\end{itemize}

\begin{examplebox}
Pour
\[
g(x)=\frac{\sqrt{x-2}}{x-5},
\]
il faut $x\ge2$ et $x\neq5$. Ainsi,
\[
\dom(g)=[2,5[\ \cup\ ]5,+\infty[.
\]
\end{examplebox}

\section{Graphe et lecture graphique}

Le graphe de $f$ est l'ensemble des points $(x,f(x))$. Une courbe représente une fonction si toute droite verticale la coupe au plus une fois.

\begin{center}
\begin{tikzpicture}[scale=0.8]
\draw[-{Stealth},thick] (-4,0)--(4.3,0) node[right] {$x$};
\draw[-{Stealth},thick] (0,-2.5)--(0,3.2) node[above] {$y$};
\draw[teal,very thick,domain=-3.2:3.2,samples=100] plot(\x,{0.25*(\x)^2-1.4});
\fill[deepblue] (0,-1.4) circle (2pt) node[below right] {minimum};
\draw[dashed,uqamblue] (-3.6,1)--(3.6,1);
\node[anchor=west,uqamblue] at (3.1,1.15) {$y=1$};
\end{tikzpicture}
\end{center}

Sur un graphe, on peut lire :
\begin{itemize}
\item le domaine par projection sur l'axe des $x$;
\item l'image par projection sur l'axe des $y$;
\item les zéros aux intersections avec l'axe horizontal;
\item les intervalles de croissance et de décroissance;
\item les maximums, minimums, symétries et périodicités.
\end{itemize}

\section{Transformations de graphes}

À partir du graphe de $y=f(x)$ :
\begin{align*}
y=f(x)+k&\quad\text{translation verticale de }k,\\
y=f(x-h)&\quad\text{translation horizontale de }h,\\
y=af(x)&\quad\text{étirement vertical de facteur }\abs a,\\
y=-f(x)&\quad\text{réflexion par rapport à l'axe des }x,\\
y=f(-x)&\quad\text{réflexion par rapport à l'axe des }y.
\end{align*}

\begin{warningbox}
Le signe d'une translation horizontale semble inversé : $f(x-3)$ déplace le graphe de trois unités vers la droite, car il faut augmenter $x$ de $3$ pour retrouver la même entrée de $f$.
\end{warningbox}

\begin{examplebox}
Le graphe de
\[
g(x)=-2(x-3)^2+1
\]
s'obtient à partir de $y=x^2$ en effectuant une translation de $3$ vers la droite, un étirement vertical de facteur $2$, une réflexion par rapport à l'axe des $x$, puis une translation de $1$ vers le haut.
\end{examplebox}

\section{Opérations et composition}

Pour deux fonctions $f$ et $g$ :
\begin{align*}
(f+g)(x)&=f(x)+g(x),\\
(f-g)(x)&=f(x)-g(x),\\
(fg)(x)&=f(x)g(x),\\
\left(\frac fg\right)(x)&=\frac{f(x)}{g(x)},\qquad g(x)\neq0.
\end{align*}
Le domaine est l'intersection des domaines, avec la restriction supplémentaire $g(x)\neq0$ pour le quotient.

\begin{definition}{Composition}{}
La composition de $f$ suivie de $g$ est
\[
(g\circ f)(x)=g(f(x)).
\]
On applique d'abord la fonction la plus proche de $x$.
\end{definition}

\begin{examplebox}
Si $f(x)=x^2$ et $g(x)=2x+1$, alors
\[
(g\circ f)(x)=2x^2+1,
\qquad
(f\circ g)(x)=(2x+1)^2.
\]
L'ordre change le résultat.
\end{examplebox}

\begin{counterbox}
La composition n'est généralement pas commutative : $f\circ g\neq g\circ f$. Elle ressemble davantage à une suite d'instructions qu'à une multiplication de nombres.
\end{counterbox}

\section{Fonction réciproque}

\begin{intuitionbox}
Une fonction réciproque défait l'action de la fonction initiale. Si $f$ transforme $x$ en $y$, alors $f^{-1}$ transforme $y$ en $x$.
\end{intuitionbox}

\begin{definition}{Fonction réciproque}{}
Une fonction $f:A\to B$ est inversible s'il existe $f^{-1}:B\to A$ telle que
\[
f^{-1}(f(x))=x
\quad\text{et}\quad
f(f^{-1}(y))=y.
\]
Cela exige que chaque sortie corresponde à une unique entrée.
\end{definition}

Les graphes de $f$ et de $f^{-1}$ sont symétriques par rapport à la droite $y=x$.

\begin{examplebox}
Pour $f(x)=3x-5$,
\[
y=3x-5\quad\Longrightarrow\quad x=\frac{y+5}{3}.
\]
Ainsi,
\[
f^{-1}(x)=\frac{x+5}{3}.
\]
\end{examplebox}

La fonction $x\mapsto x^2$ n'est pas inversible sur tout $\R$, car $2$ et $-2$ ont la même image. Restreinte à $[0,+\infty[$, elle admet la réciproque $x\mapsto\sqrt x$.

\subsection{Le test de la droite horizontale}

Une fonction possède une réciproque sur son image exactement lorsque chaque valeur de sortie est atteinte au plus une fois. Sur un graphe, cela se vérifie en traçant des droites horizontales.

\begin{center}
\begin{tikzpicture}[scale=0.72]
  \draw[-{Stealth},thick,slate] (-2.8,0)--(3,0);
  \draw[-{Stealth},thick,slate] (0,-2.3)--(0,3);
  \draw[navy,very thick,domain=-2.2:2.2,samples=70] plot(\x,{0.85*\x+0.2});
  \draw[teal,thick,dashed] (-2.5,1.2)--(2.6,1.2);
  \fill[terracotta] (1.176,1.2) circle (2.3pt);
  \node[navy] at (0,-2.75) {une intersection : réciproque possible};
\end{tikzpicture}
\qquad
\begin{tikzpicture}[scale=0.72]
  \draw[-{Stealth},thick,slate] (-2.8,0)--(3,0);
  \draw[-{Stealth},thick,slate] (0,-0.8)--(0,3.5);
  \draw[teal,very thick,domain=-2.1:2.1,samples=80] plot(\x,{0.65*\x*\x});
  \draw[terracotta,thick,dashed] (-2.5,1.6)--(2.6,1.6);
  \fill[terracotta] (-1.569,1.6) circle (2.3pt);
  \fill[terracotta] (1.569,1.6) circle (2.3pt);
  \node[teal] at (0,-1.25) {deux intersections : restriction nécessaire};
\end{tikzpicture}
\end{center}

\begin{connectionbox}
Si $f:A\to B$ est bijective, alors
\[
\dom(f^{-1})=\im(f),
\qquad
\im(f^{-1})=\dom(f).
\]
Sur les graphes, le point $(a,b)$ de $f$ devient le point $(b,a)$ de $f^{-1}$ : les coordonnées échangent leurs rôles.
\end{connectionbox}

\subsection{Une réciproque rationnelle avec contrôle des domaines}

Considérons
\[
f(x)=\frac{2x+1}{x-3},
\qquad x\ne3.
\]
Posons $y=f(x)$ et résolvons pour $x$ :
\[
y(x-3)=2x+1,
\]
\[
x(y-2)=3y+1,
\]
\[
x=\frac{3y+1}{y-2}.
\]
Ainsi,
\[
f^{-1}(x)=\frac{3x+1}{x-2},
\qquad x\ne2.
\]
La valeur $2$ est exclue du domaine de la réciproque parce qu'elle n'appartient pas à l'image de $f$. Une substitution dans $f(f^{-1}(x))$ permet de vérifier le résultat sur le domaine autorisé.


\section{Situation, formule, tableau et graphe}

\begin{visualbox}
Une fonction n'est pas seulement une formule. C'est une correspondance entre des entrées et des sorties. La même fonction peut être décrite par une règle, un tableau ou un graphe; chacune de ces représentations répond à des questions différentes.
\end{visualbox}

\begin{center}
\begin{tikzpicture}[node distance=1.6cm and 1.5cm,every node/.style={font=\small}]
  \node[draw=navy,rounded corners,fill=mistblue,minimum width=3.2cm,minimum height=1.2cm,align=center] (formula) {$f(x)=2x-1$\\règle générale};
  \node[draw=teal,rounded corners,fill=mistteal,minimum width=3.2cm,minimum height=1.2cm,right=of formula] (table) {$\begin{array}{c|ccc}x&0&1&2\\\hline f(x)&-1&1&3\end{array}$};
  \node[draw=gold,rounded corners,fill=mistgold,minimum width=3.2cm,minimum height=1.2cm,below=of formula,xshift=2.35cm,align=center] (graph) {graphe :\\tendance et intersections};
  \draw[thick,teal,<->] (formula)--(table);
  \draw[thick,gold!80!black,<->] (formula)--(graph);
  \draw[thick,navy,<->] (table)--(graph);
\end{tikzpicture}
\end{center}

\begin{center}
\begin{tikzpicture}[scale=0.75]
  \draw[-{Stealth},thick,slate] (-1,0)--(5.2,0) node[right] {$x$};
  \draw[-{Stealth},thick,slate] (0,-2)--(0,4.5) node[above] {$y$};
  \draw[navy,very thick,domain=-0.4:2.6,samples=80] plot(\x,{(\x-1)^2});
  \draw[terracotta,thick,dashed] (1.8,-1.8)--(1.8,4.2);
  \fill[terracotta] (1.8,0.64) circle (2.3pt);
  \node[anchor=west] at (2.1,3.5) {une entrée};
  \node[anchor=west] at (2.1,3.0) {une seule sortie};
\end{tikzpicture}
\qquad
\begin{tikzpicture}[scale=0.75]
  \draw[-{Stealth},thick,slate] (-2.6,0)--(2.8,0) node[right] {$x$};
  \draw[-{Stealth},thick,slate] (0,-2.6)--(0,2.8) node[above] {$y$};
  \draw[teal,very thick] (0,0) circle (1.7);
  \draw[terracotta,thick,dashed] (1,-2.3)--(1,2.3);
  \fill[terracotta] (1,{sqrt(1.89)}) circle (2.3pt);
  \fill[terracotta] (1,{-sqrt(1.89)}) circle (2.3pt);
  \node at (0,-3.0) {pas une fonction de $x$};
\end{tikzpicture}
\end{center}

\begin{connectionbox}
Le domaine se lit avant le calcul : il indique quelles entrées sont permises. L'image se lit après l'application de la fonction : elle indique quelles sorties sont réellement atteintes. Pour une fonction réciproque, ces rôles s'échangent; c'est pourquoi les graphes de $f$ et $f^{-1}$ sont symétriques par rapport à $y=x$.
\end{connectionbox}

\subsection{Composition : suivre le trajet d'une valeur}
Dans $(g\circ f)(x)=g(f(x))$, la valeur entre d'abord dans $f$, puis la sortie de $f$ devient l'entrée de $g$. Dessiner deux machines en série évite presque toujours l'erreur d'ordre.


\section{Lire une fonction à partir de son graphe}

Un graphe condense une infinité de couples $(x,f(x))$. Il permet de répondre immédiatement à des questions qu'une formule peut rendre moins visibles : où la fonction est positive, quelles valeurs elle atteint, où elle augmente et si un modèle semble périodique.

\begin{center}
\begin{tikzpicture}[scale=0.86]
 \draw[-{Stealth},thick,slate] (-4.2,0)--(4.6,0) node[right] {$x$};
 \draw[-{Stealth},thick,slate] (0,-2.6)--(0,4.2) node[above] {$y$};
 \draw[navy,very thick,smooth] plot coordinates {(-3.8,-1.2) (-3,0.2) (-2.2,2.4) (-1.2,3.1) (0,2.0) (1,-0.4) (2,-1.5) (3,0.2) (4,2.6)};
 \draw[teal,thick,dashed] (-4,2)--(4.2,2);
 \fill[terracotta] (-3.15,0) circle (2.4pt);
 \fill[terracotta] (0.85,0) circle (2.4pt);
 \fill[terracotta] (2.9,0) circle (2.4pt);
 \node[terracotta,below] at (0.85,-0.25) {zéros};
 \node[teal,right] at (4.15,2) {$f(x)=2$};
 \node[gold!75!black] at (-1.2,3.55) {maximum local};
 \node[gold!75!black] at (2,-1.95) {minimum local};
\end{tikzpicture}
\end{center}

\subsection{Lire une équation sur le graphe}
Résoudre $f(x)=2$ revient à chercher les intersections avec la droite horizontale $y=2$. Résoudre $f(x)>0$ revient à repérer les portions situées au-dessus de l'axe des $x$. Résoudre $f(x)=g(x)$ revient à chercher les points d'intersection des deux graphes.

\subsection{Domaine, image et contexte}
Le domaine est la projection du graphe sur l'axe horizontal; l'image est sa projection sur l'axe vertical. Toutefois, un graphe imprimé n'est qu'une fenêtre. Il faut distinguer les bornes réellement imposées par le problème des limites du dessin.

\begin{visualbox}
Une représentation est choisie selon la question :
\begin{itemize}
\item la formule facilite les calculs exacts et les transformations;
\item le tableau révèle des régularités dans des données discrètes;
\item le graphe révèle l'allure globale et les intersections;
\item une description verbale donne le sens des variables et des unités.
\end{itemize}
La modélisation consiste souvent à passer plusieurs fois de l'une à l'autre.
\end{visualbox}

\subsection{Pourquoi une réciproque exige l'injectivité}
Si une droite horizontale coupe le graphe de $f$ en deux points, une même sortie correspond à deux entrées. Après réflexion par rapport à $y=x$, ces deux points se retrouvent sur une même droite verticale : la relation réfléchie ne peut pas être une fonction. Le test horizontal n'est donc que le test vertical appliqué au graphe réfléchi.



\subsection{Une fonction transporte aussi des unités}
Si $t$ est mesuré en heures et $d(t)$ en kilomètres, alors la pente d'un modèle affine $d(t)=mt+b$ s'exprime en kilomètres par heure. Les unités ne sont pas décoratives : elles révèlent la nature d'un paramètre. Dans une composition, les unités doivent s'enchaîner correctement. Si $C(d)$ transforme une distance en coût et $d(t)$ transforme un temps en distance, alors
\[
(C\circ d)(t)=C(d(t))
\]
transforme bien un temps en coût. L'ordre inverse $d(C(t))$ n'a généralement aucun sens.

\begin{center}
\begin{tikzpicture}[node distance=1.0cm,every node/.style={font=\small,rounded corners,inner sep=7pt,align=center}]
 \node[draw=navy,fill=mistblue] (t) {temps\\heures};
 \node[draw=teal,fill=mistteal,right=of t] (d) {distance\\kilomètres};
 \node[draw=gold,fill=mistgold,right=of d] (c) {coût\\dollars};
 \draw[->,very thick,teal] (t)--node[above] {$d$}(d);
 \draw[->,very thick,gold!80!black] (d)--node[above] {$C$}(c);
 \draw[->,very thick,terracotta,bend left=35] (t) to node[above] {$C\circ d$} (c);
\end{tikzpicture}
\end{center}

\subsection{Un graphe est une relation globale}
Une formule comme $f(x)=1/x$ donne un calcul local pour chaque $x\ne0$. Le graphe montre simultanément deux branches, une asymptote verticale, une asymptote horizontale et une symétrie. Il révèle ainsi des relations entre des valeurs éloignées que le calcul point par point ne met pas en évidence.

\begin{connectionbox}
Pour analyser une fonction, poser successivement quatre questions : quelles entrées sont permises? quelles sorties sont possibles? comment une petite variation de l'entrée affecte-t-elle la sortie? quelles structures globales - symétrie, périodicité, asymptote - organisent le graphe?
\end{connectionbox}


\section{Analyser une fonction sous plusieurs représentations}

Une fonction n'est ni seulement une formule ni seulement une courbe. C'est une règle qui associe à chaque entrée admissible une unique sortie. La formule, le tableau, le graphe et la description verbale sont quatre représentations d'un même objet; chacune rend certaines propriétés plus visibles que les autres.

\subsection{Domaine, codomaine et image}

Dans
\[
f:A\longrightarrow B,
\]
$A$ est le domaine et $B$ le codomaine. L'image est l'ensemble des valeurs réellement atteintes. Par exemple, pour
\[
f:[-2,2]\longrightarrow\R,
\qquad f(x)=x^2,
\]
on a
\[
\dom(f)=[-2,2],
\qquad
\im(f)=[0,4].
\]
Le codomaine annoncé est $\R$, mais l'image est plus petite. Cette distinction devient essentielle lorsqu'on étudie la surjectivité et les fonctions réciproques.

\subsection{Même formule, fonctions différentes}

Les fonctions
\[
f:\R\to\R,
\qquad f(x)=x^2,
\]
et
\[
g:[0,+\infty[\to\R,
\qquad g(x)=x^2
\]
ont la même formule mais pas le même domaine. La première n'est pas injective, car $f(2)=f(-2)$; la seconde est injective. Le domaine fait donc partie de l'identité de la fonction.

De même,
\[
\frac{x^2-1}{x-1}
\]
et $x+1$ donnent les mêmes valeurs pour $x\neq1$, mais la première expression n'est pas définie en $1$.

\subsection{Passer du tableau au modèle}

Considérons les données
\[
\begin{array}{c|cccc}
x&0&1&2&3\\\hline
f(x)&5&8&11&14
\end{array}
\]
La sortie augmente de $3$ lorsque l'entrée augmente de $1$. On reconnaît une variation additive constante :
\[
f(x)=3x+5.
\]
Un tableau ne prouve pas qu'une formule reste valable pour toutes les valeurs, mais il permet de proposer un modèle et d'en vérifier les paramètres.

\subsection{Lire un graphe sans connaître la formule}

Un graphe permet de lire :
\begin{itemize}
\item le domaine par projection sur l'axe horizontal;
\item l'image par projection sur l'axe vertical;
\item les zéros aux intersections avec l'axe horizontal;
\item $f(0)$ à l'intersection avec l'axe vertical;
\item les intervalles de croissance ou de décroissance;
\item les maxima, minima, symétries et répétitions.
\end{itemize}
Une valeur lue graphiquement est généralement approximative. Il faut donc écrire $x\approx2{,}4$ si le point n'est pas déterminé exactement.

\subsection{Transformations et déplacement des points}

Si le point $(a,f(a))$ appartient au graphe de $f$, alors :
\begin{itemize}
\item le graphe de $f(x)+k$ contient $(a,f(a)+k)$;
\item le graphe de $f(x-h)$ contient $(a+h,f(a))$;
\item le graphe de $-f(x)$ contient $(a,-f(a))$;
\item le graphe de $f(-x)$ contient $(-a,f(a))$.
\end{itemize}
Cette règle sur les points explique le signe apparemment inversé dans une translation horizontale.

\begin{examplebox}
À partir de $f(x)=\abs x$, construisons
\[
g(x)=-2\abs{x-3}+1.
\]
Le sommet $(0,0)$ devient $(3,1)$. Le facteur $2$ étire verticalement la courbe et le signe négatif la réfléchit par rapport à l'axe horizontal.
\end{examplebox}

\subsection{Opérations sur les fonctions et domaines}

Pour deux fonctions $f$ et $g$,
\[
(f+g)(x)=f(x)+g(x),
\qquad
(fg)(x)=f(x)g(x).
\]
Le domaine est au moins l'intersection des domaines de $f$ et de $g$. Pour le quotient,
\[
\left(\frac fg\right)(x)=\frac{f(x)}{g(x)},
\]
il faut en plus exclure les valeurs où $g(x)=0$.

Par exemple, si
\[
f(x)=\sqrt{x+1},
\qquad
g(x)=x-2,
\]
alors
\[
\dom\left(\frac fg\right)=[-1,+\infty[\setminus\{2\}.
\]

\subsection{Composition : suivre la chaîne des entrées}

Dans
\[
(g\circ f)(x)=g(f(x)),
\]
la fonction $f$ agit d'abord. Le domaine contient les $x$ tels que $x\in\dom(f)$ et $f(x)\in\dom(g)$.

Prenons
\[
f(x)=x^2-4,
\qquad g(u)=\sqrt u.
\]
Alors
\[
(g\circ f)(x)=\sqrt{x^2-4}.
\]
Il faut imposer $x^2-4\geq0$, donc
\[
\dom(g\circ f)=]-\infty,-2]\cup[2,+\infty[.
\]
L'ordre inverse donne
\[
(f\circ g)(x)=x-4,
\]
mais seulement pour $x\geq0$, car $g(x)=\sqrt x$ doit d'abord être défini.

\subsection{Fonction réciproque et test horizontal}

Une fonction possède une réciproque sur son image si chaque sortie provient d'une seule entrée. Graphiquement, aucune droite horizontale ne doit couper la courbe en plus d'un point.

Pour
\[
f(x)=2x-5,
\]
on pose $y=2x-5$, puis on résout pour $x$ :
\[
x=\frac{y+5}{2}.
\]
Ainsi
\[
f^{-1}(x)=\frac{x+5}{2}.
\]
La vérification par composition donne
\[
f^{-1}(f(x))=x,
\qquad
f(f^{-1}(x))=x.
\]

\begin{warningbox}
La notation $f^{-1}$ ne signifie pas $1/f$. La fonction réciproque inverse l'action de $f$; le quotient $1/f(x)$ prend le réciproque numérique de chaque sortie.
\end{warningbox}

\subsection{Étude intégrée d'une fonction définie par morceaux}

Considérons
\[
h(x)=
\begin{cases}
-x-1,&x<0,\\
x^2-1,&0\leq x\leq2,\\
3,&x>2.
\end{cases}
\]
Le domaine est $\R$. Il faut traiter séparément chaque règle et porter attention aux points de raccord. En $x=0$, la valeur est donnée par la deuxième règle : $h(0)=-1$. En $x=2$, $h(2)=3$, et la troisième règle commence seulement pour $x>2$. Les disques pleins et vides du graphique traduisent précisément ces inclusions.

\begin{summarybox}
Une fonction est déterminée par sa règle et son domaine. Les graphes, tableaux et formules se complètent. Pour les opérations, les compositions et les réciproques, le contrôle du domaine est aussi important que le calcul algébrique.
\end{summarybox}


\chapter{Modèles linéaires et quadratiques}

\begin{questionbox}
Comment reconnaître si une situation varie à rythme constant ou si elle possède plutôt un sommet, un maximum ou un minimum?
\end{questionbox}

\section{Le modèle affine}

\begin{definition}{Fonction affine}{}
Une fonction affine est de la forme
\[
f(x)=mx+b.
\]
Le nombre $m$ est la pente et $b=f(0)$ est l'ordonnée à l'origine.
\end{definition}

La pente entre deux points $(x_1,y_1)$ et $(x_2,y_2)$, avec $x_1\neq x_2$, est
\[
m=\frac{y_2-y_1}{x_2-x_1}.
\]
Elle mesure la variation de $y$ pour une unité de variation de $x$.

\begin{center}
\begin{tikzpicture}[scale=0.8]
\draw[-{Stealth},thick] (-1,0)--(6,0) node[right] {$x$};
\draw[-{Stealth},thick] (0,-1)--(0,5) node[above] {$y$};
\draw[uqamblue,very thick] (-0.3,0.55)--(5.5,4.6);
\fill (1,1.45) circle (2pt) node[below right] {$A$};
\fill (4,3.55) circle (2pt) node[above left] {$B$};
\draw[dashed] (1,1.45)--(4,1.45)--(4,3.55);
\node at (2.5,1.15) {$\Delta x$};
\node at (4.45,2.5) {$\Delta y$};
\end{tikzpicture}
\end{center}

\begin{examplebox}
Un taxi demande $4{,}50\,$\$ au départ et $1{,}80\,$\$ par kilomètre. Le coût en fonction de la distance $d$ est
\[
C(d)=1{,}80d+4{,}50.
\]
La pente $1{,}80$ est le coût marginal par kilomètre; l'ordonnée à l'origine est le coût fixe.
\end{examplebox}

\subsection{Déterminer une droite}

Une droite non verticale est déterminée par deux points, ou par un point et une pente. La forme point-pente est
\[
y-y_0=m(x-x_0).
\]

\begin{examplebox}
La droite passant par $(2,5)$ et de pente $-3$ vérifie
\[
y-5=-3(x-2),
\]
donc $y=-3x+11$.
\end{examplebox}

\section{Le modèle quadratique}

\begin{definition}{Fonction quadratique}{}
Une fonction quadratique est de la forme
\[
f(x)=ax^2+bx+c,
\qquad a\neq0.
\]
Son graphe est une parabole.
\end{definition}

La forme canonique
\[
f(x)=a(x-h)^2+k
\]
révèle directement le sommet $(h,k)$ et l'axe de symétrie $x=h$.

\begin{center}
\begin{tikzpicture}[scale=0.75]
\draw[-{Stealth},thick] (-4,0)--(5,0) node[right] {$x$};
\draw[-{Stealth},thick] (0,-3)--(0,5) node[above] {$y$};
\draw[teal,very thick,domain=-2.5:4.5,samples=100] plot(\x,{0.45*(\x-1)^2-2});
\draw[dashed,uqamblue] (1,-2.6)--(1,4.5) node[above] {$x=h$};
\fill[deepblue] (1,-2) circle (2.4pt) node[below right] {$(h,k)$};
\end{tikzpicture}
\end{center}

\begin{methodbox}
Pour passer de $ax^2+bx+c$ à la forme canonique :
\begin{enumerate}
\item factoriser $a$ dans les termes contenant $x$;
\item compléter le carré;
\item regrouper les constantes.
\end{enumerate}
Le sommet est aussi obtenu par
\[
h=-\frac{b}{2a},\qquad k=f(h).
\]
\end{methodbox}

\begin{examplebox}
\begin{align*}
f(x)&=2x^2-8x+5\\
&=2(x^2-4x)+5\\
&=2\bigl((x-2)^2-4\bigr)+5\\
&=2(x-2)^2-3.
\end{align*}
Le sommet est $(2,-3)$ et la parabole s'ouvre vers le haut.
\end{examplebox}

\section{Zéros et formule quadratique}

Les zéros d'une fonction quadratique sont les solutions de
\[
ax^2+bx+c=0.
\]
Le discriminant
\[
\Delta=b^2-4ac
\]
détermine le nombre de solutions réelles : deux si $\Delta>0$, une solution double si $\Delta=0$, aucune si $\Delta<0$.

\begin{theorem}{Formule quadratique}{}
Pour $a\neq0$ et $\Delta=b^2-4ac\geq0$, les solutions réelles de $ax^2+bx+c=0$ sont
\[
x=\frac{-b\pm\sqrt{b^2-4ac}}{2a}.
\]
Lorsque $\Delta<0$, il n'y a pas de solution réelle.
\end{theorem}

\begin{examplebox}
Pour $2x^2-3x-2=0$,
\[
\Delta=(-3)^2-4(2)(-2)=25,
\]
donc
\[
x=\frac{3\pm5}{4}.
\]
Les solutions sont $2$ et $-\frac12$, et
\[
2x^2-3x-2=(2x+1)(x-2).
\]
\end{examplebox}

\section{Modéliser et interpréter}

\begin{examplebox}
La hauteur d'un objet lancé est modélisée par
\[
h(t)=-4{,}9t^2+19{,}6t+1{,}5.
\]
Le coefficient négatif indique que la parabole s'ouvre vers le bas. Le maximum est atteint à
\[
t=-\frac{19{,}6}{2(-4{,}9)}=2\text{ secondes}.
\]
La hauteur maximale vaut
\[
h(2)=21{,}1\text{ mètres}.
\]
Les racines indiquent les instants où le modèle atteint la hauteur zéro; seule la racine positive possède un sens physique après le lancement.
\end{examplebox}


\section{Voir la variation : différences premières et secondes}

\begin{visualbox}
Un modèle affine ajoute toujours la même quantité lorsque $x$ augmente d'une unité. Un modèle quadratique n'a pas une variation constante, mais ses \emph{variations} changent à rythme constant. Cette différence apparaît avant même de connaître la formule.
\end{visualbox}

\begin{center}
\begin{tabular}{c@{\hspace{1.4cm}}c}
\begin{tabular}{c|ccccc}
$x$&0&1&2&3&4\\\hline
$2x+1$&1&3&5&7&9\\
$\Delta$&&2&2&2&2
\end{tabular}
&
\begin{tabular}{c|ccccc}
$x$&0&1&2&3&4\\\hline
$x^2$&0&1&4&9&16\\
$\Delta$&&1&3&5&7\\
$\Delta^2$&&&2&2&2
\end{tabular}
\end{tabular}
\end{center}

\begin{center}
\begin{tikzpicture}[scale=0.72]
  \draw[-{Stealth},thick,slate] (-1,0)--(6,0) node[right] {$x$};
  \draw[-{Stealth},thick,slate] (0,-1)--(0,5.2) node[above] {$y$};
  \draw[navy,very thick] (-0.2,0.4)--(5.4,4.6);
  \node[navy] at (4.2,3.2) {variation constante};
\end{tikzpicture}
\qquad
\begin{tikzpicture}[scale=0.72]
  \draw[-{Stealth},thick,slate] (-3,0)--(3.5,0) node[right] {$x$};
  \draw[-{Stealth},thick,slate] (0,-1)--(0,5.2) node[above] {$y$};
  \draw[teal,very thick,domain=-2.6:2.6,samples=80] plot(\x,{0.55*\x*\x+0.3});
  \node[teal] at (1.7,3.7) {sommet et courbure};
\end{tikzpicture}
\end{center}

\begin{connectionbox}
Les trois formes d'un trinôme mettent en évidence trois questions différentes :
\[
ax^2+bx+c \quad\text{(valeur initiale et coefficients)},
\]
\[
a(x-h)^2+k \quad\text{(sommet)},
\qquad
 a(x-r_1)(x-r_2) \quad\text{(zéros)}.
\]
Changer de forme ne change pas la fonction; cela change ce que l'on voit immédiatement.
\end{connectionbox}

\subsection{Le domaine vient du contexte}
Une formule peut être définie pour tout réel alors que le modèle ne l'est pas. Si $t$ représente le temps après un lancement, on impose $t\ge0$; si $h(t)$ représente une hauteur, on ne conserve que la portion où $h(t)\ge0$. Un modèle mathématique est toujours une formule accompagnée de son domaine et de ses unités.


\section{Passer de données à un modèle}

La première tâche n'est pas d'ajuster une formule compliquée, mais de reconnaître le type de variation présent dans les données.

\subsection{Différences et rapports successifs}
Considérons deux tableaux.
\[
\begin{array}{c|ccccc}
t&0&1&2&3&4\\\hline
A(t)&12&17&22&27&32
\end{array}
\qquad
\begin{array}{c|ccccc}
t&0&1&2&3&4\\\hline
B(t)&12&18&27&40{,}5&60{,}75
\end{array}
\]
Pour $A$, les différences successives sont toutes égales à $5$ : le modèle affine $A(t)=12+5t$ est naturel. Pour $B$, les rapports successifs valent $1{,}5$ : le modèle exponentiel $B(t)=12(1{,}5)^t$ est naturel. Un même jeu de points peut parfois être approximé par plusieurs modèles, mais ces diagnostics indiquent la structure la plus simple.

\subsection{Résidus : regarder ce que le modèle n'explique pas}
Si une observation vaut $y_i$ et que le modèle prédit $\widehat y_i$, le résidu est
\[
r_i=y_i-\widehat y_i.
\]
Un bon modèle ne signifie pas forcément des résidus nuls. Il signifie que les résidus sont petits par rapport à l'échelle du problème et qu'ils ne présentent pas une structure évidente.

\begin{center}
\begin{tikzpicture}[scale=0.82]
 \draw[-{Stealth},thick,slate] (-0.4,0)--(6.5,0) node[right] {$x$};
 \draw[-{Stealth},thick,slate] (0,-0.5)--(0,5.5) node[above] {$y$};
 \draw[navy,very thick] (0.2,0.9)--(6,4.7);
 \foreach \x/\y in {0.8/1.6,1.8/1.75,2.8/3.15,3.8/3.2,4.8/4.55,5.6/4.35}{
   \fill[terracotta] (\x,\y) circle (2.5pt);
   \draw[terracotta,thin,dashed] (\x,\y)--(\x,{0.655*\x+0.77});
 }
 \node[terracotta,align=left] at (5.3,1.3) {écarts verticaux\\= résidus};
\end{tikzpicture}
\end{center}

\subsection{Le sommet comme compromis}
Dans un modèle quadratique $f(x)=a(x-h)^2+k$, le sommet $(h,k)$ donne directement l'optimum. Cette forme est particulièrement naturelle lorsqu'une quantité diminue puis augmente, ou augmente puis diminue : coût total en fonction d'une production, hauteur d'un projectile, aire soumise à une contrainte.

\begin{connectionbox}
Un modèle n'est pas jugé seulement par la proximité numérique. Il doit aussi respecter le contexte : unités correctes, domaine réaliste, comportement plausible et paramètres interprétables. Une parabole peut ajuster quatre points de population sans pour autant décrire une croissance crédible à long terme.
\end{connectionbox}



\subsection{Compléter le carré comme déplacement d'une parabole}
La transformation
\[
x^2+6x+5=(x+3)^2-4
\]
peut être lue comme une opération géométrique. Le graphe de $y=x^2$ est déplacé de $3$ unités vers la gauche et de $4$ unités vers le bas. Le sommet, difficile à lire dans la forme développée, devient immédiatement visible.

\begin{center}
\begin{tikzpicture}[scale=0.72]
 \draw[-{Stealth},thick,slate] (-5,0)--(4.5,0) node[right] {$x$};
 \draw[-{Stealth},thick,slate] (0,-5)--(0,5.5) node[above] {$y$};
 \draw[navy!45,thick,dashed,domain=-2.1:2.1,samples=80] plot(\x,{0.65*\x*\x});
 \draw[teal,very thick,domain=-5:-1,samples=80] plot(\x,{0.65*(\x+3)*(\x+3)-4});
 \fill[terracotta] (-3,-4) circle (3pt);
 \draw[gold,thick,->] (0,0.4)--(-3,0.4);
 \draw[gold,thick,->] (-3,0.2)--(-3,-4);
 \node[gold!75!black,above] at (-1.5,0.4) {$-3$ horizontalement};
 \node[gold!75!black,rotate=90] at (-3.25,-1.8) {$-4$ verticalement};
 \node[teal] at (-1.5,3.4) {$y=(x+3)^2-4$};
\end{tikzpicture}
\end{center}

La même manipulation, appliquée à $ax^2+bx+c=0$, conduit à la formule quadratique. La formule n'est donc pas un résultat extérieur : elle condense une complétion du carré effectuée avec des coefficients généraux.

\subsection{Le discriminant comme distance au niveau zéro}
Dans la forme canonique
\[
f(x)=a(x-h)^2+k,
\]
les racines existent lorsque la parabole atteint le niveau $y=0$. Le discriminant encode cette situation : $\Delta>0$ correspond à deux intersections, $\Delta=0$ à un contact au sommet, et $\Delta<0$ à aucune intersection réelle. Cette lecture graphique précède utilement le calcul des racines.


\section{Construire, comparer et valider un modèle}

Un modèle mathématique ne reproduit pas toute la réalité. Il retient une structure utile : variation constante pour le modèle affine, courbure régulière et extremum pour le modèle quadratique. La qualité d'un modèle dépend autant de l'interprétation de ses paramètres que de l'exactitude des calculs.

\subsection{Déterminer un modèle affine à partir de deux points}

Une droite passant par $(x_1,y_1)$ et $(x_2,y_2)$, avec $x_1\neq x_2$, a pour pente
\[
a=\frac{y_2-y_1}{x_2-x_1}.
\]
Une fois $a$ connu, l'ordonnée à l'origine est obtenue par $b=y_1-ax_1$.

\begin{examplebox}
Un réservoir contient $620$ L au temps $t=0$ et $470$ L après $3$ h. Si la diminution est approximativement constante,
\[
a=\frac{470-620}{3-0}=-50\ \mathrm{L/h}.
\]
Le modèle est
\[
V(t)=620-50t.
\]
La pente est un taux de variation avec unité; l'ordonnée à l'origine est la quantité initiale.
\end{examplebox}

Le modèle prédit $V(t)=0$ pour $t=12{,}4$ h. Au-delà, il donnerait un volume négatif, ce qui montre que son domaine contextuel doit être limité.

\subsection{Comparer deux modèles affines}

Deux forfaits sont représentés par
\[
A(n)=18+0{,}08n,
\qquad
B(n)=30.
\]
Le seuil d'égalité est déterminé par
\[
18+0{,}08n=30,
\qquad n=150.
\]
Pour $n<150$, le forfait A coûte moins cher; pour $n>150$, le forfait B coûte moins cher. Le point d'intersection du graphique représente donc un seuil de décision.

\subsection{Trois formes d'une fonction quadratique}

Une même fonction quadratique peut être écrite sous trois formes :
\[
f(x)=ax^2+bx+c,
\]
\[
f(x)=a(x-h)^2+k,
\]
\[
f(x)=a(x-r_1)(x-r_2)
\]
lorsque deux racines réelles existent. Chaque forme répond à une question différente :
\begin{itemize}
\item la forme développée montre les coefficients et l'ordonnée à l'origine;
\item la forme canonique montre le sommet et l'axe de symétrie;
\item la forme factorisée montre les zéros et le signe.
\end{itemize}

\begin{examplebox}
Pour
\[
f(x)=2x^2-8x+5,
\]
la complétion du carré donne
\[
f(x)=2(x-2)^2-3.
\]
Le sommet est $(2,-3)$. La forme développée donne immédiatement $f(0)=5$. Les deux représentations décrivent la même parabole, mais rendent visibles des informations différentes.
\end{examplebox}

\subsection{Dérivation de la formule quadratique}

Partons de
\[
ax^2+bx+c=0,
\qquad a\neq0.
\]
En divisant par $a$ et en complétant le carré,
\[
x^2+\frac ba x=-\frac ca,
\]
\[
x^2+\frac ba x+\left(\frac b{2a}\right)^2
=
-\frac ca+\left(\frac b{2a}\right)^2,
\]
\[
\left(x+\frac b{2a}\right)^2
=
\frac{b^2-4ac}{4a^2}.
\]
En multipliant par $4a^2$, on obtient la forme plus directe
\[
(2ax+b)^2=b^2-4ac=\Delta.
\]
Lorsque $\Delta\geq0$,
\[
2ax+b=\pm\sqrt\Delta,
\]
d'où
\[
x=\frac{-b\pm\sqrt{b^2-4ac}}{2a}.
\]
La formule n'est donc pas une règle isolée : elle résume la complétion du carré.

\subsection{Le discriminant comme information graphique}

Le nombre
\[
\Delta=b^2-4ac
\]
indique le nombre d'intersections de la parabole avec l'axe horizontal :
\begin{itemize}
\item $\Delta>0$ : deux racines réelles distinctes;
\item $\Delta=0$ : une racine double, la parabole est tangente à l'axe;
\item $\Delta<0$ : aucune intersection réelle.
\end{itemize}
Le discriminant relie donc l'équation, la factorisation et le graphe.

\subsection{Étude complète d'une trajectoire}

Une balle est lancée depuis $1{,}5$ m avec une vitesse verticale initiale de $14$ m/s :
\[
h(t)=-4{,}9t^2+14t+1{,}5.
\]
Le sommet est atteint lorsque
\[
t_s=-\frac b{2a}=-\frac{14}{2(-4{,}9)}\approx1{,}43\ \mathrm{s}.
\]
La hauteur maximale est
\[
h(1{,}43)\approx11{,}5\ \mathrm{m}.
\]
Le temps d'impact est la racine positive de $h(t)=0$, environ $2{,}96$ s. La racine négative est une solution de l'équation algébrique, mais elle n'appartient pas au domaine temporel du modèle.

Cette distinction entre solution mathématique et solution contextuelle est centrale en modélisation.

\subsection{Optimisation quadratique}

Si $a<0$, la fonction
\[
f(x)=a(x-h)^2+k
\]
possède un maximum $k$ atteint en $x=h$. Si $a>0$, elle possède un minimum. Dans un problème d'aire ou de profit, la forme canonique permet de lire directement l'optimum.

\begin{examplebox}
Un rectangle a un périmètre de $40$ m. Si un côté vaut $x$, l'autre vaut $20-x$ et l'aire est
\[
A(x)=x(20-x)=-x^2+20x.
\]
En complétant le carré,
\[
A(x)=-(x-10)^2+100.
\]
L'aire maximale est $100\ \mathrm{m}^2$, obtenue pour un carré de côté $10$ m.
\end{examplebox}

\subsection{Critères de choix entre modèles affine et quadratique}

Un modèle affine est naturel lorsque les différences successives sont approximativement constantes. Un modèle quadratique est naturel lorsque les premières différences changent régulièrement ou lorsqu'un maximum/minimum est visible. La décision doit aussi considérer le contexte : une croissance réelle peut sembler affine sur un court intervalle sans l'être à long terme.

\begin{center}
\begin{tabularx}{\textwidth}{@{}lYY@{}}
\toprule
Critère & Modèle affine & Modèle quadratique\\\midrule
Variation & additive constante & variation du taux elle-même régulière\\
Graphe & droite & parabole\\
Paramètres visibles & pente, valeur initiale & sommet, ouverture, zéros\\
Usage typique & coût unitaire, vitesse constante & trajectoire, aire, optimisation\\
\bottomrule
\end{tabularx}
\end{center}

\subsection{Contrôler un modèle}

Un modèle doit être vérifié à quatre niveaux :
\begin{enumerate}
\item les unités des paramètres sont cohérentes;
\item les données utilisées sont reproduites;
\item le domaine contextuel est annoncé;
\item les prédictions restent plausibles sur l'intervalle étudié.
\end{enumerate}

\begin{summarybox}
Le modèle affine décrit une variation additive constante; le modèle quadratique décrit une courbure et un extremum. Les différentes formes algébriques rendent visibles des propriétés complémentaires, et le contexte décide quelles prédictions sont admissibles.
\end{summarybox}


\chapter[Polynômes et fonctions rationnelles]{Fonctions polynomiales\\et rationnelles}

\begin{questionbox}
Comment prévoir l'allure d'un graphe à partir de ses facteurs, de ses zéros et de ses valeurs interdites?
\end{questionbox}

\section{Architecture des fonctions polynomiales et rationnelles}

\begin{visualbox}
Une expression algébrique contient déjà une grande partie de son graphe. Pour un polynôme, le \textbf{terme dominant} décrit les extrémités et les \textbf{facteurs} décrivent les contacts avec l'axe horizontal. Pour une fonction rationnelle, le \textbf{dénominateur} ajoute des valeurs interdites, des trous ou des asymptotes.
\end{visualbox}

\begin{center}
\begin{tabularx}{\textwidth}{@{}lYY@{}}
\toprule
Information & Polynôme $P(x)$ & Fonction rationnelle $P(x)/Q(x)$\\\midrule
Domaine & tout $\R$ & $Q(x)\ne0$\\
Zéros & facteurs du numérateur & zéros du numérateur admissibles\\
Comportement lointain & terme dominant & comparaison des degrés, puis division si nécessaire\\
Comportement local & multiplicité des racines & trou ou asymptote près d'une valeur interdite\\
\bottomrule
\end{tabularx}
\end{center}

\begin{center}
\begin{tikzpicture}[scale=0.68]
  \draw[-{Stealth},thick,slate] (-3.3,0)--(3.5,0);
  \draw[-{Stealth},thick,slate] (0,-2.4)--(0,3.2);
  \draw[navy,very thick,domain=-2.8:2.8,samples=100] plot(\x,{0.17*(\x+2)*(\x-0.4)*(\x-0.4)});
  \fill[terracotta] (-2,0) circle (2.2pt);
  \fill[gold] (0.4,0) circle (2.2pt);
  \node[below] at (-2,-0.15) {racine simple};
  \node[above] at (0.4,0.15) {racine double};
\end{tikzpicture}
\qquad
\begin{tikzpicture}[scale=0.68]
  \draw[-{Stealth},thick,slate] (-3.3,0)--(3.5,0);
  \draw[-{Stealth},thick,slate] (0,-2.4)--(0,3.2);
  \draw[terracotta,thick,dashed] (1,-2.2)--(1,3.0);
  \draw[teal,very thick,domain=-2.7:0.82,samples=80] plot(\x,{0.75/(\x-1)});
  \draw[teal,very thick,domain=1.18:3,samples=80] plot(\x,{0.75/(\x-1)});
  \node[terracotta,rotate=90] at (1.23,1.65) {valeur interdite};
\end{tikzpicture}
\end{center}

\begin{connectionbox}
Une simplification conserve une formule sur le domaine autorisé, mais elle ne restaure jamais une valeur exclue. Ainsi,
\[
\frac{(x-1)(x+2)}{x-1}=x+2\qquad (x\ne1),
\]
et le graphe possède un trou en $(1,3)$.
\end{connectionbox}

\section{Prévoir un graphe à partir de l'expression}

\begin{visualbox}
Pour un polynôme, le terme dominant détermine l'allure aux extrémités et les facteurs déterminent les contacts avec l'axe des $x$. Pour une fonction rationnelle, le dénominateur indique où le graphe peut se briser; il faut ensuite distinguer un facteur simplifiable, qui produit un trou, d'un facteur restant, qui produit une asymptote verticale.
\end{visualbox}

\begin{center}
\begin{tikzpicture}[scale=0.68]
  \draw[-{Stealth},thick,slate] (-3.2,0)--(3.4,0);
  \draw[-{Stealth},thick,slate] (0,-2.5)--(0,3.2);
  \draw[navy,very thick,domain=-2.7:2.7,samples=90] plot(\x,{0.18*(\x+2)*(\x-0.5)*(\x-0.5)});
  \fill[terracotta] (-2,0) circle (2pt);
  \fill[gold] (0.5,0) circle (2pt);
  \node[below] at (-2,-0.15) {traverse};
  \node[above] at (0.5,0.15) {touche};
\end{tikzpicture}
\qquad
\begin{tikzpicture}[scale=0.68]
  \draw[-{Stealth},thick,slate] (-3.2,0)--(3.4,0);
  \draw[-{Stealth},thick,slate] (0,-2.5)--(0,3.2);
  \draw[terracotta,thick,dashed] (1,-2.3)--(1,3.0);
  \draw[teal,very thick,domain=-2.7:0.82,samples=90] plot(\x,{0.75/(\x-1)});
  \draw[teal,very thick,domain=1.18:3,samples=90] plot(\x,{0.75/(\x-1)});
  \node[terracotta,rotate=90] at (1.22,1.8) {asymptote};
\end{tikzpicture}
\end{center}

\begin{center}
\begin{tikzpicture}[scale=0.75]
  \draw[-{Stealth},thick,slate] (-3.4,0)--(3.5,0);
  \draw[-{Stealth},thick,slate] (0,-2)--(0,3.2);
  \draw[navy,very thick,domain=-3:3,samples=80] plot(\x,{0.75*\x+0.5});
  \draw[fill=white,draw=terracotta,very thick] (1,1.25) circle (3pt);
  \node[terracotta,anchor=west] at (1.25,1.6) {trou : valeur exclue};
\end{tikzpicture}
\end{center}

\begin{connectionbox}
La simplification algébrique explique la forme visible, mais le domaine initial explique les points absents. Ainsi,
\[
\frac{x^2-1}{x-1}=x+1\quad\text{pour }x\ne1.
\]
Le graphe est la droite $y=x+1$ avec un trou en $(1,2)$. Dans $1/(x-1)$, le facteur $x-1$ ne disparaît pas : le graphe s'éloigne sans borne près de $x=1$.
\end{connectionbox}

\subsection{Une procédure de lecture}
Avant de calculer beaucoup de points, déterminer : le domaine, les zéros, l'ordonnée à l'origine, les multiplicités, les asymptotes et le comportement à l'infini. Quelques informations structurales valent mieux qu'une longue table de valeurs.


\section{Comportement local et comportement global}

Deux informations différentes construisent le graphe d'une fonction algébrique. Les facteurs et les valeurs interdites décrivent ce qui se passe près de points particuliers; le terme dominant ou le rapport des degrés décrit ce qui se passe loin de l'origine.

\subsection{Près d'une racine}
Pour $P(x)=(x-1)^mQ(x)$ avec $Q(1)\ne0$, la parité de $m$ détermine le contact avec l'axe. Si $m$ est impair, le facteur change de signe et le graphe traverse l'axe. Si $m$ est pair, le facteur reste positif de part et d'autre et le graphe touche l'axe sans le traverser.

\begin{center}
\begin{tikzpicture}[scale=0.75]
 \draw[-{Stealth},thick,slate] (-2.7,0)--(2.9,0);
 \draw[-{Stealth},thick,slate] (0,-2.2)--(0,2.7);
 \draw[navy,very thick,domain=-2:2,samples=80] plot(\x,{0.55*\x});
 \fill[terracotta] (0,0) circle (2.5pt);
 \node[below] at (0,-2.55) {multiplicité impaire};
\end{tikzpicture}
\qquad
\begin{tikzpicture}[scale=0.75]
 \draw[-{Stealth},thick,slate] (-2.7,0)--(2.9,0);
 \draw[-{Stealth},thick,slate] (0,-0.6)--(0,3.6);
 \draw[teal,very thick,domain=-2:2,samples=80] plot(\x,{0.65*\x*\x});
 \fill[terracotta] (0,0) circle (2.5pt);
 \node[below] at (0,-0.95) {multiplicité paire};
\end{tikzpicture}
\end{center}

\subsection{Près d'une valeur interdite}
Pour $R(x)=P(x)/Q(x)$, une racine commune de $P$ et $Q$ peut disparaître après simplification, mais le point reste exclu : on obtient un trou. Une racine du dénominateur qui ne s'annule pas dans le numérateur produit généralement une asymptote verticale.

\subsection{Loin de l'origine}
Pour un polynôme, les termes de degré inférieur deviennent négligeables devant le terme dominant. Par exemple,
\[
P(x)=-2x^5+7x^2-1
\]
se comporte comme $-2x^5$ lorsque $|x|$ est grand. Pour une fonction rationnelle, la comparaison des degrés fournit l'asymptote horizontale dans les cas usuels.

\begin{visualbox}
Une esquisse fiable se construit dans cet ordre :
\begin{enumerate}
\item domaine et valeurs interdites;
\item zéros et multiplicités;
\item signe entre les points critiques;
\item asymptotes;
\item comportement aux extrémités;
\item quelques valeurs seulement pour ajuster la forme.
\end{enumerate}
La table de valeurs intervient à la fin, non au début.
\end{visualbox}


\section{Du facteur au graphe}

Les fonctions polynomiales et rationnelles se lisent à plusieurs échelles. Le terme dominant décrit le comportement très loin de l'origine; les facteurs décrivent ce qui se passe près des zéros; le dénominateur révèle les valeurs interdites et les asymptotes. Une étude complète combine ces trois niveaux.

\subsection{Le terme dominant et les extrémités du graphe}

Pour un polynôme
\[
P(x)=a_nx^n+a_{n-1}x^{n-1}+\cdots+a_0,
\]
le terme $a_nx^n$ domine lorsque $\abs x$ devient grand. Le degré et le signe de $a_n$ déterminent donc les extrémités du graphe :
\begin{itemize}
\item degré pair, $a_n>0$ : les deux extrémités montent;
\item degré pair, $a_n<0$ : les deux extrémités descendent;
\item degré impair, $a_n>0$ : gauche vers le bas, droite vers le haut;
\item degré impair, $a_n<0$ : gauche vers le haut, droite vers le bas.
\end{itemize}

Cette information ne donne pas les détails au centre du graphe, mais elle fournit son orientation globale.

\subsection{Racines et multiplicités}

Si
\[
P(x)=a(x-r)^mQ(x)
\]
et $Q(r)\neq0$, alors $r$ est une racine de multiplicité $m$. Une multiplicité impaire change le signe du facteur $(x-r)^m$ : la courbe traverse l'axe. Une multiplicité paire conserve le signe : la courbe touche l'axe puis repart du même côté.

\begin{examplebox}
Pour
\[
P(x)=-(x+2)(x-1)^2(x-4)^3,
\]
les racines sont $-2$, $1$ et $4$. La courbe traverse l'axe en $-2$ et $4$, mais le touche seulement en $1$. Le degré total est $6$ et le coefficient dominant est négatif; les deux extrémités descendent.
\end{examplebox}

\subsection{Construire un tableau de signes à partir des facteurs}

Dans une forme factorisée, le signe est constant entre deux racines consécutives. Il suffit de choisir un nombre test dans chaque intervalle ou de suivre les changements imposés par les multiplicités.

Pour
\[
P(x)=(x+3)(x-2)^2,
\]
le facteur carré est non négatif. Le signe de $P$ est donc essentiellement celui de $x+3$, sauf aux zéros. On obtient
\[
P(x)<0\quad\text{pour }x<-3,
\]
\[
P(x)=0\quad\text{pour }x=-3\text{ ou }x=2,
\]
\[
P(x)>0\quad\text{pour }x>-3,\ x\neq2.
\]

\subsection{Fonctions rationnelles : domaine avant tout}

Une fonction rationnelle
\[
f(x)=\frac{P(x)}{Q(x)}
\]
est définie lorsque $Q(x)\neq0$. Les zéros du dénominateur doivent être notés avant toute simplification.

Considérons
\[
f(x)=\frac{x^2-4}{x^2-x-2}
=
\frac{(x-2)(x+2)}{(x-2)(x+1)}.
\]
Le domaine exclut $x=2$ et $x=-1$. Pour les autres valeurs,
\[
f(x)=\frac{x+2}{x+1}.
\]
Le facteur simplifié produit un trou en $x=2$; le facteur non simplifié du dénominateur produit une asymptote verticale en $x=-1$.

La hauteur du trou est obtenue à partir de la formule simplifiée :
\[
\frac{2+2}{2+1}=\frac43.
\]
Le point absent est donc $(2,4/3)$.

\subsection{Asymptotes verticales, horizontales et obliques}

Une asymptote verticale apparaît généralement lorsqu'un facteur du dénominateur demeure après simplification. Le comportement de part et d'autre dépend du signe des facteurs.

Pour le comportement à l'infini :
\begin{itemize}
\item si $\deg P<\deg Q$, l'asymptote horizontale est $y=0$;
\item si $\deg P=\deg Q$, elle est le rapport des coefficients dominants;
\item si $\deg P=\deg Q+1$, la division polynomiale donne une asymptote oblique.
\end{itemize}

\begin{examplebox}
Pour
\[
g(x)=\frac{2x^2+3x-1}{x^2-4},
\]
les degrés sont égaux; l'asymptote horizontale est
\[
y=\frac21=2.
\]
Cela signifie que $g(x)-2$ tend vers $0$ lorsque $\abs x$ grandit, non que la fonction ne peut jamais couper la droite $y=2$.
\end{examplebox}

\subsection{Division polynomiale et lecture d'une asymptote}

Étudions
\[
h(x)=\frac{x^2+1}{x-1}.
\]
La division donne
\[
x^2+1=(x-1)(x+1)+2,
\]
donc
\[
h(x)=x+1+\frac{2}{x-1}.
\]
Lorsque $\abs x$ est grand, le dernier terme devient petit. La droite
\[
y=x+1
\]
est une asymptote oblique. Cette décomposition montre aussi de quel côté de l'asymptote se trouve la courbe selon le signe de $2/(x-1)$.

\subsection{Résoudre une équation rationnelle}

Résolvons
\[
\frac{2}{x-1}+\frac{1}{x+2}=1.
\]
Les valeurs interdites sont $x=1$ et $x=-2$. En multipliant par $(x-1)(x+2)$,
\[
2(x+2)+(x-1)=(x-1)(x+2).
\]
Après développement,
\[
3x+3=x^2+x-2,
\]
\[
x^2-2x-5=0.
\]
Les candidats sont
\[
x=1\pm\sqrt6.
\]
Aucun n'est interdit; les deux sont solutions. La multiplication par le dénominateur commun n'est équivalente qu'après avoir exclu ses zéros.

\subsection{Étude complète d'une fonction rationnelle}

Pour analyser
\[
r(x)=\frac{(x-1)(x+3)}{(x-2)(x+1)},
\]
on suit une séquence stable :
\begin{enumerate}
\item domaine : $x\neq2,-1$;
\item zéros : $x=1,-3$;
\item ordonnée à l'origine : $r(0)=3/2$;
\item asymptotes verticales : $x=2$ et $x=-1$;
\item asymptote horizontale : $y=1$;
\item signe : tableau construit avec les quatre valeurs critiques;
\item comportement près des asymptotes : étude du signe du numérateur et du dénominateur.
\end{enumerate}
Cette procédure évite de tracer la courbe à partir d'une simple table de valeurs, qui pourrait manquer les changements rapides près des valeurs interdites.

\subsection{Erreurs fréquentes et contrôles}

\begin{itemize}
\item Une racine du dénominateur simplifiée reste exclue du domaine.
\item Une asymptote est une description de comportement, non une barrière infranchissable dans tous les cas.
\item Les degrés donnent l'asymptote à l'infini, mais pas les zéros ni les trous.
\item Une solution d'équation rationnelle doit être vérifiée dans le domaine initial.
\end{itemize}

\begin{summarybox}
Le terme dominant décrit les extrémités d'un polynôme; les facteurs décrivent ses zéros et son signe. Pour une fonction rationnelle, le domaine doit être établi avant la simplification, puis les trous, asymptotes, zéros et signes sont étudiés séparément.
\end{summarybox}


\chapter[Fonctions exponentielles et logarithmiques]{Fonctions exponentielles\\et logarithmiques}

\begin{questionbox}
Comment modéliser une variation multiplicative répétée, puis retrouver le temps ou l'exposant nécessaire pour atteindre une valeur donnée?
\end{questionbox}

\section{Variation additive et variation multiplicative}

\begin{intuitionbox}
Dans un modèle linéaire, on ajoute toujours la même quantité. Dans un modèle exponentiel, on multiplie toujours par le même facteur.
\end{intuitionbox}

\begin{definition}{Fonction exponentielle}{}
Pour $b>0$ et $b\neq1$, la fonction exponentielle de base $b$ est
\[
f(x)=b^x.
\]
Son domaine est $\R$ et son image est $]0,+\infty[$.
\end{definition}

Si $b>1$, la fonction est croissante. Si $0<b<1$, elle est décroissante. Dans tous les cas, $b^0=1$.

\begin{center}
\begin{tikzpicture}[scale=0.78]
\draw[-{Stealth},thick] (-4,0)--(4.5,0) node[right] {$x$};
\draw[-{Stealth},thick] (0,-0.5)--(0,5) node[above] {$y$};
\draw[teal,very thick,domain=-3.2:2.15,samples=100] plot(\x,{2^(\x)});
\draw[gold,very thick,dashed,domain=-2.15:3.2,samples=100] plot(\x,{0.5^(\x)});
\node[teal] at (2.6,4.3) {$2^x$};
\node[gold] at (-2.7,4.3) {$(1/2)^x$};
\fill (0,1) circle (2pt) node[above right] {$(0,1)$};
\end{tikzpicture}
\end{center}

\begin{definition}{Modèle exponentiel}{}
Un modèle exponentiel s'écrit
\[
Q(t)=Q_0a^t
\]
ou, si la variation relative continue est utilisée,
\[
Q(t)=Q_0\e^{kt}.
\]
Le nombre $Q_0$ est la valeur initiale.
\end{definition}

\begin{examplebox}
Une quantité de $1200$ augmente de $6\%$ par année. Le facteur annuel est $1{,}06$, donc
\[
Q(t)=1200(1{,}06)^t.
\]
Après cinq années,
\[
Q(5)\approx1605{,}87.
\]
\end{examplebox}

Pour une diminution de $p\%$, le facteur est $1-p$ lorsque $p$ est écrit sous forme décimale. Une baisse de $18\%$ correspond donc à un facteur $0{,}82$.

\section{Le logarithme comme fonction réciproque}

\begin{definition}{Logarithme}{}
Pour $b>0$ et $b\neq1$,
\[
\log_b(x)=y
\quad\Longleftrightarrow\quad
b^y=x.
\]
Le domaine de $\log_b$ est $]0,+\infty[$ et son image est $\R$.
\end{definition}

\begin{examplebox}
\[
\log_2(32)=5
\quad\text{car}\quad
2^5=32,
\]
et
\[
\log_{10}(0{,}001)=-3
\quad\text{car}\quad
10^{-3}=0{,}001.
\]
\end{examplebox}

Les graphes de $b^x$ et $\log_b x$ sont symétriques par rapport à $y=x$.

\section{Lois des logarithmes}

Pour $M,N>0$ et $r\in\R$ :
\begin{align*}
\log_b(MN)&=\log_bM+\log_bN,\\
\log_b\left(\frac MN\right)&=\log_bM-\log_bN,\\
\log_b(M^r)&=r\log_bM.
\end{align*}

\begin{warningbox}
Il n'existe pas de règle simple pour le logarithme d'une somme :
\[
\log_b(M+N)\neq\log_bM+\log_bN.
\]
Les lois logarithmiques transforment produits en sommes, quotients en différences et puissances en facteurs.
\end{warningbox}

\subsection{Passage d'une base à une autre}

\begin{theorem}{Formule de changement de base}{}
Pour $a>0$, $a\neq1$, $b>0$, $b\neq1$ et $x>0$,
\[
\log_b x=\frac{\log_a x}{\log_a b}.
\]
En particulier,
\[
\log_bx=\frac{\ln x}{\ln b}.
\]
\end{theorem}

\begin{examplebox}
Une calculatrice peut ne pas offrir directement $\log_2$. On calcule alors
\[
\log_2(7)=\frac{\ln7}{\ln2}\approx2{,}807.
\]
\end{examplebox}

\section{Équations exponentielles et logarithmiques}

\begin{methodbox}
Pour une équation exponentielle :
\begin{enumerate}
\item essayer d'écrire les deux membres avec la même base;
\item sinon, isoler l'exponentielle et appliquer un logarithme;
\item interpréter la solution dans le contexte.
\end{enumerate}
\end{methodbox}

\begin{examplebox}
Résolvons $3\cdot2^x=20$ :
\[
2^x=\frac{20}{3},
\qquad
x=\log_2\left(\frac{20}{3}\right)=\frac{\ln(20/3)}{\ln2}\approx2{,}737.
\]
\end{examplebox}

Pour une équation logarithmique, il faut commencer par écrire les conditions de domaine.

\begin{examplebox}
Résolvons
\[
\ln(x-1)+\ln(x+1)=\ln8.
\]
Le domaine impose $x>1$. Les lois donnent
\[
\ln(x^2-1)=\ln8,
\]
donc $x^2-1=8$, soit $x=\pm3$. Le domaine conserve seulement $x=3$.
\end{examplebox}

\section{Applications}

\begin{examplebox}[title={Exemple : demi-vie}]
Une substance possède une demi-vie de $12$ ans. Sa quantité est
\[
Q(t)=Q_0\left(\frac12\right)^{t/12}.
\]
Pour savoir quand il reste $10\%$ de la quantité initiale, on résout
\[
0{,}1=\left(\frac12\right)^{t/12}.
\]
Ainsi,
\[
t=12\frac{\ln(0{,}1)}{\ln(1/2)}\approx39{,}86\text{ ans}.
\]
\end{examplebox}

\begin{examplebox}[title={Exemple : intérêts composés}]
Un capital $C_0$ placé au taux annuel $r$, capitalisé $n$ fois par année, devient
\[
C(t)=C_0\left(1+\frac rn\right)^{nt}.
\]
Cette formule distingue clairement la valeur initiale, le facteur de croissance par période et le nombre total de périodes.
\end{examplebox}


\section{Interprétation d'un facteur multiplicatif}

\begin{visualbox}
La distinction fondamentale n'est pas « droite ou courbe », mais \textbf{ajouter la même quantité} ou \textbf{multiplier par le même facteur}. Une hausse de $5$ unités par période produit un modèle affine; une hausse de $5\%$ par période produit un modèle exponentiel.
\end{visualbox}

\begin{center}
\begin{tikzpicture}[node distance=0.9cm,every node/.style={font=\small,minimum width=1.45cm,minimum height=0.75cm,rounded corners}]
  \node[draw=navy,fill=mistblue] (a0) {$100$};
  \node[draw=navy,fill=mistblue,right=of a0] (a1) {$105$};
  \node[draw=navy,fill=mistblue,right=of a1] (a2) {$110$};
  \node[draw=navy,fill=mistblue,right=of a2] (a3) {$115$};
  \draw[navy,thick,->] (a0)--node[above] {$+5$}(a1);
  \draw[navy,thick,->] (a1)--node[above] {$+5$}(a2);
  \draw[navy,thick,->] (a2)--node[above] {$+5$}(a3);
  \node[draw=teal,fill=mistteal,below=1.3cm of a0] (m0) {$100$};
  \node[draw=teal,fill=mistteal,right=of m0] (m1) {$105$};
  \node[draw=teal,fill=mistteal,right=of m1] (m2) {$110{,}25$};
  \node[draw=teal,fill=mistteal,right=of m2] (m3) {$115{,}76$};
  \draw[teal,thick,->] (m0)--node[above] {$\times1{,}05$}(m1);
  \draw[teal,thick,->] (m1)--node[above] {$\times1{,}05$}(m2);
  \draw[teal,thick,->] (m2)--node[above] {$\times1{,}05$}(m3);
\end{tikzpicture}
\end{center}

\begin{connectionbox}
Le logarithme répond à une question inverse : combien de multiplications identiques faut-il effectuer pour atteindre une valeur donnée? Si $b^x=y$, alors $x=\log_b y$. Les lois logarithmiques traduisent les lois des exposants : un produit devient une somme parce que les exposants s'additionnent.
\end{connectionbox}

\begin{center}
\begin{tikzpicture}[scale=0.72]
  \draw[-{Stealth},thick,slate] (-2.5,0)--(5.2,0) node[right] {$x$};
  \draw[-{Stealth},thick,slate] (0,-2.2)--(0,5) node[above] {$y$};
  \draw[navy,very thick,domain=-2:2.15,samples=80] plot(\x,{pow(2,\x)});
  \draw[teal,very thick,domain=0.16:4.5,samples=80] plot(\x,{ln(\x)/ln(2)});
  \draw[gold,thick,dashed] (-1.8,-1.8)--(4.5,4.5);
  \node[navy] at (2.6,4.1) {$y=2^x$};
  \node[teal] at (3.6,1.5) {$y=\log_2x$};
  \node[gold!75!black,rotate=45] at (2.9,3.15) {$y=x$};
\end{tikzpicture}
\end{center}

\subsection{Temps caractéristiques}
Le temps de doublement ou la demi-vie est souvent plus intuitif qu'un taux. Pour $Q(t)=Q_0b^t$, le temps $T$ de doublement vérifie $b^T=2$, donc $T=\log_b2$. Cette formule n'est pas un résultat isolé : elle exprime directement le rôle inverse du logarithme.


\section{Trois façons de voir une exponentielle}

Une fonction exponentielle peut être comprise comme une suite de multiplications, une courbe dont le taux relatif est constant ou une droite après changement d'échelle logarithmique.

\subsection{Le facteur par période}
Dans $Q(t)=Q_0b^t$, le rapport
\[
\frac{Q(t+1)}{Q(t)}=b
\]
est constant. Si $b=1+r$, alors $r$ est le taux de variation par période. Une diminution de $12\%$ correspond à $b=0{,}88$, et non à $-0{,}12$.

\subsection{Le logarithme redresse la courbe}
En prenant le logarithme,
\[
\ln Q(t)=\ln Q_0+t\ln b.
\]
Le graphe de $\ln Q$ en fonction de $t$ est donc affine. Cette relation explique pourquoi les échelles semi-logarithmiques sont utiles pour reconnaître une croissance exponentielle.

\begin{center}
\begin{tikzpicture}[scale=0.78]
 \draw[-{Stealth},thick,slate] (-0.3,0)--(5.8,0) node[right] {$t$};
 \draw[-{Stealth},thick,slate] (0,-0.3)--(0,5.4) node[above] {$Q$};
 \draw[navy,very thick,domain=0:4.5,samples=80] plot(\x,{0.55*pow(1.65,\x)});
 \node[navy] at (4.1,4.6) {échelle ordinaire};
\end{tikzpicture}
\qquad
\begin{tikzpicture}[scale=0.78]
 \draw[-{Stealth},thick,slate] (-0.3,0)--(5.8,0) node[right] {$t$};
 \draw[-{Stealth},thick,slate] (0,-0.3)--(0,5.4) node[above] {$\ln Q$};
 \draw[teal,very thick] (0,0.5)--(4.8,4.7);
 \node[teal] at (3.8,3.0) {droite};
\end{tikzpicture}
\end{center}

\subsection{Valeur d'équilibre}
Certaines situations exponentielles ne tendent pas vers zéro. Une température se rapproche de la température ambiante $T_a$ :
\[
T(t)=T_a+(T_0-T_a)b^t,
\qquad 0<b<1.
\]
C'est l'écart $T(t)-T_a$, et non la température elle-même, qui est multiplié par un facteur constant.

\begin{connectionbox}
Pour choisir un modèle exponentiel, chercher une quantité dont le \textbf{rapport} entre deux périodes est stable. Pour choisir un modèle affine, chercher une quantité dont la \textbf{différence} est stable. Cette comparaison est plus robuste que l'impression visuelle « la courbe monte vite ».
\end{connectionbox}


\section{Modéliser une variation multiplicative}

Une variation affine ajoute une quantité fixe; une variation exponentielle multiplie par un facteur fixe. Cette différence explique pourquoi une croissance exponentielle peut sembler lente au début puis devenir très rapide, et pourquoi le logarithme apparaît dès qu'on cherche le temps nécessaire pour atteindre un seuil.

\subsection{Du pourcentage au modèle exponentiel}

Si une quantité augmente de $r\%$ à chaque période, elle est multipliée par
\[
b=1+\frac r{100}.
\]
Après $n$ périodes,
\[
Q_n=Q_0b^n.
\]
Pour une baisse de $r\%$, le facteur est $1-r/100$.

\begin{examplebox}
Un capital de $5000\,$\$ augmente de $4{,}2\%$ par année. Le modèle est
\[
C(t)=5000(1{,}042)^t.
\]
Après $8$ ans,
\[
C(8)\approx5000(1{,}042)^8\approx6948{,}57\ \$.
\]
Le taux $4{,}2\%$ n'est pas ajouté chaque année au capital initial; il agit sur la valeur accumulée.
\end{examplebox}

\subsection{Déterminer un modèle à partir de deux mesures}

Supposons
\[
Q(t)=Ab^t
\]
et que $Q(2)=180$, $Q(5)=486$. Alors
\[
Ab^2=180,
\qquad
Ab^5=486.
\]
En divisant,
\[
b^3=\frac{486}{180}=2{,}7,
\qquad
b=\sqrt[3]{2{,}7}.
\]
Puis
\[
A=\frac{180}{b^2}.
\]
La division des deux équations élimine la valeur initiale et isole le facteur de croissance.

\subsection{Temps de doublement et demi-vie}

Pour un modèle $Q(t)=Q_0b^t$ avec $b>1$, le temps de doublement $T_d$ satisfait
\[
b^{T_d}=2,
\]
d'où
\[
T_d=\frac{\ln2}{\ln b}.
\]
Pour une décroissance, la demi-vie satisfait
\[
b^{T_{1/2}}=\frac12,
\qquad
T_{1/2}=\frac{\ln(1/2)}{\ln b}.
\]
Ces temps caractéristiques ne dépendent pas de $Q_0$; ils décrivent le rythme du modèle.

\subsection{Le logarithme comme opération réciproque}

L'égalité
\[
\log_bx=y
\]
signifie exactement
\[
b^y=x.
\]
Le logarithme répond donc à la question : « à quel exposant faut-il élever $b$ pour obtenir $x$? » Cette définition explique les valeurs fondamentales
\[
\log_b1=0,
\qquad
\log_bb=1.
\]
Elle explique aussi le domaine $x>0$, puisque $b^y$ est toujours positif lorsque $b>0$.

\subsection{Justification des lois logarithmiques}

Posons
\[
u=\log_bx,
\qquad v=\log_by.
\]
Alors $x=b^u$ et $y=b^v$. Par les lois des exposants,
\[
xy=b^{u+v}.
\]
Donc
\[
\log_b(xy)=u+v=\log_bx+\log_by.
\]
De la même façon,
\[
\log_b\left(\frac xy\right)=\log_bx-\log_by,
\]
et
\[
\log_b(x^p)=p\log_bx.
\]
Ces lois concernent les produits, quotients et puissances; il n'existe aucune loi analogue pour une somme :
\[
\log_b(x+y)\neq\log_bx+\log_by.
\]

\subsection{Dérivation du changement de base}

Si
\[
y=\log_bx,
\]
alors $b^y=x$. En appliquant le logarithme naturel,
\[
y\ln b=\ln x.
\]
Donc
\[
\log_bx=\frac{\ln x}{\ln b}.
\]
La même démonstration fonctionne avec n'importe quelle base admissible. Cette formule permet à une calculatrice possédant seulement les touches $\ln$ et $\log$ de calculer un logarithme en base quelconque.

\subsection{Équations exponentielles}

Si les deux membres peuvent être écrits avec la même base, on compare les exposants :
\[
9^{x-1}=27^{2x}
\]
devient
\[
3^{2x-2}=3^{6x},
\]
d'où $2x-2=6x$ et $x=-1/2$.

Sinon, on prend un logarithme :
\[
5\cdot1{,}08^t=12
\]
donne
\[
t=\frac{\ln(12/5)}{\ln1{,}08}.
\]
Le logarithme est appliqué après avoir isolé l'exponentielle.

\subsection{Équations logarithmiques et contrôle du domaine}

Résolvons
\[
\ln(x-1)+\ln(x+2)=\ln12.
\]
Le domaine exige $x>1$. En regroupant,
\[
\ln((x-1)(x+2))=\ln12,
\]

aussi
\[
(x-1)(x+2)=12.
\]
On obtient
\[
x^2+x-14=0,
\]

donc
\[
x=\frac{-1\pm\sqrt{57}}2.
\]
Seule la racine supérieure à $1$ est admissible. Le domaine doit être établi avant l'utilisation des lois logarithmiques.

\subsection{Refroidissement : un modèle avec valeur d'équilibre}

Un objet à $90^\circ$C est placé dans une pièce à $20^\circ$C. Un modèle simple est
\[
T(t)=20+70\e^{-0{,}18t}.
\]
La quantité qui décroît exponentiellement n'est pas la température totale, mais l'écart à la température ambiante :
\[
T(t)-20=70\e^{-0{,}18t}.
\]
Pour trouver quand $T(t)=35$, on résout
\[
15=70\e^{-0{,}18t},
\]
\[
t=-\frac1{0{,}18}\ln\left(\frac{15}{70}\right)\approx8{,}56.
\]
Le modèle ne passe jamais exactement sous $20^\circ$C et s'en approche progressivement.

\subsection{Échelles logarithmiques}

Les logarithmes transforment une multiplication en addition. C'est pourquoi ils servent à représenter des grandeurs couvrant plusieurs ordres de grandeur. Une augmentation d'une unité sur une échelle logarithmique de base $10$ correspond à une multiplication par $10$ de la grandeur mesurée.

Cette compression est utile pour les intensités sonores, certaines mesures de concentration et les données dont les valeurs s'étendent sur une très grande plage.

\subsection{Utilisation fiable de la calculatrice}

\begin{itemize}
\item Entrer les parenthèses dans le changement de base : $\ln(7)/\ln(2)$.
\item Vérifier que l'argument d'un logarithme est positif.
\item Conserver plusieurs chiffres avant l'arrondi final.
\item Tester le résultat dans l'équation initiale.
\item Comparer avec une estimation : si $2^3<7<2^4$, alors $3<\log_2 7<4$.
\end{itemize}

\begin{summarybox}
Une exponentielle traduit un facteur multiplicatif constant. Le logarithme est son opération inverse et permet de résoudre les problèmes où l'inconnue est dans l'exposant. Les lois logarithmiques proviennent directement des lois des exposants et exigent des arguments positifs.
\end{summarybox}


\part{Dénombrement et probabilités}
% ============================================================

\chapter{Principes de dénombrement}

\begin{questionbox}
Comment compter efficacement un grand nombre de possibilités, et comment décider si l'ordre ou la répétition doivent être pris en compte?
\end{questionbox}

\section{Le principe fondamental du dénombrement}

\begin{intuitionbox}
Lorsqu'une procédure se déroule en plusieurs étapes indépendantes, chaque choix possible à une étape peut être combiné avec tous les choix de l'étape suivante. C'est pourquoi les nombres de possibilités se multiplient.
\end{intuitionbox}

\begin{theorem}{Principe multiplicatif}{}
Si une première étape peut être réalisée de $n_1$ façons, une deuxième de $n_2$ façons, et ainsi de suite, alors la procédure complète peut être réalisée de
\[
n_1n_2\cdots n_k
\]
façons.
\end{theorem}

\begin{examplebox}
Un code contient deux lettres suivies de trois chiffres. Si la répétition est permise, il existe
\[
26^2\cdot10^3=676\,000
\]
codes possibles.
\end{examplebox}

Un arbre de choix rend le principe visible.

\begin{center}
\begin{tikzpicture}[grow=right,level distance=2.4cm,sibling distance=8mm,
  every node/.style={font=\small},
  edge from parent/.style={draw,-{Stealth[length=2mm]}}]
\node {repas}
 child {node {poulet}
   child {node {riz}}
   child {node {salade}}
 }
 child {node {poisson}
   child {node {riz}}
   child {node {salade}}
 }
 child {node {tofu}
   child {node {riz}}
   child {node {salade}}
 };
\end{tikzpicture}
\end{center}
Il y a $3\times2=6$ repas possibles.

\section{Factorielle et permutations}

\begin{definition}{Factorielle}{}
Pour $n\in\N$,
\[
n!=n(n-1)\cdots2\cdot1,
\qquad 0!=1.
\]
\end{definition}

\begin{definition}{Permutation}{}
Une permutation de $n$ objets distincts est un ordre utilisant tous les objets. Le nombre de permutations est
\[
P_n=n!.
\]
\end{definition}

\begin{examplebox}
Six personnes peuvent s'asseoir sur une rangée de
\[
6!=720
\]
façons.
\end{examplebox}

Si certains objets sont identiques, les permutations qui ne diffèrent que par l'échange d'objets identiques doivent être regroupées.

\begin{proposition}{Permutations avec répétitions}{}
Si $n$ objets contiennent $n_1$ objets identiques d'un premier type, $n_2$ d'un deuxième type, etc., avec $n_1+\cdots+n_r=n$, le nombre d'ordres distincts est
\[
\frac{n!}{n_1!n_2!\cdots n_r!}.
\]
\end{proposition}

\begin{examplebox}
Le mot \textsc{BANANE} contient six lettres : trois A, deux N et un B. Le nombre d'anagrammes distinctes est
\[
\frac{6!}{3!2!}=60.
\]
\end{examplebox}

\section{Arrangements}

\begin{definition}{Arrangement sans répétition}{}
Un arrangement de $k$ objets choisis parmi $n$ objets distincts est une sélection ordonnée sans répétition. Son nombre est
\[
A_n^k=\frac{n!}{(n-k)!}.
\]
\end{definition}

\begin{examplebox}
Parmi dix finalistes, on attribue une médaille d'or, une d'argent et une de bronze. L'ordre compte et une personne ne reçoit qu'une médaille :
\[
A_{10}^3=10\cdot9\cdot8=720.
\]
\end{examplebox}

Avec répétition autorisée, le nombre de suites ordonnées de longueur $k$ formées à partir de $n$ symboles est $n^k$.

\section{Combinaisons}

\begin{definition}{Combinaison}{}
Une combinaison de $k$ objets choisis parmi $n$ objets distincts est une sélection où l'ordre ne compte pas. Son nombre est
\[
\binom nk=\frac{n!}{k!(n-k)!}.
\]
\end{definition}

\begin{intuitionbox}
L'arrangement $A_n^k$ compte chaque groupe de $k$ objets dans les $k!$ ordres possibles. Pour oublier l'ordre, on divise donc par $k!$.
\end{intuitionbox}

\begin{examplebox}
Former un comité de quatre personnes parmi douze donne
\[
\binom{12}{4}=495
\]
comités. Aucun poste particulier n'est attribué, donc l'ordre ne compte pas.
\end{examplebox}

\section{Critères de sélection d'une méthode}

\begin{methodbox}
Avant de calculer, poser trois questions :
\begin{enumerate}
\item Utilise-t-on tous les objets ou seulement une partie?
\item L'ordre change-t-il le résultat?
\item La répétition est-elle permise?
\end{enumerate}
\begin{center}
\begin{tabular}{lll}
\toprule
Situation & Ordre & Nombre \\
\midrule
Tous les objets distincts & compte & $n!$ \\
$k$ objets parmi $n$, sans répétition & compte & $A_n^k$ \\
$k$ objets parmi $n$, sans répétition & ne compte pas & $\binom nk$ \\
$k$ positions, $n$ choix par position & compte & $n^k$ \\
\bottomrule
\end{tabular}
\end{center}
\end{methodbox}

\begin{counterbox}
Choisir trois livres à lire et choisir quel livre sera lu en premier, deuxième et troisième sont deux problèmes différents. Le premier est une combinaison; le second est un arrangement.
\end{counterbox}

\section[Comptage par complément et inclusion-exclusion]{Comptage par complément\\et inclusion-exclusion}

Il est parfois plus simple de compter ce qu'on ne veut pas.

\begin{examplebox}
Combien de chaînes de cinq chiffres contiennent au moins un zéro? Il existe $10^5$ chaînes au total. Les chaînes sans zéro utilisent neuf chiffres à chaque position, donc il y en a $9^5$. Le nombre recherché est
\[
10^5-9^5=40\,951.
\]
\end{examplebox}

Pour deux ensembles finis $A$ et $B$,
\[
\abs{A\cup B}=\abs A+\abs B-\abs{A\cap B}.
\]
On soustrait l'intersection parce qu'elle a été comptée deux fois.


\section{Voir les choix comme un arbre}

\begin{visualbox}
Avant toute formule, il faut décrire une issue. Un arbre matérialise la succession des choix. Si chaque branche du premier niveau se prolonge par le même nombre de branches, le nombre total de feuilles est un produit. Si l'on rassemble des familles de feuilles incompatibles, leurs nombres s'additionnent.
\end{visualbox}

\begin{center}
\begin{tikzpicture}[grow=right,level distance=2.0cm,sibling distance=7mm,
 every node/.style={font=\small},edge from parent/.style={draw,teal,-{Stealth[length=2mm]}}]
\node[draw=navy,rounded corners,fill=mistblue] {tenue}
 child {node {chemise A}
   child {node {pantalon 1}}
   child {node {pantalon 2}}
   child {node {pantalon 3}}
 }
 child {node {chemise B}
   child {node {pantalon 1}}
   child {node {pantalon 2}}
   child {node {pantalon 3}}
 };
\end{tikzpicture}
\end{center}

\begin{center}
\resizebox{0.88\textwidth}{!}{%
\begin{tikzpicture}[node distance=0.65cm and 0.9cm,every node/.style={font=\small,rounded corners,inner sep=5pt}]
 \node[draw=navy,fill=mistblue] (q1) {Utilise-t-on tous les objets?};
 \node[draw=teal,fill=mistteal,below left=of q1] (all) {oui};
 \node[draw=teal,fill=mistteal,below right=of q1] (some) {non};
 \node[draw=gold,fill=mistgold,below=of some] (q2) {L'ordre distingue-t-il les résultats?};
 \node[draw=navy,fill=mistblue,below=of all] (perm) {permutation};
 \node[draw=navy,fill=mistblue,below left=of q2] (arr) {arrangement};
 \node[draw=navy,fill=mistblue,below right=of q2] (comb) {combinaison};
 \draw[->,thick] (q1)--(all); \draw[->,thick] (q1)--(some); \draw[->,thick] (all)--(perm);
 \draw[->,thick] (some)--(q2); \draw[->,thick] (q2)--node[left]{oui}(arr); \draw[->,thick] (q2)--node[right]{non}(comb);
\end{tikzpicture}}
\end{center}

\begin{connectionbox}
La division par $k!$ dans $\binom nk$ corrige une surcomptabilisation : chaque groupe de $k$ objets apparaît sous $k!$ ordres différents parmi les arrangements. La formule n'est donc pas à mémoriser isolément; elle raconte exactement ce que l'on décide d'ignorer.
\end{connectionbox}

\subsection{Trois contrôles rapides}
Un résultat de dénombrement doit être entier, positif et compatible avec un cas simple. Pour $k=1$, on doit retrouver $\binom n1=n$; pour $k=n$, une seule sélection est possible; et $\binom nk=\binom n{n-k}$ exprime que choisir les objets retenus revient à choisir ceux que l'on laisse.


\section{Pourquoi l'ordre change le comptage}

La même collection d'objets peut conduire à des réponses très différentes selon ce que l'on considère comme une issue distincte.

\subsection{Sélectionner trois personnes}
Avec les personnes $A,B,C,D$, les six écritures
\[
ABC,\,ACB,\,BAC,\,BCA,\,CAB,\,CBA
\]
représentent six arrangements si les positions sont différentes, mais un seul groupe si l'ordre n'a pas de rôle. Chaque groupe de trois personnes est donc compté $3!=6$ fois parmi les arrangements.

\begin{center}
\begin{tikzpicture}[node distance=0.55cm,every node/.style={font=\small,rounded corners,inner sep=5pt}]
 \node[draw=navy,fill=mistblue] (abc) {$ABC$};
 \node[draw=navy,fill=mistblue,right=of abc] (acb) {$ACB$};
 \node[draw=navy,fill=mistblue,right=of acb] (bac) {$BAC$};
 \node[draw=navy,fill=mistblue,right=of bac] (bca) {$BCA$};
 \node[draw=navy,fill=mistblue,right=of bca] (cab) {$CAB$};
 \node[draw=navy,fill=mistblue,right=of cab] (cba) {$CBA$};
 \draw[teal,thick,decorate,decoration={brace,amplitude=6pt,mirror}] (abc.south west) -- (cba.south east)
   node[midway,below=9pt] {un même groupe $\{A,B,C\}$};
\end{tikzpicture}
\end{center}

Ainsi,
\[
\binom n3=\frac{n(n-1)(n-2)}{3!}.
\]
Le dénominateur ne vient pas d'une convention : il retire les permutations internes que le problème ne distingue pas.

\subsection{Compter le complément}
Il est parfois plus simple de compter les résultats interdits. Pour des codes de quatre chiffres contenant au moins un zéro, compter directement oblige à séparer plusieurs cas. Le complément « aucun zéro » possède $9^4$ codes parmi $10^4$ codes au total, donc
\[
10^4-9^4
\]
codes contiennent au moins un zéro.

\begin{visualbox}
Avant d'écrire une formule, compléter une phrase précise :
\begin{quote}
Une issue est déterminée par \dots; deux issues sont différentes lorsque \dots; la répétition est \dots; l'ordre est \dots.
\end{quote}
Cette phrase contient souvent toute la méthode de dénombrement.
\end{visualbox}


\section{Construire un comptage fiable}

Le dénombrement ne commence pas par une formule, mais par une description précise des décisions successives. Il faut identifier les étapes, les rôles, l'ordre, la possibilité de répétition et les restrictions. Une formule correcte appliquée au mauvais modèle donne un résultat incorrect.

\subsection{Principe additif et principe multiplicatif}

Le principe multiplicatif s'applique lorsque toutes les étapes doivent être réalisées. Si une tenue est formée d'un haut, d'un pantalon et d'une paire de chaussures, les nombres de choix se multiplient.

Le principe additif s'applique à des cas incompatibles. Si un déplacement se fait soit en bus selon $4$ lignes, soit en métro selon $3$ itinéraires, il existe $4+3=7$ choix, à condition qu'un trajet ne soit pas compté dans les deux catégories.

\begin{examplebox}
Une plaque contient deux lettres suivies de trois chiffres. Sans restriction et avec répétition,
\[
26^2\cdot10^3
\]
plaques sont possibles. Le facteur $26$ apparaît deux fois parce que les deux positions de lettre sont des étapes distinctes.
\end{examplebox}

\subsection{Arbre de choix et validation sur un cas élémentaire}

Un arbre rend visibles les étapes et aide à décider si le nombre de choix reste constant. Pour trois lancers d'une pièce, chaque nœud possède deux branches; il y a
\[
2^3=8
\]
suites. Pour choisir deux personnes sans remise parmi cinq, le premier choix offre $5$ possibilités et le second $4$, donc $5\cdot4=20$ sélections ordonnées.

Avant d'utiliser une formule générale, un petit cas sert de contrôle. Avec trois objets $A,B,C$, les arrangements de deux objets sont
\[
AB,AC,BA,BC,CA,CB,
\]
soit $6=3\cdot2$.

\subsection{Permutations : tous les objets, ordre important}

Le nombre de façons d'ordonner $n$ objets distincts est
\[
n!=n(n-1)\cdots2\cdot1.
\]
Le premier emplacement peut recevoir $n$ objets, le deuxième $n-1$, etc.

Si certains objets sont identiques, les permutations visibles sont moins nombreuses. Le mot \emph{BANANE} contient six lettres : B, A, N, A, N, E. Les deux A sont identiques, tout comme les deux N. Le nombre d'ordres distincts est
\[
\frac{6!}{2!2!}=180.
\]
On divise par les permutations internes des lettres identiques, qui ne produisent aucune nouvelle disposition visible.

\subsection{Arrangements : choisir et attribuer des rôles}

Choisir un président, un vice-président et un secrétaire parmi $10$ personnes est un choix ordonné de $3$ personnes :
\[
A_{10}^3=10\cdot9\cdot8=720.
\]
Les rôles rendent l'ordre important. Le groupe $(A,B,C)$ n'est pas équivalent à $(B,A,C)$, car les responsabilités changent.

\subsection{Combinaisons : choisir sans ordre}

Choisir un comité de $3$ personnes parmi $10$ ne tient pas compte de l'ordre. Les $3!$ ordres d'un même groupe ont été comptés dans les arrangements; on divise donc par $3!$ :
\[
\binom{10}{3}=\frac{10\cdot9\cdot8}{3\cdot2\cdot1}=120.
\]
Cette explication est plus importante que la formule elle-même : la division élimine les répétitions artificielles créées par l'ordre.

\subsection{Symétrie des coefficients binomiaux}

Choisir $k$ personnes qui participent revient à choisir les $n-k$ personnes qui ne participent pas. Ainsi,
\[
\binom nk=\binom n{n-k}.
\]
Cette symétrie permet souvent de simplifier les calculs. Par exemple,
\[
\binom{20}{18}=\binom{20}{2}=190.
\]

\subsection{Comptage par complément}

Les expressions « au moins un » ou « pas tous » sont souvent plus simples par complément.

\begin{examplebox}
Un code contient deux lettres suivies de quatre chiffres. Les répétitions sont permises, mais au moins un chiffre doit être impair. Le nombre total est
\[
26^2\cdot10^4.
\]
Le complément contient seulement des chiffres pairs :
\[
26^2\cdot5^4.
\]
Le nombre admissible est donc
\[
26^2(10^4-5^4)=6\,337\,500.
\]
\end{examplebox}

\subsection{Inclusion-exclusion}

Lorsque deux catégories se chevauchent,
\[
\abs{A\cup B}=\abs A+\abs B-\abs{A\cap B}.
\]
La soustraction corrige les éléments comptés deux fois.

Supposons que, parmi $80$ personnes, $45$ parlent français, $38$ parlent anglais et $20$ parlent les deux langues. Le nombre qui parle au moins une des deux langues est
\[
45+38-20=63.
\]
Il reste $17$ personnes qui ne parlent aucune des deux selon les données.

\subsection{Guide de décision}

\begin{center}
\begin{tabularx}{\textwidth}{@{}lY@{}}
\toprule
Structure de la situation & Outil naturel\\\midrule
plusieurs étapes successives & principe multiplicatif\\
cas incompatibles & principe additif\\
tous les objets distincts sont ordonnés & permutation\\
$k$ objets choisis avec rôles ou ordre & arrangement\\
$k$ objets choisis sans ordre & combinaison\\
condition « au moins un » & complément\\
deux propriétés qui se chevauchent & inclusion-exclusion\\
\bottomrule
\end{tabularx}
\end{center}

\subsection{Contrôles de plausibilité}

Le résultat ne peut pas dépasser le nombre de suites sans restriction. Une combinaison doit être inférieure ou égale à l'arrangement correspondant. La symétrie
\[
\binom nk=\binom n{n-k}
\]
fournit un contrôle immédiat. Enfin, l'énumération d'un petit cas permet de vérifier la structure choisie.

\begin{summarybox}
Compter exige d'abord de décrire les étapes et de décider si l'ordre et la répétition comptent. Les permutations, arrangements et combinaisons ne sont que des formes spécialisées du principe multiplicatif; le complément et l'inclusion-exclusion traitent les restrictions efficacement.
\end{summarybox}


\chapter{Modèles probabilistes élémentaires}

\begin{questionbox}
Comment traduire une expérience aléatoire en un modèle mathématique et calculer la probabilité d'un événement composé?
\end{questionbox}

\section{Expérience, univers et événement}

\begin{definition}{Univers et événement}{}
L'univers $\Omega$ est l'ensemble de toutes les issues possibles d'une expérience aléatoire. Un événement est un sous-ensemble $A\subseteq\Omega$.
\end{definition}

Pour un dé à six faces,
\[
\Omega=\{1,2,3,4,5,6\}.
\]
L'événement « obtenir un nombre pair » est $A=\{2,4,6\}$.

Le complément $A^c$ contient les issues où $A$ ne se réalise pas. L'union $A\cup B$ signifie « $A$ ou $B$ », et l'intersection $A\cap B$ signifie « $A$ et $B$ ».

\section{Axiomes et équiprobabilité}

Une probabilité $P$ satisfait :
\begin{align*}
0&\le P(A)\le1,\\
P(\Omega)&=1,\\
P(A\cup B)&=P(A)+P(B)\quad\text{si }A\cap B=\varnothing.
\end{align*}
On en déduit
\[
P(A^c)=1-P(A)
\]
et
\[
P(A\cup B)=P(A)+P(B)-P(A\cap B).
\]

\begin{theorem}{Cas équiprobable}{}
Si toutes les issues de $\Omega$ sont également probables et si $\Omega$ est fini, alors
\[
P(A)=\frac{\abs A}{\abs\Omega}.
\]
\end{theorem}

\begin{examplebox}
On lance deux dés distinguables. Il y a $6\cdot6=36$ couples équiprobables. Six couples ont une somme égale à $7$ :
\[
(1,6),(2,5),(3,4),(4,3),(5,2),(6,1).
\]
Ainsi,
\[
P(\text{somme }7)=\frac6{36}=\frac16.
\]
\end{examplebox}

\begin{warningbox}
Les sommes possibles $2,3,\ldots,12$ ne sont pas équiprobables. Par exemple, la somme $2$ provient d'un seul couple, tandis que la somme $7$ provient de six couples. Il faut choisir un univers dont les issues élémentaires sont réellement équiprobables.
\end{warningbox}

\section{Probabilité conditionnelle}

\begin{definition}{Probabilité conditionnelle}{}
Si $P(B)>0$, la probabilité de $A$ sachant que $B$ s'est réalisé est
\[
P(A\mid B)=\frac{P(A\cap B)}{P(B)}.
\]
\end{definition}

Cette formule restreint l'univers à $B$. Elle se réécrit
\[
P(A\cap B)=P(B)P(A\mid B).
\]

\begin{examplebox}
Un sac contient trois boules rouges et deux boules bleues. On tire deux boules sans remise. La probabilité de tirer deux rouges est
\[
P(R_1\cap R_2)=P(R_1)P(R_2\mid R_1)=\frac35\cdot\frac24=\frac3{10}.
\]
\end{examplebox}

\begin{center}
\begin{tikzpicture}[grow=right,level distance=2.7cm,sibling distance=12mm,
 every node/.style={font=\small},edge from parent/.style={draw,-{Stealth[length=2mm]}}]
\node {départ}
 child {node {$R$} edge from parent node[above] {$3/5$}
   child {node {$R$} edge from parent node[above] {$2/4$}}
   child {node {$B$} edge from parent node[below] {$2/4$}}
 }
 child {node {$B$} edge from parent node[below] {$2/5$}
   child {node {$R$} edge from parent node[above] {$3/4$}}
   child {node {$B$} edge from parent node[below] {$1/4$}}
 };
\end{tikzpicture}
\end{center}

Dans un arbre, on multiplie les probabilités le long d'une branche et on additionne les branches correspondant à des cas disjoints.

\section{Indépendance}

\begin{definition}{Événements indépendants}{}
Deux événements $A$ et $B$ sont indépendants si
\[
P(A\cap B)=P(A)P(B).
\]
Lorsque $P(B)>0$, cela équivaut à $P(A\mid B)=P(A)$.
\end{definition}

\begin{examplebox}
Deux lancers successifs d'une pièce équilibrée sont indépendants. La probabilité d'obtenir deux faces est
\[
\frac12\cdot\frac12=\frac14.
\]
En revanche, deux tirages sans remise dans un petit sac ne sont généralement pas indépendants, car le premier tirage modifie la composition du sac.
\end{examplebox}

\section{Épreuve de Bernoulli et modèle binomial}

\begin{definition}{Épreuve de Bernoulli}{}
Une épreuve de Bernoulli possède deux issues : succès, de probabilité $p$, et échec, de probabilité $1-p$.
\end{definition}

Si l'on répète indépendamment l'épreuve $n$ fois et si $X$ compte les succès, alors
\[
P(X=k)=\binom nk p^k(1-p)^{n-k}.
\]

\begin{intuitionbox}
Le facteur $p^k(1-p)^{n-k}$ donne la probabilité d'une suite particulière contenant $k$ succès. Le coefficient $\binom nk$ compte les positions possibles de ces succès.
\end{intuitionbox}

\begin{examplebox}
Une question à choix binaire est répondue au hasard cinq fois, indépendamment. La probabilité d'obtenir exactement trois bonnes réponses est
\[
\binom53\left(\frac12\right)^3\left(\frac12\right)^2
=\frac{10}{32}=\frac5{16}.
\]
\end{examplebox}

\section{Espérance mathématique}

\begin{definition}{Espérance d'une variable discrète}{}
Si une variable aléatoire $X$ prend les valeurs $x_1,\ldots,x_m$ avec probabilités $p_1,\ldots,p_m$, alors
\[
\mathbb E[X]=\sum_{i=1}^m x_i p_i.
\]
\end{definition}

L'espérance est une moyenne à long terme, pas nécessairement une valeur possible lors d'une seule expérience.

\begin{examplebox}
Pour un dé équilibré,
\[
\mathbb E[X]=\frac{1+2+3+4+5+6}{6}=3{,}5.
\]
On n'obtient jamais $3{,}5$ sur un lancer, mais la moyenne de nombreux lancers tend vers cette valeur.
\end{examplebox}


\section{Proportion, fréquence et probabilité}

\begin{visualbox}
Une probabilité mesure la part de l'univers occupée par un événement. Dans un modèle équiprobable, compter les issues favorables revient à mesurer une aire discrète. Le conditionnement ne change pas l'événement observé; il réduit l'univers aux issues compatibles avec l'information reçue.
\end{visualbox}

\begin{center}
\begin{tikzpicture}[scale=0.62]
  \foreach \i in {0,...,5}{\foreach \j in {0,...,5}{
    \pgfmathtruncatemacro{\s}{\i+\j+2}
    \ifnum\s=7 \fill[mistgold] (\i,\j) rectangle ++(1,1);\fi
    \draw[slate!60] (\i,\j) rectangle ++(1,1);
  }}
  \node[below] at (3,-0.35) {premier dé};
  \node[rotate=90] at (-0.55,3) {second dé};
  \node[gold!75!black] at (7.6,3.4) {$6$ cases sur $36$};
  \node[gold!75!black] at (7.6,2.7) {$P(\text{somme }7)=1/6$};
\end{tikzpicture}
\end{center}

\begin{center}
\begin{tikzpicture}[every node/.style={font=\small,rounded corners,inner sep=5pt,align=center}]
 \node[draw=navy,fill=mistblue] (p) at (0,0) {population};
 \node[draw=teal,fill=mistteal] (m) at (3,1.25) {malade};
 \node[draw=teal,fill=mistteal] (n) at (3,-1.25) {non malade};
 \node[draw=gold,fill=mistgold] (mp) at (6,2.05) {test $+$};
 \node[draw=gold,fill=mistgold] (mm) at (6,0.45) {test $-$};
 \node[draw=gold,fill=mistgold] (np) at (6,-0.45) {test $+$};
 \node[draw=gold,fill=mistgold] (nm) at (6,-2.05) {test $-$};
 \draw[->,thick,teal] (p)--node[above left] {$0{,}02$}(m);
 \draw[->,thick,teal] (p)--node[below left] {$0{,}98$}(n);
 \draw[->,thick,gold!80!black] (m)--node[above] {$0{,}95$}(mp);
 \draw[->,thick,gold!80!black] (m)--node[below] {$0{,}05$}(mm);
 \draw[->,thick,gold!80!black] (n)--node[above] {$0{,}04$}(np);
 \draw[->,thick,gold!80!black] (n)--node[below] {$0{,}96$}(nm);
\end{tikzpicture}
\end{center}

\begin{connectionbox}
Sur un arbre, on \textbf{multiplie le long d'une branche} parce que les étapes se succèdent; on \textbf{additionne les branches compatibles avec l'événement} parce qu'elles représentent des façons distinctes de le réaliser. Cette règle relie directement dénombrement et probabilité.
\end{connectionbox}

\subsection{Indépendance n'est pas incompatibilité}
Deux événements incompatibles ne peuvent pas se produire ensemble. Deux événements indépendants peuvent se produire ensemble, mais connaître l'un ne change pas la probabilité de l'autre. La différence se voit dans les formules : incompatibilité donne $P(A\cap B)=0$, indépendance donne $P(A\cap B)=P(A)P(B)$.


\section{Le conditionnement avec des effectifs concrets}

Les probabilités conditionnelles sont souvent plus faciles à comprendre avec une population fictive qu'avec des pourcentages isolés.

Supposons une maladie présente chez $2\%$ des personnes. Un test détecte $95\%$ des personnes malades et donne un résultat faussement positif chez $4\%$ des personnes non malades. Sur $10\,000$ personnes :
\begin{itemize}
\item $200$ sont malades; environ $190$ ont un test positif;
\item $9\,800$ ne sont pas malades; environ $392$ ont un test positif.
\end{itemize}
Parmi les $582$ tests positifs, seulement $190$ correspondent à des personnes malades. Ainsi,
\[
P(\text{malade}\mid +)=\frac{190}{582}\approx0{,}326.
\]
Un test positif n'implique donc pas une probabilité de maladie de $95\%$.

\begin{center}
\begin{tikzpicture}[scale=0.95]
 \draw[slate,thick] (0,0) rectangle (8,4.5);
 \draw[teal,thick] (0,0)--(0.16,4.5);
 \fill[mistred] (0,0) rectangle (0.16,4.275);
 \fill[mistgold] (0.16,0) rectangle (8,0.18);
 \node[rotate=90,terracotta] at (0.08,2.2) {malades};
 \node[gold!75!black] at (4.1,0.09) {faux positifs};
 \node[navy,align=left] at (4.3,2.5) {la grande taille du groupe non malade\\peut produire beaucoup de faux positifs};
\end{tikzpicture}
\end{center}

\subsection{Pourquoi la probabilité de base compte}
Le taux de maladie initial, appelé probabilité de base, détermine la taille des deux branches avant même le test. Un test très précis peut produire une proportion importante de faux positifs lorsque l'événement recherché est rare.

\subsection{Indépendance : une égalité à vérifier}
Dire que $A$ et $B$ sont indépendants signifie
\[
P(A\mid B)=P(A),
\]
à condition que $P(B)>0$. Cette définition se traduit par $P(A\cap B)=P(A)P(B)$. L'indépendance n'est jamais conclue seulement parce que deux événements « semblent sans rapport ».

\begin{connectionbox}
Le conditionnement est un changement de dénominateur. Dans $P(A)=|A|/|\Omega|$, on compare $A$ à tout l'univers. Dans $P(A\mid B)=|A\cap B|/|B|$, on compare la partie de $A$ située dans $B$ au nouvel univers $B$.
\end{connectionbox}



\subsection{Tableau à double entrée ou arbre?}
Les deux représentations contiennent la même information, mais elles mettent l'accent sur des aspects différents. Un arbre suit l'ordre temporel ou logique des étapes. Un tableau à double entrée compare directement les intersections de deux catégories.

\begin{center}
\begin{tabular}{c|cc|c}
 & $B$ & $B^c$ & total\\\hline
$A$ & $n(A\cap B)$ & $n(A\cap B^c)$ & $n(A)$\\
$A^c$ & $n(A^c\cap B)$ & $n(A^c\cap B^c)$ & $n(A^c)$\\\hline
total & $n(B)$ & $n(B^c)$ & $n(\Omega)$
\end{tabular}
\end{center}

Dans ce tableau,
\[
P(A\mid B)=\frac{n(A\cap B)}{n(B)}.
\]
La formule se lit littéralement : parmi les individus de la colonne $B$, quelle proportion appartient aussi à la ligne $A$?

\subsection{Simuler pour contrôler, non pour remplacer le raisonnement}
Une simulation informatique peut vérifier qu'une probabilité théorique est plausible. Elle produit toutefois une fréquence variable, tandis que le modèle exact produit une probabilité. Avec peu d'essais, une fréquence peut s'éloigner fortement de la valeur théorique; avec beaucoup d'essais, elle tend à se stabiliser autour d'elle.

\begin{visualbox}
Lorsqu'un modèle probabiliste donne une réponse surprenante, trois contrôles sont particulièrement utiles : dessiner l'univers, refaire le calcul avec des effectifs entiers et simuler un grand nombre d'expériences. Ces contrôles éclairent le raisonnement; ils ne dispensent pas de définir correctement l'expérience.
\end{visualbox}


\section{Interpréter une information probabiliste}

Une probabilité n'est pas seulement un nombre entre $0$ et $1$. Elle dépend d'un modèle : quelles sont les issues possibles, lesquelles sont considérées comme élémentaires, et comment leur attribue-t-on des probabilités? Une description incorrecte de l'univers produit des calculs trompeurs même si les formules sont appliquées correctement.

\subsection{Construire l'univers}

Pour deux lancers d'une pièce, l'univers ordonné est
\[
\Omega=\{PP,PF,FP,FF\}.
\]
Les quatre issues sont équiprobables pour une pièce équilibrée. En revanche, si l'on décrit seulement le nombre de piles, les résultats $0,1,2$ ne sont pas équiprobables : « un pile » correspond à deux issues élémentaires.

Cette distinction explique pourquoi la formule
\[
P(A)=\frac{\abs A}{\abs\Omega}
\]
n'est valide que lorsque les issues élémentaires de $\Omega$ sont équiprobables.

\subsection{Événements et opérations ensemblistes}

Un événement est un sous-ensemble de $\Omega$. Le complément $A^c$ signifie que $A$ ne se produit pas; l'union $A\cup B$ signifie qu'au moins un des événements se produit; l'intersection $A\cap B$ signifie qu'ils se produisent ensemble.

Les règles fondamentales sont
\[
P(A^c)=1-P(A),
\]
\[
P(A\cup B)=P(A)+P(B)-P(A\cap B).
\]
La soustraction de l'intersection évite le double comptage.

\subsection{Tirage sans remise : les probabilités changent}

Une urne contient cinq boules rouges et trois bleues. On tire deux boules sans remise. La probabilité de deux rouges est
\[
\frac58\cdot\frac47=\frac5{14}.
\]
La probabilité d'une boule de chaque couleur possède deux chemins :
\[
\frac58\cdot\frac37+\frac38\cdot\frac57
=\frac{15}{56}+\frac{15}{56}
=\frac{30}{56}=\frac{15}{28}.
\]
L'arbre permet de voir les deux ordres et le changement des dénominateurs après le premier tirage.

\subsection{Probabilité conditionnelle}

La probabilité
\[
P(A\mid B)=\frac{P(A\cap B)}{P(B)}
\]
se lit « probabilité de $A$ sachant $B$ ». L'information $B$ réduit l'univers de référence aux situations où $B$ s'est produit.

\begin{examplebox}
Dans un groupe, $60\%$ des personnes utilisent le transport collectif et $25\%$ utilisent à la fois le transport collectif et le vélo. Parmi les personnes utilisant le transport collectif, la proportion qui utilise aussi le vélo est
\[
P(V\mid T)=\frac{P(V\cap T)}{P(T)}=\frac{0{,}25}{0{,}60}=\frac5{12}.
\]
Le dénominateur est la population conditionnelle, non la population totale.
\end{examplebox}

\subsection{Indépendance et incompatibilité}

Deux événements sont indépendants lorsque
\[
P(A\cap B)=P(A)P(B).
\]
L'information sur l'un ne modifie alors pas la probabilité de l'autre. Deux événements incompatibles satisfont $A\cap B=\varnothing$; ils ne peuvent pas se produire ensemble.

Ces notions sont différentes. Si $A$ et $B$ sont incompatibles et ont tous deux une probabilité positive, ils ne sont pas indépendants, car
\[
P(A\cap B)=0\neq P(A)P(B).
\]

\subsection{Arbres de probabilités}

Dans un arbre :
\begin{itemize}
\item les probabilités issues d'un même nœud totalisent $1$;
\item on multiplie les probabilités le long d'un chemin;
\item on additionne les chemins correspondant au même événement final.
\end{itemize}

Cette règle est une traduction visuelle de la probabilité conditionnelle :
\[
P(A\cap B)=P(A)P(B\mid A).
\]

\subsection{Essais répétés élémentaires}

Lorsqu'une expérience possède deux résultats, succès ou échec, avec une probabilité de succès $p$ constante et des essais indépendants, la probabilité d'une suite précise contenant $k$ succès et $n-k$ échecs est
\[
p^k(1-p)^{n-k}.
\]
Il existe
\[
\binom nk
\]
positions possibles pour les $k$ succès. Ainsi, la probabilité d'obtenir exactement $k$ succès est
\[
\binom nkp^k(1-p)^{n-k}.
\]
Cette formule est une application directe du dénombrement, et non une nouvelle règle sans lien avec le chapitre précédent.

\begin{examplebox}
Une composante a une probabilité $0{,}9$ de fonctionner. Pour quatre composantes indépendantes, la probabilité que exactement trois fonctionnent est
\[
\binom43(0{,}9)^3(0{,}1)=4(0{,}729)(0{,}1)=0{,}2916.
\]
\end{examplebox}

\subsection{Interprétation de l'espérance}

Si une variable aléatoire $X$ prend les valeurs $x_i$ avec probabilités $p_i$, sa valeur moyenne théorique est
\[
E(X)=\sum_i x_ip_i.
\]
Elle ne prédit pas nécessairement le résultat d'une expérience unique; elle décrit une moyenne à long terme.

Pour un jeu qui rapporte $8\,$\$ avec probabilité $1/4$ et fait perdre $3\,$\$ sinon,
\[
E(X)=8\cdot\frac14-3\cdot\frac34=-\frac14.
\]
À long terme, la perte moyenne est de $0{,}25\,$\$ par partie.

\subsection{Diagnostic d'un modèle}

Avant de calculer :
\begin{enumerate}
\item définir clairement l'univers;
\item vérifier l'équiprobabilité ou donner les probabilités individuelles;
\item décider si l'ordre des étapes est pertinent;
\item identifier le remplacement ou l'absence de remplacement;
\item distinguer indépendance et incompatibilité;
\item vérifier que le résultat appartient à $[0,1]$.
\end{enumerate}

\subsection{Contrôle par complément}

Pour « au moins un succès » dans $n$ essais indépendants,
\[
P(\text{au moins un})=1-P(\text{aucun}).
\]
Avec une probabilité de succès $p$,
\[
P(\text{au moins un})=1-(1-p)^n.
\]
Cette voie est plus courte que l'addition de tous les cas contenant un ou plusieurs succès.

\begin{summarybox}
Une probabilité dépend d'un univers bien construit. Le conditionnement change le groupe de référence, l'indépendance signifie absence d'influence probabiliste et les arbres traduisent les étapes successives. Le dénombrement explique les modèles d'essais répétés.
\end{summarybox}


\part{Vecteurs, matrices et optimisation}
% ============================================================

\chapter{Vecteurs dans \texorpdfstring{$\R^2$ et $\R^3$}{R2 et R3}}

\begin{questionbox}
Comment représenter mathématiquement un déplacement qui possède à la fois une longueur et une direction?
\end{questionbox}

\section{Du point au vecteur}

\begin{intuitionbox}
Un point indique une position. Un vecteur indique un déplacement. Le même déplacement peut être dessiné à plusieurs endroits : ce qui compte est la différence entre le point d'arrivée et le point de départ.
\end{intuitionbox}

\begin{definition}{Vecteur entre deux points}{}
Dans $\R^2$, si $A(x_A,y_A)$ et $B(x_B,y_B)$, alors
\[
\vect{AB}=(x_B-x_A,\,y_B-y_A).
\]
Dans $\R^3$, si $A(x_A,y_A,z_A)$ et $B(x_B,y_B,z_B)$, alors
\[
\vect{AB}=(x_B-x_A,\,y_B-y_A,\,z_B-z_A).
\]
\end{definition}

\begin{center}
\begin{tikzpicture}[scale=0.8]
\draw[-{Stealth},thick] (-1,0)--(6,0) node[right] {$x$};
\draw[-{Stealth},thick] (0,-1)--(0,5) node[above] {$y$};
\fill (1,1) circle (2pt) node[below left] {$A$};
\fill (5,4) circle (2pt) node[above right] {$B$};
\draw[very thick,uqamblue,-{Stealth[length=3mm]}] (1,1)--(5,4);
\draw[dashed] (1,1)--(5,1)--(5,4);
\node at (3,0.7) {$4$};
\node at (5.35,2.5) {$3$};
\node[teal] at (3.1,2.9) {$\vect{AB}=(4,3)$};
\end{tikzpicture}
\end{center}

\section{Opérations sur les vecteurs}

Pour $\mathbf u=(u_1,u_2,u_3)$, $\mathbf v=(v_1,v_2,v_3)$ et $\lambda\in\R$ :
\begin{align*}
\mathbf u+\mathbf v&=(u_1+v_1,u_2+v_2,u_3+v_3),\\
\mathbf u-\mathbf v&=(u_1-v_1,u_2-v_2,u_3-v_3),\\
\lambda\mathbf u&=(\lambda u_1,\lambda u_2,\lambda u_3).
\end{align*}
L'addition se visualise par la règle du triangle ou du parallélogramme.

\begin{center}
\begin{tikzpicture}[scale=0.8]
\coordinate (O) at (0,0);
\coordinate (U) at (3,1);
\coordinate (V) at (1,2.5);
\coordinate (S) at (4,3.5);
\draw[very thick,uqamblue,-{Stealth[length=3mm]}] (O)--(U) node[midway,below] {$\mathbf u$};
\draw[very thick,teal,-{Stealth[length=3mm]}] (O)--(V) node[midway,left] {$\mathbf v$};
\draw[very thick,teal,-{Stealth[length=3mm]}] (U)--(S);
\draw[very thick,uqamblue,-{Stealth[length=3mm]}] (V)--(S);
\draw[very thick,deepblue,-{Stealth[length=3mm]}] (O)--(S) node[midway,above left] {$\mathbf u+\mathbf v$};
\end{tikzpicture}
\end{center}

\begin{examplebox}
Si $\mathbf u=(2,-1,3)$ et $\mathbf v=(-4,5,1)$, alors
\[
\mathbf u+\mathbf v=(-2,4,4),
\qquad
2\mathbf u-\mathbf v=(8,-7,5).
\]
\end{examplebox}

\section{Norme, distance et vecteur unitaire}

\begin{definition}{Norme euclidienne}{}
La norme de $\mathbf v=(v_1,v_2)$ ou $(v_1,v_2,v_3)$ est
\[
\norm{\mathbf v}=\sqrt{v_1^2+v_2^2}
\quad\text{ou}\quad
\norm{\mathbf v}=\sqrt{v_1^2+v_2^2+v_3^2}.
\]
\end{definition}

La distance entre deux points est $\norm{\vect{AB}}$. Un vecteur unitaire dans la direction de $\mathbf v\neq\mathbf0$ est
\[
\widehat{\mathbf v}=\frac{\mathbf v}{\norm{\mathbf v}}.
\]

\begin{examplebox}
Pour $\mathbf v=(3,4)$,
\[
\norm{\mathbf v}=5,
\qquad
\widehat{\mathbf v}=\left(\frac35,\frac45\right).
\]
\end{examplebox}

\section{Produit scalaire et angle}

\begin{definition}{Produit scalaire}{}
Pour $\mathbf u,\mathbf v\in\R^n$,
\[
\mathbf u\cdot\mathbf v=u_1v_1+\cdots+u_nv_n.
\]
Il satisfait aussi
\[
\mathbf u\cdot\mathbf v=\norm{\mathbf u}\norm{\mathbf v}\cos\theta,
\]
où $\theta$ est l'angle entre les deux vecteurs non nuls.
\end{definition}

Ainsi,
\[
\cos\theta=\frac{\mathbf u\cdot\mathbf v}{\norm{\mathbf u}\norm{\mathbf v}}.
\]
Deux vecteurs non nuls sont perpendiculaires exactement lorsque leur produit scalaire est nul.

\begin{examplebox}
Pour $\mathbf u=(1,2,-1)$ et $\mathbf v=(2,0,2)$,
\[
\mathbf u\cdot\mathbf v=2+0-2=0.
\]
Les vecteurs sont orthogonaux.
\end{examplebox}

\section{Combinaisons linéaires}

Une expression
\[
a\mathbf u+b\mathbf v
\]
est une combinaison linéaire de $\mathbf u$ et $\mathbf v$. Dans le plan, deux vecteurs non parallèles permettent d'atteindre tout vecteur. Dans l'espace, trois vecteurs non coplanaires peuvent jouer le même rôle.

\begin{examplebox}
Avec $\mathbf i=(1,0)$ et $\mathbf j=(0,1)$,
\[
(x,y)=x\mathbf i+y\mathbf j.
\]
Les coordonnées sont donc les coefficients de la décomposition dans la base canonique.
\end{examplebox}


\section{Interprétation géométrique : un déplacement libre}

\begin{visualbox}
Deux flèches placées à des endroits différents représentent le même vecteur si elles ont les mêmes composantes. Le vecteur ne mémorise pas son point de départ : il mémorise seulement le déplacement horizontal, vertical et, dans l'espace, le déplacement selon $z$.
\end{visualbox}

\begin{center}
\begin{tikzpicture}[scale=0.85]
  \draw[-{Stealth},thick,slate] (-1,0)--(7,0) node[right] {$x$};
  \draw[-{Stealth},thick,slate] (0,-1)--(0,5.5) node[above] {$y$};
  \draw[navy,very thick,-{Stealth[length=3mm]}] (0.8,0.8)--(3.8,2.8);
  \draw[teal,very thick,-{Stealth[length=3mm]}] (2.5,2.6)--(5.5,4.6);
  \draw[dashed,slate] (0.8,0.8)--(3.8,0.8)--(3.8,2.8);
  \draw[dashed,slate] (2.5,2.6)--(5.5,2.6)--(5.5,4.6);
  \node[navy] at (2.1,2.3) {$(3,2)$};
  \node[teal] at (4.0,4.4) {$(3,2)$};
\end{tikzpicture}
\end{center}

\begin{center}
\begin{tikzpicture}[scale=0.85]
  \coordinate (O) at (0,0); \coordinate (U) at (4,0.8); \coordinate (V) at (1.5,3.0);
  \draw[navy,very thick,-{Stealth[length=3mm]}] (O)--(U) node[midway,below] {$\mathbf u$};
  \draw[teal,very thick,-{Stealth[length=3mm]}] (O)--(V) node[midway,left] {$\mathbf v$};
  \draw[gold,thick,dashed] (V)--($(U)+(V)$)--(U);
  \draw[terracotta,very thick,-{Stealth[length=3mm]}] (O)--($(U)+(V)$) node[midway,above left] {$\mathbf u+\mathbf v$};
\end{tikzpicture}
\qquad
\begin{tikzpicture}[scale=0.85]
  \coordinate (O) at (0,0); \coordinate (U) at (4,0.7); \coordinate (V) at (2.1,2.8);
  \draw[navy,very thick,-{Stealth[length=3mm]}] (O)--(U) node[right] {$\mathbf u$};
  \draw[teal,very thick,-{Stealth[length=3mm]}] (O)--(V) node[above] {$\mathbf v$};
  \draw[gold,thick,dashed] (V)--($(O)!(V)!(U)$);
  \fill[gold] ($(O)!(V)!(U)$) circle (2pt);
  \node[gold!75!black,below] at ($(O)!(V)!(U)$) {projection};
\end{tikzpicture}
\end{center}

\begin{connectionbox}
Les composantes permettent le calcul; la flèche donne le sens géométrique. Le produit scalaire relie les deux :
\[
\mathbf u\cdot\mathbf v=u_1v_1+u_2v_2
=\norm{\mathbf u}\norm{\mathbf v}\cos\theta.
\]
Il mesure la part d'un vecteur orientée dans la direction de l'autre. Un produit scalaire nul signifie qu'aucune composante ne pointe dans cette direction : les vecteurs sont perpendiculaires.
\end{connectionbox}


\section{Le produit scalaire comme longueur d'une ombre}

Le produit scalaire mesure combien un vecteur avance dans la direction d'un autre. Si $\mathbf e$ est un vecteur unitaire, alors $\mathbf u\cdot\mathbf e$ est la longueur signée de la projection de $\mathbf u$ sur cette direction.

\begin{center}
\begin{tikzpicture}[scale=0.95]
 \coordinate (O) at (0,0); \coordinate (U) at (4.5,3.1); \coordinate (E) at (5.5,0.8);
 \draw[navy,very thick,-{Stealth[length=3mm]}] (O)--(U) node[above] {$\mathbf u$};
 \draw[teal,very thick,-{Stealth[length=3mm]}] (O)--(E) node[right] {$\mathbf e$};
 \coordinate (H) at ($(O)!(U)!(E)$);
 \draw[gold,thick,dashed] (U)--(H);
 \draw[terracotta,very thick] (O)--(H);
 \fill[terracotta] (H) circle (2.5pt);
 \node[terracotta,below] at ($(O)!0.5!(H)$) {$\mathbf u\cdot\mathbf e$};
 \draw pic[draw=gold,angle radius=7mm] {angle=E--O--U};
 \node[gold!75!black] at (1.25,0.65) {$\theta$};
\end{tikzpicture}
\end{center}

Comme la longueur de cette projection vaut $\norm{\mathbf u}\cos\theta$, on obtient
\[
\mathbf u\cdot\mathbf v
=\norm{\mathbf v}\bigl(\norm{\mathbf u}\cos\theta\bigr)
=\norm{\mathbf u}\norm{\mathbf v}\cos\theta.
\]

\subsection{Pourquoi la formule en coordonnées fonctionne}
Dans une base orthonormée,
\[
\mathbf u=u_1\mathbf e_1+u_2\mathbf e_2,
\qquad
\mathbf v=v_1\mathbf e_1+v_2\mathbf e_2.
\]
La distributivité et les relations $\mathbf e_1\cdot\mathbf e_1=1$, $\mathbf e_2\cdot\mathbf e_2=1$, $\mathbf e_1\cdot\mathbf e_2=0$ donnent
\[
\mathbf u\cdot\mathbf v=u_1v_1+u_2v_2.
\]
La formule algébrique est donc la somme des contributions selon deux directions perpendiculaires.

\subsection{Signe et angle}
Si le produit scalaire est positif, l'angle est aigu et les vecteurs pointent globalement dans le même sens. S'il est nul, ils sont orthogonaux. S'il est négatif, l'angle est obtus et une partie du déplacement s'oppose à la direction choisie.

\begin{visualbox}
Pour vérifier un angle calculé, regarder d'abord le signe du produit scalaire. Une calculatrice donnant $37^\circ$ alors que le produit scalaire est négatif signale une erreur de saisie ou de formule : l'angle devrait être obtus.
\end{visualbox}


\section{Relier déplacements et coordonnées}

Un point indique une position; un vecteur décrit un déplacement. Deux flèches placées à des endroits différents représentent le même vecteur si elles ont les mêmes composantes. Cette indépendance du point de départ permet d'utiliser les vecteurs pour la navigation, les forces et la géométrie analytique.

\subsection{Points et vecteurs : ne pas confondre les rôles}

Pour deux points $A(x_A,y_A)$ et $B(x_B,y_B)$,
\[
\vect{AB}=(x_B-x_A,\,y_B-y_A).
\]
Le vecteur contient l'information nécessaire pour passer de $A$ à $B$. Si
\[
A=(-2,1),\qquad B=(3,4),
\]
alors
\[
\vect{AB}=(5,3).
\]
À partir d'un autre point $C$, ajouter le même vecteur produit une translation parallèle de même longueur.

Dans $\R^3$, la même idée donne
\[
\vect{AB}=(x_B-x_A,\,y_B-y_A,\,z_B-z_A).
\]
La troisième composante ne modifie pas les règles algébriques; elle ajoute une direction spatiale.

\subsection{Addition : enchaîner des déplacements}

Si
\[
\mathbf u=(u_1,u_2),
\qquad
\mathbf v=(v_1,v_2),
\]
alors
\[
\mathbf u+\mathbf v=(u_1+v_1,u_2+v_2).
\]
Géométriquement, on place l'origine de $\mathbf v$ à l'extrémité de $\mathbf u$. Le vecteur résultant relie le point de départ initial au point d'arrivée final.

\begin{examplebox}
Un randonneur se déplace de $4$ km vers l'est et $1$ km vers le nord, puis de $2$ km vers l'ouest et $5$ km vers le nord. Les déplacements sont
\[
\mathbf u=(4,1),
\qquad
\mathbf v=(-2,5).
\]
Le déplacement total est
\[
\mathbf u+\mathbf v=(2,6).
\]
La distance en ligne droite jusqu'au point de départ est
\[
\sqrt{2^2+6^2}=2\sqrt{10}\ \mathrm{km}.
\]
\end{examplebox}

\subsection{Multiplication par un scalaire}

Le vecteur $k\mathbf u$ a une longueur multipliée par $\abs k$. Si $k>0$, le sens est conservé; si $k<0$, le sens est inversé. La relation
\[
\mathbf v=k\mathbf u
\]
caractérise la colinéarité de deux vecteurs non nuls.

Par exemple,
\[
(6,-3)=3(2,-1),
\]
donc les deux vecteurs sont parallèles. En revanche, $(6,-2)$ n'est pas un multiple de $(2,-1)$.

\subsection{Norme, distance et vecteur unitaire}

La norme vient du théorème de Pythagore :
\[
\norm{\mathbf v}=\sqrt{v_1^2+v_2^2}
\]
dans le plan, et
\[
\norm{\mathbf v}=\sqrt{v_1^2+v_2^2+v_3^2}
\]
dans l'espace.

La distance entre deux points est la norme du vecteur qui les relie :
\[
d(A,B)=\norm{\vect{AB}}.
\]
Si $\mathbf v\neq\mathbf0$, le vecteur
\[
\widehat{\mathbf v}=\frac{\mathbf v}{\norm{\mathbf v}}
\]
a la même direction et une norme égale à $1$.

\subsection{Produit scalaire en coordonnées}

Dans une base orthonormée,
\[
\mathbf u\cdot\mathbf v=u_1v_1+u_2v_2
\]
ou, dans $\R^3$,
\[
\mathbf u\cdot\mathbf v=u_1v_1+u_2v_2+u_3v_3.
\]
Cette formule provient de la distributivité et du fait que les vecteurs de base sont unitaires et perpendiculaires.

Si
\[
\mathbf u=u_1\mathbf e_1+u_2\mathbf e_2,
\qquad
\mathbf v=v_1\mathbf e_1+v_2\mathbf e_2,
\]
alors
\begin{align*}
\mathbf u\cdot\mathbf v
&=(u_1\mathbf e_1+u_2\mathbf e_2)\cdot(v_1\mathbf e_1+v_2\mathbf e_2)\\
&=u_1v_1(\mathbf e_1\cdot\mathbf e_1)
+u_1v_2(\mathbf e_1\cdot\mathbf e_2)
+u_2v_1(\mathbf e_2\cdot\mathbf e_1)
+u_2v_2(\mathbf e_2\cdot\mathbf e_2)\\
&=u_1v_1+u_2v_2.
\end{align*}

\subsection{Angle et orthogonalité}

La définition géométrique est
\[
\mathbf u\cdot\mathbf v
=
orm{\mathbf u}\norm{\mathbf v}\cos\theta.
\]
Pour deux vecteurs non nuls,
\[
\cos\theta=\frac{\mathbf u\cdot\mathbf v}{\norm{\mathbf u}\norm{\mathbf v}}.
\]
Si le produit scalaire est positif, l'angle est aigu; s'il est nul, l'angle est droit; s'il est négatif, l'angle est obtus.

\begin{examplebox}
Pour
\[
\mathbf u=(2,1),
\qquad
\mathbf v=(-1,2),
\]
on obtient
\[
\mathbf u\cdot\mathbf v=2(-1)+1(2)=0.
\]
Les vecteurs sont perpendiculaires. Ce test évite de calculer explicitement l'angle.
\end{examplebox}

\subsection{Projection d'un vecteur}

La projection de $\mathbf u$ sur une direction non nulle $\mathbf v$ est
\[
\operatorname{proj}_{\mathbf v}(\mathbf u)
=
\frac{\mathbf u\cdot\mathbf v}{\mathbf v\cdot\mathbf v}\mathbf v.
\]
Elle représente la composante de $\mathbf u$ parallèle à $\mathbf v$. Le reste
\[
\mathbf u-\operatorname{proj}_{\mathbf v}(\mathbf u)
\]
est perpendiculaire à $\mathbf v$.

\begin{examplebox}
Pour $\mathbf u=(4,3)$ et $\mathbf v=(1,1)$,
\[
\mathbf u\cdot\mathbf v=7,
\qquad
\mathbf v\cdot\mathbf v=2.
\]
Donc
\[
\operatorname{proj}_{\mathbf v}(\mathbf u)=\frac72(1,1)=\left(\frac72,\frac72\right).
\]
La composante perpendiculaire est
\[
\left(\frac12,-\frac12\right),
\]
dont le produit scalaire avec $(1,1)$ est nul.
\end{examplebox}

\subsection{Combinaisons linéaires et directions accessibles}

Une combinaison linéaire de $\mathbf u$ et $\mathbf v$ est un vecteur
\[
a\mathbf u+b\mathbf v.
\]
Dans le plan, deux vecteurs non colinéaires permettent de produire tout vecteur du plan. Deux vecteurs colinéaires ne produisent qu'une seule droite de directions.

Résoudre
\[
a\mathbf u+b\mathbf v=\mathbf w
\]
revient à résoudre un système linéaire sur les composantes. Cette idée prépare naturellement les chapitres sur les matrices et les systèmes.

\subsection{Vérifications essentielles}

\begin{itemize}
\item Une norme est toujours non négative.
\item Un vecteur unitaire doit avoir une norme $1$.
\item Le quotient utilisé dans la formule de l'angle doit appartenir à $[-1,1]$.
\item Une relation de colinéarité exige le même facteur pour toutes les composantes.
\item Une projection parallèle à $\mathbf v$ doit être un multiple de $\mathbf v$.
\end{itemize}

\begin{summarybox}
Les vecteurs traduisent les déplacements en calculs sur les composantes. La norme mesure la longueur, le produit scalaire relie coordonnées et angles, et les combinaisons linéaires décrivent les directions que l'on peut construire.
\end{summarybox}


\chapter{Équations de la droite et du plan}

\begin{questionbox}
Quelles données suffisent pour décrire une droite dans le plan ou l'espace, puis un plan dans $\R^3$?
\end{questionbox}

\section{Droite dans le plan}

Une droite peut être décrite de plusieurs façons. La forme affine
\[
y=mx+b
\]
convient aux droites non verticales. La forme générale
\[
ax+by=c
\]
inclut aussi les droites verticales.

\begin{definition}{Équation vectorielle d'une droite}{}
Une droite passant par un point $P_0$ et de vecteur directeur $\mathbf v\neq\mathbf0$ est
\[
\mathbf r(t)=\mathbf r_0+t\mathbf v,
\qquad t\in\R.
\]
\end{definition}

Dans $\R^2$, si $P_0=(x_0,y_0)$ et $\mathbf v=(a,b)$, cela donne les équations paramétriques
\[
x=x_0+at,
\qquad
y=y_0+bt.
\]

\begin{examplebox}
La droite passant par $P_0=(2,-1)$ et de direction $(3,4)$ possède l'équation
\[
(x,y)=(2,-1)+t(3,4),
\]
soit
\[
x=2+3t,
\qquad y=-1+4t.
\]
Comme $m=4/3$, sa forme affine est
\[
y+1=\frac43(x-2).
\]
\end{examplebox}

\section{Vecteur normal}

Pour la droite
\[
ax+by=c,
\]
le vecteur $\mathbf n=(a,b)$ est normal, donc perpendiculaire à la droite. Un vecteur directeur possible est $(b,-a)$.

\begin{examplebox}
La droite $2x-3y=6$ a pour vecteur normal $(2,-3)$ et pour vecteur directeur $(3,2)$. Sa pente est $2/3$.
\end{examplebox}

\section{Positions relatives de deux droites}

Pour deux droites dans $\R^2$ :
\begin{itemize}
\item directions non parallèles : une intersection;
\item directions parallèles mais points incompatibles : aucune intersection;
\item même direction et un point commun : droites confondues.
\end{itemize}
Dans $\R^3$, deux droites peuvent aussi être gauches : ni parallèles ni sécantes.

\section{Plan dans \texorpdfstring{$\R^3$}{R3}}

\begin{definition}{Équation cartésienne d'un plan}{}
Un plan de vecteur normal $\mathbf n=(a,b,c)\neq\mathbf0$ passant par $P_0=(x_0,y_0,z_0)$ est décrit par
\[
a(x-x_0)+b(y-y_0)+c(z-z_0)=0.
\]
Il s'écrit aussi
\[
ax+by+cz=d.
\]
\end{definition}

\begin{intuitionbox}
Le vecteur normal indique la direction perpendiculaire au plan. Tous les déplacements à l'intérieur du plan ont un produit scalaire nul avec ce vecteur.
\end{intuitionbox}

\begin{examplebox}
Le plan passant par $(1,2,-1)$ et de normale $(2,-1,3)$ vérifie
\[
2(x-1)-(y-2)+3(z+1)=0.
\]
Après simplification,
\[
2x-y+3z=-3.
\]
\end{examplebox}

Un plan peut aussi être décrit par un point et deux vecteurs directeurs non parallèles :
\[
\mathbf r(s,t)=\mathbf r_0+s\mathbf u+t\mathbf v.
\]

\section{Intersection d'une droite et d'un plan}

On substitue les équations paramétriques de la droite dans l'équation du plan.

\begin{examplebox}
Considérons la droite
\[
(x,y,z)=(1,0,2)+t(2,1,-1)
\]
et le plan
\[
x+y+z=4.
\]
La substitution donne
\[
(1+2t)+t+(2-t)=4,
\]
donc $3+2t=4$ et $t=1/2$. Le point d'intersection est
\[
\left(2,\frac12,\frac32\right).
\]
\end{examplebox}

Si le coefficient de $t$ disparaît, la droite est parallèle au plan; elle peut être entièrement contenue dans le plan ou ne jamais le rencontrer.


\section{Point, direction et vecteur normal}

\begin{visualbox}
Une droite naît lorsqu'on part d'un point et qu'on se déplace dans une seule direction. Un plan naît lorsqu'on part d'un point et qu'on combine deux directions non colinéaires. Les équations paramétriques décrivent ce mouvement; les équations cartésiennes décrivent une contrainte satisfaite par tous les points de l'objet.
\end{visualbox}

\begin{center}
\begin{tikzpicture}[scale=0.85]
  \fill[navy] (0,0) circle (2.5pt) node[below] {$P_0$};
  \draw[teal,very thick,-{Stealth[length=3mm]}] (0,0)--(2.2,1.2) node[midway,above] {$\mathbf v$};
  \draw[navy,very thick] (-2.2,-1.2)--(5,2.7);
  \node[navy] at (3.6,1.3) {$P_0+t\mathbf v$};
\end{tikzpicture}
\qquad
\begin{tikzpicture}[scale=0.75]
  \fill[mistteal] (-2,-1)--(2,-0.1)--(3,2)--(-1,1.1)--cycle;
  \draw[teal,thick] (-2,-1)--(2,-0.1)--(3,2)--(-1,1.1)--cycle;
  \fill[navy] (0,0) circle (2.5pt) node[below] {$P_0$};
  \draw[navy,very thick,-{Stealth[length=3mm]}] (0,0)--(1.7,0.4) node[below] {$\mathbf u$};
  \draw[gold,very thick,-{Stealth[length=3mm]}] (0,0)--(0.7,1.3) node[left] {$\mathbf v$};
  \draw[terracotta,very thick,-{Stealth[length=3mm]}] (0,0)--(-0.2,2.2) node[left] {$\mathbf n$};
\end{tikzpicture}
\end{center}

\begin{connectionbox}
Dans $ax+by=c$, le vecteur $(a,b)$ est normal à la droite. Dans $ax+by+cz=d$, le vecteur $(a,b,c)$ est normal au plan. Cette lecture explique immédiatement les tests de parallélisme et de perpendicularité : on compare des directions ou des normales, plutôt que de manipuler des coefficients sans figure.
\end{connectionbox}

\subsection{Choisir la représentation}
La forme paramétrique est naturelle pour générer des points et calculer une intersection. La forme cartésienne est naturelle pour vérifier si un point appartient à l'objet et pour lire un vecteur normal. Passer de l'une à l'autre est une traduction, non une nouvelle notion.

\section{Choisir et exploiter une représentation}

Une droite ou un plan peut être décrit de plusieurs façons. La forme paramétrique met en évidence le mouvement le long de directions; la forme cartésienne met en évidence une normale et facilite les tests d'appartenance. Passer d'une forme à l'autre permet de choisir la représentation la plus efficace.

\subsection{Droite par un point et une direction}

Une droite du plan passant par $P_0(x_0,y_0)$ et dirigée par $\mathbf v=(a,b)\neq\mathbf0$ s'écrit
\[
\begin{cases}
x=x_0+at,\\y=y_0+bt,
\end{cases}
\qquad t\in\R.
\]
Le paramètre $t$ indique combien de fois on ajoute le vecteur directeur.

\begin{examplebox}
La droite passant par $P_0=(1,-2)$ et dirigée par $(3,1)$ est
\[
\begin{cases}
x=1+3t,\\y=-2+t.
\end{cases}
\]
Pour $t=0$, on retrouve $P_0$; pour $t=2$, on obtient le point $(7,0)$.
\end{examplebox}

\subsection{Forme cartésienne et vecteur normal}

Une équation
\[
Ax+By=C
\]
a pour vecteur normal
\[
\mathbf n=(A,B).
\]
Un vecteur directeur peut être choisi comme
\[
\mathbf v=(-B,A),
\]
car
\[
(A,B)\cdot(-B,A)=0.
\]
Cette relation entre direction et normale est le pont entre la géométrie vectorielle et l'équation cartésienne.

Pour convertir la droite précédente, on prend un normal à $(3,1)$, par exemple $(1,-3)$. L'équation a donc la forme
\[
x-3y=C.
\]
En utilisant le point $(1,-2)$,
\[
C=1-3(-2)=7.
\]
Ainsi
\[
x-3y=7.
\]

\subsection{Pente et droites verticales}

Lorsque $B\neq0$, l'équation cartésienne devient
\[
y=-\frac ABx+\frac CB.
\]
La pente est $-A/B$. Mais une droite verticale a une équation $x=c$ et n'a pas de pente définie. La forme paramétrique et la forme cartésienne traitent donc plus uniformément les droites que la seule forme $y=mx+b$.

\subsection{Positions relatives de deux droites}

Deux droites sont parallèles si leurs vecteurs directeurs sont colinéaires. Elles sont confondues si, en plus, un point de l'une appartient à l'autre. Sinon, dans le plan, elles se coupent en un point unique.

Pour deux équations cartésiennes
\[
A_1x+B_1y=C_1,
\qquad
A_2x+B_2y=C_2,
\]
les droites sont parallèles lorsque les normales sont colinéaires. Résoudre le système permet de déterminer l'intersection éventuelle.

\subsection{Appartenance et position relative}

Pour savoir si un point $P(x_0,y_0)$ appartient à la droite $Ax+By=C$, on substitue ses coordonnées. Pour comparer deux droites, on examine d'abord leurs directions : des vecteurs directeurs colinéaires indiquent des droites parallèles ou confondues; sinon elles sont sécantes.

\begin{examplebox}
Le point $P=(2,1)$ appartient à la droite $3x-y=5$, car $3(2)-1=5$. La droite $6x-2y=7$ a la même normale $(6,-2)=2(3,-1)$, mais elle ne contient pas $P$; les deux droites sont donc parallèles et distinctes.
\end{examplebox}

\subsection{Plan dans l'espace}

Un plan passant par $P_0(x_0,y_0,z_0)$ et de vecteur normal
\[
\mathbf n=(A,B,C)\neq\mathbf0
\]
a pour équation
\[
A(x-x_0)+B(y-y_0)+C(z-z_0)=0.
\]
Sous forme développée,
\[
Ax+By+Cz=D.
\]
Tout vecteur reliant deux points du plan est perpendiculaire à $\mathbf n$.

\subsection{Plan à partir de trois points}

Trois points non alignés $A,B,C$ déterminent un plan. On forme deux directions
\[
\mathbf u=\vect{AB},
\qquad
\mathbf v=\vect{AC}.
\]
On cherche ensuite un vecteur $\mathbf n=(a,b,c)$ tel que
\[
\mathbf n\cdot\mathbf u=0,
\qquad
\mathbf n\cdot\mathbf v=0.
\]
Ce petit système fournit une normale, puis l'équation du plan est obtenue avec l'un des points.

\subsection{Intersection d'une droite et d'un plan}

Soit la droite
\[
\mathbf r(t)=\mathbf r_0+t\mathbf v
\]
et le plan
\[
Ax+By+Cz=D.
\]
On remplace $x,y,z$ par les expressions paramétriques. Une équation en $t$ apparaît.

\begin{examplebox}
Considérons
\[
\mathbf r(t)=(1,0,2)+t(2,-1,3)
\]
et le plan
\[
x+2y-z=4.
\]
La substitution donne
\[
(1+2t)+2(-t)-(2+3t)=4,
\]
\[
-1-3t=4,
\qquad t=-\frac53.
\]
Le point d'intersection est
\[
\left(1-\frac{10}{3},\frac53,2-5\right)
=
\left(-\frac73,\frac53,-3\right).
\]
\end{examplebox}

Si l'équation en $t$ devient une contradiction, la droite est parallèle au plan sans y être contenue. Si elle devient une identité, toute la droite appartient au plan.

\subsection{Comparer droites et plans dans l'espace}

Deux plans sont parallèles lorsque leurs normales sont colinéaires. S'ils ne sont pas parallèles, leur intersection est généralement une droite. Deux droites de l'espace peuvent être parallèles, sécantes ou gauches : des droites gauches ne sont ni parallèles ni situées dans un même plan.

\subsection{Lire une position relative avec les normales}

Deux plans sont parallèles si leurs vecteurs normaux sont colinéaires. Une droite de direction $\mathbf v$ est parallèle au plan de normale $\mathbf n$ lorsque $\mathbf v\cdot\mathbf n=0$. Si, en plus, un point de la droite satisfait l'équation du plan, toute la droite est contenue dans le plan.

\begin{connectionbox}
Cette lecture évite d'accumuler des formules supplémentaires : direction et normale suffisent pour décider d'un parallélisme, d'une perpendicularité ou d'une inclusion. La substitution sert ensuite à confirmer l'appartenance.
\end{connectionbox}

\subsection{Choix de représentation}

\begin{center}
\begin{tabularx}{\textwidth}{@{}lY@{}}
\toprule
Question & Représentation utile\\\midrule
produire des points d'une droite & paramétrique\\
tester l'appartenance d'un point & cartésienne\\
comparer des directions & vecteurs directeurs\\
tester une perpendicularité & vecteurs normaux et produit scalaire\\
chercher une intersection & substitution ou système\\
calculer une distance & forme cartésienne normalisée\\
\bottomrule
\end{tabularx}
\end{center}

\begin{summarybox}
Une droite est déterminée par un point et une direction; un plan par un point et une normale. Les formes paramétrique et cartésienne sont complémentaires. Les intersections se ramènent à des substitutions ou à des systèmes linéaires.
\end{summarybox}


\chapter{Matrices et transformations}

\begin{questionbox}
Pourquoi organiser des nombres en tableau, et comment une matrice peut-elle transformer un vecteur ou représenter un système d'équations?
\end{questionbox}

\section{Définition et dimensions}

\begin{definition}{Matrice}{}
Une matrice de taille $m\times n$ est un tableau rectangulaire de $m$ lignes et $n$ colonnes :
\[
A=(a_{ij})=
\begin{pmatrix}
a_{11}&\cdots&a_{1n}\\
\vdots&\ddots&\vdots\\
a_{m1}&\cdots&a_{mn}
\end{pmatrix}.
\]
\end{definition}

Deux matrices sont égales si elles ont la même taille et les mêmes entrées correspondantes.

\begin{examplebox}
La matrice
\[
A=\begin{pmatrix}2&-1&0\\3&4&5\end{pmatrix}
\]
est de taille $2\times3$. L'entrée $a_{23}$ vaut $5$.
\end{examplebox}

\section{Addition et multiplication par un scalaire}

Des matrices de même taille s'additionnent entrée par entrée. Pour un scalaire $\lambda$,
\[
(\lambda A)_{ij}=\lambda a_{ij}.
\]

\begin{examplebox}
\[
\begin{pmatrix}1&2\\-1&3\end{pmatrix}
+2\begin{pmatrix}0&4\\2&-1\end{pmatrix}
=
\begin{pmatrix}1&10\\3&1\end{pmatrix}.
\]
\end{examplebox}

\section{Produit matriciel}

Le produit $AB$ est défini lorsque le nombre de colonnes de $A$ égale le nombre de lignes de $B$. Si $A$ est $m\times n$ et $B$ est $n\times p$, alors $AB$ est $m\times p$ et
\[
(AB)_{ij}=\sum_{k=1}^n a_{ik}b_{kj}.
\]

\begin{intuitionbox}
Chaque entrée de $AB$ est le produit scalaire d'une ligne de $A$ avec une colonne de $B$. Les dimensions intérieures doivent donc correspondre.
\end{intuitionbox}

\begin{examplebox}
\[
\begin{pmatrix}1&2\\0&-1\end{pmatrix}
\begin{pmatrix}3&1\\4&2\end{pmatrix}
=
\begin{pmatrix}
1\cdot3+2\cdot4&1\cdot1+2\cdot2\\
0\cdot3-1\cdot4&0\cdot1-1\cdot2
\end{pmatrix}
=
\begin{pmatrix}11&5\\-4&-2\end{pmatrix}.
\]
\end{examplebox}

\begin{warningbox}
Le produit matriciel n'est généralement pas commutatif : $AB$ et $BA$ peuvent être différents, ou l'un des deux peut ne pas être défini.
\end{warningbox}

\section{Matrice identité et puissances}

La matrice identité $I_n$ possède des $1$ sur la diagonale et des $0$ ailleurs. Elle vérifie
\[
AI_n=A,
\qquad I_mA=A
\]
pour toute matrice $A$ de taille $m\times n$.

Une puissance $A^k$ représente la répétition de la même transformation matricielle.

\section{Matrices comme transformations du plan}

Une matrice $2\times2$
\[
A=\begin{pmatrix}a&b\\c&d\end{pmatrix}
\]
transforme un vecteur $(x,y)$ en
\[
A\begin{pmatrix}x\\y\end{pmatrix}
=
\begin{pmatrix}ax+by\\cx+dy\end{pmatrix}.
\]
Les colonnes de $A$ sont les images des vecteurs de base $\mathbf e_1=(1,0)$ et $\mathbf e_2=(0,1)$.

\begin{examplebox}
La matrice
\[
A=\begin{pmatrix}2&0\\0&\frac12\end{pmatrix}
\]
étire horizontalement d'un facteur $2$ et comprime verticalement d'un facteur $1/2$.
\end{examplebox}

\begin{center}
\begin{tikzpicture}[scale=0.75]
\begin{scope}[xshift=-3.8cm]
\draw[step=1cm,gray!35,very thin] (-1.2,-1.2) grid (2.2,2.2);
\draw[-{Stealth},thick] (-1.3,0)--(2.4,0);
\draw[-{Stealth},thick] (0,-1.3)--(0,2.4);
\draw[very thick,uqamblue] (0,0)--(1,0)--(1,1)--(0,1)--cycle;
\node at (0.5,-1.7) {carré unité};
\end{scope}
\begin{scope}[xshift=3.2cm]
\draw[xstep=2cm,ystep=0.5cm,gray!35,very thin] (-2.2,-0.7) grid (4.2,1.2);
\draw[-{Stealth},thick] (-2.3,0)--(4.4,0);
\draw[-{Stealth},thick] (0,-0.8)--(0,1.4);
\draw[very thick,teal] (0,0)--(2,0)--(2,0.5)--(0,0.5)--cycle;
\node at (1,-1.2) {image par $A$};
\end{scope}
\draw[-{Stealth[length=3mm]},thick,deepblue] (-0.7,0.5)--(0.7,0.5) node[midway,above] {$A$};
\end{tikzpicture}
\end{center}

\begin{connectionbox}[title={Notation utile : la transposée}]
La transposée $A^{\mathsf T}$ échange lignes et colonnes : $(A^{\mathsf T})_{ij}=a_{ji}$. Dans ce cours, elle sert surtout à écrire certains produits scalaires et à changer la présentation d'un tableau; elle n'est pas un nouveau type de transformation à mémoriser.
\end{connectionbox}

\section{Une matrice déplace une grille}

\begin{visualbox}
Une matrice carrée n'est pas d'abord un tableau à multiplier : c'est une règle qui indique où vont les vecteurs de base. Les colonnes de la matrice sont précisément les images de $\mathbf e_1=(1,0)$ et $\mathbf e_2=(0,1)$. Une fois ces deux images connues, toute la grille est déterminée par linéarité.
\end{visualbox}

\begin{center}
\begin{tikzpicture}[scale=0.72]
  \foreach \x in {-2,-1,0,1,2}{\draw[slate!35] (\x,-2)--(\x,2);}
  \foreach \y in {-2,-1,0,1,2}{\draw[slate!35] (-2,\y)--(2,\y);}
  \draw[-{Stealth},navy,very thick] (0,0)--(1,0) node[below] {$\mathbf e_1$};
  \draw[-{Stealth},teal,very thick] (0,0)--(0,1) node[left] {$\mathbf e_2$};
  \node at (0,-2.45) {grille initiale};
\end{tikzpicture}
\quad $\xrightarrow{\ A=\left(\begin{smallmatrix}2&1\\0&1\end{smallmatrix}\right)\ }$ \quad
\begin{tikzpicture}[scale=0.72]
  \foreach \a in {-2,-1,0,1,2}{\draw[slate!35] ({2*\a-2},-2)--({2*\a+2},2);}
  \foreach \b in {-2,-1,0,1,2}{\draw[slate!35] ({-2+\b},-2)--({2+\b},2);}
  \draw[-{Stealth},navy,very thick] (0,0)--(2,0) node[below] {$A\mathbf e_1$};
  \draw[-{Stealth},teal,very thick] (0,0)--(1,1) node[above left] {$A\mathbf e_2$};
  \node at (0,-2.45) {grille transformée};
\end{tikzpicture}
\end{center}

\begin{connectionbox}
Le produit $AB$ signifie : appliquer d'abord $B$, puis $A$. L'ordre compte parce que deux transformations successives ne produisent généralement pas le même résultat lorsqu'on les inverse. Cette non-commutativité est la version matricielle du fait que $f\circ g$ et $g\circ f$ sont souvent différentes.
\end{connectionbox}
\begin{center}
\begin{tikzpicture}[x=0.85cm,y=0.85cm,font=\small]
  \node[navy,anchor=east] at (-0.3,2.0) {$R$ puis $S$};
  \fill[mistblue] (0,1.5) rectangle (1.6,2.3);
  \draw[navy,thick] (0,1.5) rectangle (1.6,2.3);
  \draw[-{Stealth[length=2.4mm]},thick,slate] (1.9,1.9)--(3.1,1.9)
    node[midway,above] {$R$};
  \fill[mistteal] (3.4,1.1) rectangle (4.2,2.7);
  \draw[teal,thick] (3.4,1.1) rectangle (4.2,2.7);
  \draw[-{Stealth[length=2.4mm]},thick,slate] (4.5,1.9)--(5.7,1.9)
    node[midway,above] {$S$};
  \fill[mistgold] (6.0,1.1) rectangle (7.6,2.7);
  \draw[gold!80!black,thick] (6.0,1.1) rectangle (7.6,2.7);
  \node[gold!75!black] at (6.8,0.75) {$S(R(F))$};

  \node[teal,anchor=east] at (-0.3,-0.6) {$S$ puis $R$};
  \fill[mistblue] (0,-1.0) rectangle (1.6,-0.2);
  \draw[navy,thick] (0,-1.0) rectangle (1.6,-0.2);
  \draw[-{Stealth[length=2.4mm]},thick,slate] (1.9,-0.6)--(3.1,-0.6)
    node[midway,above] {$S$};
  \fill[mistteal] (3.4,-1.0) rectangle (6.6,-0.2);
  \draw[teal,thick] (3.4,-1.0) rectangle (6.6,-0.2);
  \draw[-{Stealth[length=2.4mm]},thick,slate] (6.9,-0.6)--(8.1,-0.6)
    node[midway,above] {$R$};
  \fill[mistred] (8.4,-2.2) rectangle (9.2,1.0);
  \draw[terracotta,thick] (8.4,-2.2) rectangle (9.2,1.0);
  \node[terracotta] at (8.8,-2.55) {$R(S(F))$};

  \node[ink] at (4.5,-3.15) {$S(R(F))\neq R(S(F))$};
\end{tikzpicture}
\end{center}


\begin{center}
\begin{tikzpicture}[node distance=1.3cm,every node/.style={font=\small,rounded corners,inner sep=6pt}]
 \node[draw=navy,fill=mistblue] (x) {$\mathbf x$};
 \node[draw=teal,fill=mistteal,right=of x] (bx) {$B\mathbf x$};
 \node[draw=gold,fill=mistgold,right=of bx] (abx) {$A(B\mathbf x)=AB\mathbf x$};
 \draw[->,very thick,teal] (x)--node[above] {$B$}(bx);
 \draw[->,very thick,gold!80!black] (bx)--node[above] {$A$}(abx);
\end{tikzpicture}
\end{center}

\subsection{L'inverse comme retour en arrière}
Une matrice est inversible seulement si elle ne confond pas deux directions distinctes. Si une transformation écrase tout le plan sur une droite, l'information perdue ne peut pas être reconstruite. Le déterminant nul signale exactement cette perte de dimension.


\section{Déterminant : aire, orientation et perte d'information}

Considérons la matrice
\[
A=\begin{pmatrix}a&b\\c&d\end{pmatrix}.
\]
Ses colonnes sont les vecteurs $\mathbf u=(a,c)$ et $\mathbf v=(b,d)$. Le carré unité est transformé en un parallélogramme engendré par $\mathbf u$ et $\mathbf v$.

\begin{center}
\begin{tikzpicture}[scale=0.9]
 \fill[mistblue] (0,0)--(2,0)--(2,2)--(0,2)--cycle;
 \draw[navy,thick] (0,0)--(2,0)--(2,2)--(0,2)--cycle;
 \draw[navy,very thick,-{Stealth[length=3mm]}] (0,0)--(2,0) node[below] {$\mathbf e_1$};
 \draw[teal,very thick,-{Stealth[length=3mm]}] (0,0)--(0,2) node[left] {$\mathbf e_2$};
 \node at (1,-0.55) {aire $1$};
\end{tikzpicture}
\quad $\xrightarrow{A}$ \quad
\begin{tikzpicture}[scale=0.9]
 \coordinate (O) at (0,0); \coordinate (U) at (3.2,0.8); \coordinate (V) at (1.1,2.5);
 \fill[mistteal] (O)--(U)--($(U)+(V)$)--(V)--cycle;
 \draw[teal,thick] (O)--(U)--($(U)+(V)$)--(V)--cycle;
 \draw[navy,very thick,-{Stealth[length=3mm]}] (O)--(U) node[below] {$\mathbf u$};
 \draw[gold,very thick,-{Stealth[length=3mm]}] (O)--(V) node[left] {$\mathbf v$};
 \node at (2.0,-0.5) {aire $|\det A|$};
\end{tikzpicture}
\end{center}

La valeur absolue $|ad-bc|$ est le facteur par lequel les aires sont multipliées. Le signe indique si l'orientation du plan est conservée ou renversée.

\subsection{Pourquoi le déterminant nul interdit l'inverse}
Si $\det A=0$, les colonnes sont colinéaires. Le parallélogramme s'aplatit en un segment d'aire nulle. Plusieurs points du plan sont envoyés sur la même image; il est impossible de savoir d'où provenait une sortie donnée. Aucune transformation inverse ne peut reconstruire l'information perdue.

\subsection{Composition et facteurs d'aire}
Appliquer $B$, puis $A$, multiplie d'abord les aires par $|\det B|$, puis par $|\det A|$. Il est donc naturel que
\[
\det(AB)=\det(A)\det(B).
\]
Cette propriété est une conséquence géométrique de la composition des transformations.

\begin{connectionbox}
Les trois notions suivantes racontent la même histoire :
\begin{itemize}
\item $\det A\ne0$ : la transformation ne perd pas de dimension;
\item $A^{-1}$ existe : la transformation peut être défaite;
\item $A\mathbf x=\mathbf b$ possède une solution unique pour chaque $\mathbf b$.
\end{itemize}
\end{connectionbox}


\section{Calculer, composer et interpréter une transformation}

Une matrice peut organiser des données, encoder un système ou représenter une transformation. Le produit matriciel paraît plus naturel lorsqu'on le lit comme une composition d'actions : les colonnes indiquent ce que devient chaque vecteur de base.

\subsection{Dimensions et compatibilité}

Une matrice $A$ de taille $m\times n$ possède $m$ lignes et $n$ colonnes. Elle transforme un vecteur de $\R^n$ en un vecteur de $\R^m$ :
\[
A:\R^n\longrightarrow\R^m.
\]
Cette lecture explique la condition de multiplication. Si $A$ est $m\times n$ et $B$ est $n\times p$, alors $AB$ est défini et a taille $m\times p$.

\subsection{Colonnes et images des vecteurs de base}

Pour
\[
A=
\begin{pmatrix}
a&b\\c&d
\end{pmatrix},
\]
on a
\[
A\begin{pmatrix}1\\0\end{pmatrix}
=
\begin{pmatrix}a\\c\end{pmatrix},
\qquad
A\begin{pmatrix}0\\1\end{pmatrix}
=
\begin{pmatrix}b\\d\end{pmatrix}.
\]
Les colonnes de $A$ sont donc les images des vecteurs de base. Pour un vecteur
\[
\mathbf x=x\mathbf e_1+y\mathbf e_2,
\]
la linéarité donne
\[
A\mathbf x=xA\mathbf e_1+yA\mathbf e_2.
\]
C'est exactement la règle du produit matrice-vecteur.

\subsection{Structure du produit matriciel}

Supposons que $B$ agit d'abord, puis $A$. La transformation composée est
\[
\mathbf x\longmapsto A(B\mathbf x).
\]
La matrice qui représente cette composition est $AB$. La $j$-ième colonne de $AB$ est l'image par $A$ de la $j$-ième colonne de $B$. Cette exigence détermine la règle ligne-par-colonne.

\begin{examplebox}
Soient
\[
A=\begin{pmatrix}1&2\\0&1\end{pmatrix},
\qquad
B=\begin{pmatrix}2&0\\0&3\end{pmatrix}.
\]
Alors
\[
AB=\begin{pmatrix}2&6\\0&3\end{pmatrix},
\qquad
BA=\begin{pmatrix}2&4\\0&3\end{pmatrix}.
\]
Les produits sont différents, car étirer puis cisailler n'est pas la même action que cisailler puis étirer.
\end{examplebox}

\subsection{Transformations usuelles du plan}

Quelques matrices importantes sont :
\[
\begin{pmatrix}k&0\\0&k\end{pmatrix}
\quad\text{(homothétie)},
\]
\[
\begin{pmatrix}1&0\\0&-1\end{pmatrix}
\quad\text{(réflexion selon l'axe horizontal)},
\]
\[
\begin{pmatrix}0&-1\\1&0\end{pmatrix}
\quad\text{(rotation d'un quart de tour)},
\]
\[
\begin{pmatrix}1&s\\0&1\end{pmatrix}
\quad\text{(cisaillement)}.
\]
Pour comprendre une matrice, il suffit souvent d'observer les images de $(1,0)$ et $(0,1)$, puis de reconstruire l'image du carré unité.

\subsection{Une transformation complète}

Considérons
\[
A=\begin{pmatrix}2&1\\1&1\end{pmatrix}.
\]
Les vecteurs de base deviennent
\[
\mathbf e_1\mapsto(2,1),
\qquad
\mathbf e_2\mapsto(1,1).
\]
Le carré unité devient le parallélogramme engendré par ces deux colonnes. Pour le point $(3,-1)$,
\[
A\begin{pmatrix}3\\-1\end{pmatrix}
=
\begin{pmatrix}5\\2\end{pmatrix}.
\]
Le calcul numérique et l'interprétation géométrique décrivent la même action.

\subsection{Matrices comme tableaux de données}

Supposons que les quantités de trois produits achetées par deux magasins sont organisées dans
\[
Q=\begin{pmatrix}12&8&5\\7&11&9\end{pmatrix},
\]
et que les prix unitaires sont
\[
\mathbf p=\begin{pmatrix}4\\6\\3\end{pmatrix}.
\]
Le produit
\[
Q\mathbf p
\]
donne le coût total de chaque magasin :
\[
Q\mathbf p=\begin{pmatrix}111\\121\end{pmatrix}.
\]
Le produit matriciel rassemble plusieurs produits scalaires en un seul calcul structuré.

\subsection{Inverse comme action réciproque}

Si une matrice $A$ possède une inverse, alors
\[
A^{-1}A=AA^{-1}=I.
\]
La transformation $A^{-1}$ défait l'action de $A$. Une transformation qui écrase tout le plan sur une droite perd de l'information et ne peut pas être inversée. Le chapitre suivant reliera cette perte d'information au déterminant.

\subsection{Contrôles pratiques}

\begin{itemize}
\item Vérifier les dimensions avant toute multiplication.
\item Se rappeler que $AB$ signifie que $B$ agit d'abord.
\item Tester une égalité matricielle sur les vecteurs de base.
\item Vérifier une inverse en calculant le produit avec la matrice initiale.
\item Interpréter les colonnes avant de multiplier de nombreux vecteurs.
\end{itemize}

\begin{summarybox}
Une matrice représente une action linéaire. Ses colonnes sont les images des vecteurs de base, le produit matriciel représente la composition et l'ordre des facteurs est essentiel. Les mêmes calculs organisent aussi des tableaux de données.
\end{summarybox}


\chapter{Systèmes d'équations linéaires}

\begin{questionbox}
Comment résoudre plusieurs équations simultanément, et comment savoir si le système possède une solution unique?
\end{questionbox}

\section{Trois points de vue sur un système}

\begin{visualbox}
Un système de deux équations linéaires décrit simultanément deux droites. L'algèbre ne crée pas les trois cas possibles : elle les révèle. Deux droites peuvent se couper une fois, ne jamais se couper, ou être la même droite.
\end{visualbox}

\begin{center}
\begin{tikzpicture}[scale=0.62]
 \draw[-{Stealth},slate] (-2.2,0)--(2.3,0); \draw[-{Stealth},slate] (0,-1.8)--(0,2.0);
 \draw[navy,thick] (-2,-1)--(2,1.4); \draw[teal,thick] (-2,1.2)--(2,-1.1);
 \fill[terracotta] (0.07,0.19) circle (2.3pt); \node at (0,-2.25) {une solution};
\end{tikzpicture}
\qquad
\begin{tikzpicture}[scale=0.62]
 \draw[-{Stealth},slate] (-2.2,0)--(2.3,0); \draw[-{Stealth},slate] (0,-1.8)--(0,2.0);
 \draw[navy,thick] (-2,-1)--(2,1.2); \draw[teal,thick] (-2,-0.1)--(2,2.1);
 \node at (0,-2.25) {aucune solution};
\end{tikzpicture}
\qquad
\begin{tikzpicture}[scale=0.62]
 \draw[-{Stealth},slate] (-2.2,0)--(2.3,0); \draw[-{Stealth},slate] (0,-1.8)--(0,2.0);
 \draw[navy,very thick] (-2,-1)--(2,1.2); \draw[teal,thick,dashed] (-2,-1)--(2,1.2);
 \node at (0,-2.25) {infinité de solutions};
\end{tikzpicture}
\end{center}

Considérons
\[
\begin{cases}
2x+y=5,\\
x-y=1.
\end{cases}
\]
Ce système peut être vu comme :
\begin{itemize}
\item l'intersection de deux droites;
\item deux contraintes algébriques à satisfaire simultanément;
\item l'équation matricielle
\[
\begin{pmatrix}2&1\\1&-1\end{pmatrix}
\begin{pmatrix}x\\y\end{pmatrix}
=
\begin{pmatrix}5\\1\end{pmatrix}.
\]
\end{itemize}

\section{Élimination de variables}

\begin{methodbox}
Les trois opérations élémentaires qui conservent les solutions sont :
\begin{enumerate}
\item échanger deux équations;
\item multiplier une équation par un nombre non nul;
\item ajouter à une équation un multiple d'une autre équation.
\end{enumerate}
\end{methodbox}

\begin{examplebox}
Pour
\[
\begin{cases}
2x+y=5,\\
x-y=1,
\end{cases}
\]
on additionne les équations : $3x=6$, donc $x=2$. Puis $2-y=1$, donc $y=1$.
\end{examplebox}

\section{Matrice augmentée et élimination de Gauss}

Le système
\[
A\mathbf x=\mathbf b
\]
se représente par la matrice augmentée $[A\mid\mathbf b]$. Les opérations sur les équations deviennent des opérations sur les lignes.

\begin{examplebox}
\[
\left[
\begin{array}{cc|c}
1&2&5\\
3&-1&4
\end{array}
\right]
\xrightarrow{L_2\leftarrow L_2-3L_1}
\left[
\begin{array}{cc|c}
1&2&5\\
0&-7&-11
\end{array}
\right].
\]
Ainsi, $y=11/7$, puis $x=13/7$.
\end{examplebox}

Une ligne de la forme
\[
[0\quad0\mid c],\qquad c\neq0,
\]
indique une contradiction et donc aucune solution. Une ligne entièrement nulle peut signaler une variable libre et une infinité de solutions.

\section{Déterminant en dimension deux}

\begin{definition}{Déterminant d'une matrice $2\times2$}{}
Pour
\[
A=\begin{pmatrix}a&b\\c&d\end{pmatrix},
\]
on définit
\[
\det A=ad-bc.
\]
\end{definition}

\begin{intuitionbox}
La valeur absolue de $\det A$ est le facteur par lequel la transformation $A$ multiplie les aires. Le signe indique si l'orientation est conservée ou inversée. Si $\det A=0$, le plan est aplati sur une droite ou un point.
\end{intuitionbox}

\begin{center}
\begin{tikzpicture}[scale=0.75]
\draw[-{Stealth},thick] (-1,0)--(5.2,0);
\draw[-{Stealth},thick] (0,-1)--(0,4.3);
\coordinate (O) at (0,0);
\coordinate (U) at (3,1);
\coordinate (V) at (1,3);
\coordinate (S) at (4,4);
\fill[teal!18] (O)--(U)--(S)--(V)--cycle;
\draw[very thick,uqamblue,-{Stealth[length=3mm]}] (O)--(U) node[midway,below] {$\mathbf u$};
\draw[very thick,gold,-{Stealth[length=3mm]}] (O)--(V) node[midway,left] {$\mathbf v$};
\draw[thick] (U)--(S)--(V);
\node at (2.2,2.15) {aire $=\abs{\det A}$};
\end{tikzpicture}
\end{center}

\section{Matrice inverse}

\begin{definition}{Matrice inverse}{}
Une matrice carrée $A$ est inversible s'il existe $A^{-1}$ telle que
\[
A^{-1}A=AA^{-1}=I.
\]
\end{definition}

Pour une matrice $2\times2$,
\[
A^{-1}=\frac1{ad-bc}
\begin{pmatrix}d&-b\\-c&a\end{pmatrix},
\]
à condition que $ad-bc\neq0$.

\begin{theorem}{Critère d'inversibilité}{}
Pour une matrice carrée $A$, les propriétés suivantes sont équivalentes :
\begin{itemize}
\item $A$ est inversible;
\item $\det A\neq0$;
\item le système $A\mathbf x=\mathbf b$ possède une solution unique pour tout $\mathbf b$.
\end{itemize}
\end{theorem}

\begin{examplebox}
Pour
\[
A=\begin{pmatrix}2&1\\1&-1\end{pmatrix},
\qquad \det A=-3,
\]
on a
\[
A^{-1}=\frac1{-3}\begin{pmatrix}-1&-1\\-1&2\end{pmatrix}
=\begin{pmatrix}1/3&1/3\\1/3&-2/3\end{pmatrix}.
\]
Le système $A\mathbf x=\binom51$ donne
\[
\mathbf x=A^{-1}\mathbf b=\binom21.
\]
\end{examplebox}

\section{Règle de Cramer}

\begin{theorem}{Règle de Cramer en dimension deux}{}
Pour
\[
\begin{cases}
ax+by=e,\\
cx+dy=f,
\end{cases}
\qquad D=ad-bc\neq0,
\]
la solution unique est
\[
x=\frac{\begin{vmatrix}e&b\\f&d\end{vmatrix}}{D},
\qquad
y=\frac{\begin{vmatrix}a&e\\c&f\end{vmatrix}}{D}.
\]
\end{theorem}

\begin{examplebox}
Pour
\[
\begin{cases}
2x+y=5,\\
x-y=1,
\end{cases}
\]
on a $D=-3$,
\[
D_x=\begin{vmatrix}5&1\\1&-1\end{vmatrix}=-6,
\qquad
D_y=\begin{vmatrix}2&5\\1&1\end{vmatrix}=-3.
\]
Donc $x=D_x/D=2$ et $y=D_y/D=1$.
\end{examplebox}

\begin{warningbox}
La règle de Cramer est élégante pour de petits systèmes, mais l'élimination de Gauss est généralement plus efficace et plus informative pour des systèmes de grande taille ou pour détecter les cas sans solution et avec infinité de solutions.
\end{warningbox}


\section{Interprétation géométrique d'un système}

\begin{visualbox}
Chaque équation linéaire décrit une droite dans le plan ou un plan dans l'espace. Résoudre le système revient à chercher leur intersection. Les trois résultats possibles sont donc géométriques avant d'être algébriques : un point unique, aucune intersection, ou une infinité de points communs.
\end{visualbox}

\begin{center}
\begin{tikzpicture}[scale=0.55]
 \draw[-{Stealth},thick,slate] (-2.8,0)--(2.9,0); \draw[-{Stealth},thick,slate] (0,-2.5)--(0,2.8);
 \draw[navy,very thick] (-2.4,-1.8)--(2.4,2.2); \draw[teal,very thick] (-2.4,2)--(2.4,-1.5);
 \fill[terracotta] (0.08,0.25) circle (3pt); \node[below] at (0,-2.9) {une solution};
\end{tikzpicture}
\qquad
\begin{tikzpicture}[scale=0.55]
 \draw[-{Stealth},thick,slate] (-2.8,0)--(2.9,0); \draw[-{Stealth},thick,slate] (0,-2.5)--(0,2.8);
 \draw[navy,very thick] (-2.4,-1.8)--(2.4,2.2); \draw[teal,very thick] (-2.4,-1.1)--(2.4,2.9);
 \node[below] at (0,-2.9) {aucune solution};
\end{tikzpicture}
\qquad
\begin{tikzpicture}[scale=0.55]
 \draw[-{Stealth},thick,slate] (-2.8,0)--(2.9,0); \draw[-{Stealth},thick,slate] (0,-2.5)--(0,2.8);
 \draw[navy,very thick] (-2.4,-1.8)--(2.4,2.2); \draw[teal,very thick,dashed] (-2.4,-1.8)--(2.4,2.2);
 \node[below] at (0,-2.9) {infinité de solutions};
\end{tikzpicture}
\end{center}

\begin{connectionbox}
Les opérations élémentaires sur les lignes remplacent les équations par d'autres équations ayant le même ensemble de solutions. Elles simplifient la description de l'intersection sans la déplacer. Un pivot correspond à une variable effectivement déterminée; une ligne contradictoire révèle l'absence de solution; une variable sans pivot reste libre.
\end{connectionbox}

\begin{center}
\begin{tikzpicture}[node distance=0.7cm,every node/.style={font=\small,rounded corners,inner sep=5pt}]
 \node[draw=navy,fill=mistblue] (s1) {$\begin{cases}x+y=5\\x-y=1\end{cases}$};
 \node[right=of s1] (arr) {$\Longrightarrow$};
 \node[draw=teal,fill=mistteal,right=of arr] (s2) {$\begin{cases}x+y=5\\2y=4\end{cases}$};
 \node[right=of s2] (arr2) {$\Longrightarrow$};
 \node[draw=gold,fill=mistgold,right=of arr2] (s3) {$y=2,\ x=3$};
\end{tikzpicture}
\end{center}

\subsection{Déterminant, inverse et unicité}
Pour une matrice $2\times2$, le déterminant mesure le facteur d'aire de la transformation associée. S'il est non nul, l'aire n'est pas écrasée et la transformation peut être inversée; le système $A\mathbf x=\mathbf b$ possède alors une solution unique pour tout $\mathbf b$. La règle de Cramer et la formule de l'inverse expriment cette même structure sous deux formes calculatoires.


\section{Pourquoi l'élimination conserve les solutions}

Prenons le système
\[
\begin{cases}
2x+y=7,\\
x-y=2.
\end{cases}
\]
Le point solution appartient aux deux droites. Si l'on additionne les deux équations, tout point commun satisfait aussi $3x=9$. Cette nouvelle équation ne remplace pas une droite par une droite arbitraire : elle extrait une conséquence que tout point commun doit satisfaire.

\begin{center}
\begin{tikzpicture}[scale=0.75]
 \draw[-{Stealth},thick,slate] (-0.5,0)--(5.5,0) node[right] {$x$};
 \draw[-{Stealth},thick,slate] (0,-1.5)--(0,5.5) node[above] {$y$};
 \draw[navy,very thick,domain=0.8:4.2] plot(\x,{7-2*\x});
 \draw[teal,very thick,domain=0.4:4.5] plot(\x,{\x-2});
 \draw[terracotta,very thick,dashed] (3,-1.2)--(3,5.1);
 \fill[gold] (3,1) circle (3pt);
 \node[gold!75!black,right] at (3.1,1.3) {point commun};
 \node[terracotta,rotate=90] at (3.25,4) {$3x=9$};
\end{tikzpicture}
\end{center}

\subsection{Les lignes d'une matrice sont des équations}
La matrice augmentée
\[
\left(\begin{array}{cc|c}2&1&7\\1&-1&2\end{array}\right)
\]
conserve seulement les coefficients. L'opération $L_1\leftarrow L_1+L_2$ correspond exactement à l'addition des équations. La notation matricielle ne crée donc pas une nouvelle méthode; elle rend la méthode d'élimination plus compacte et systématique.

\subsection{Pivots et degrés de liberté}
Un pivot fixe une variable après substitution. Si chaque variable possède un pivot, la solution est unique. Une ligne de la forme
\[
0=5
\]
est contradictoire et indique qu'aucun point ne satisfait toutes les équations. Une ligne $0=0$ n'apporte aucune contrainte nouvelle; une variable peut alors rester libre, ce qui produit une infinité de solutions.

\begin{visualbox}
La vérification la plus sûre reste la substitution dans le système initial. Elle contrôle à la fois les calculs d'élimination et les éventuelles erreurs de transcription dans la matrice augmentée.
\end{visualbox}



\subsection{Choisir entre élimination, inverse et règle de Cramer}
Ces méthodes résolvent le même problème, mais ne servent pas le même objectif.
\begin{itemize}
\item L'\textbf{élimination} est la méthode générale : elle révèle les pivots et traite aussi les systèmes sans solution ou avec une infinité de solutions.
\item La \textbf{matrice inverse} met en évidence la structure $\mathbf x=A^{-1}\mathbf b$, mais elle exige que $A$ soit carrée et inversible.
\item La \textbf{règle de Cramer} fournit des formules explicites avec des déterminants; elle est surtout utile pour de petits systèmes ou pour un raisonnement théorique.
\end{itemize}

\begin{center}
\begin{tikzpicture}[node distance=0.75cm and 0.8cm,every node/.style={font=\small,rounded corners,inner sep=6pt,align=center}]
 \node[draw=navy,fill=mistblue] (q) {Que veut-on comprendre?};
 \node[draw=teal,fill=mistteal,below left=of q] (g) {tous les cas\\et les pivots};
 \node[draw=gold,fill=mistgold,below=of q] (i) {structure\\$A^{-1}\mathbf b$};
 \node[draw=terracotta,fill=mistred,below right=of q] (c) {formules explicites\\petit système};
 \node[below=0.45cm of g] {élimination};
 \node[below=0.45cm of i] {inverse};
 \node[below=0.45cm of c] {Cramer};
 \draw[->] (q)--(g); \draw[->] (q)--(i); \draw[->] (q)--(c);
\end{tikzpicture}
\end{center}

\subsection{Une vérification matricielle compacte}
Après avoir obtenu $\mathbf x$, le calcul de $A\mathbf x$ doit redonner $\mathbf b$. Cette vérification teste toutes les équations simultanément et rappelle que la solution n'est pas une suite de nombres isolés, mais le vecteur envoyé sur $\mathbf b$ par la transformation $A$.


\section{Résoudre et interpréter un système}

Un système linéaire peut être lu comme un ensemble d'équations, comme une intersection d'objets géométriques ou comme une équation matricielle. Ces trois lectures expliquent les trois possibilités : une solution unique, aucune solution ou une infinité de solutions.

\subsection{Trois formes du même problème}

Le système
\[
\begin{cases}
2x+y=7,\\
x-y=2
\end{cases}
\]
représente l'intersection de deux droites et l'équation matricielle
\[
\begin{pmatrix}2&1\\1&-1\end{pmatrix}
\begin{pmatrix}x\\y\end{pmatrix}
=
\begin{pmatrix}7\\2\end{pmatrix}.
\]
L'élimination donne $3x=9$, donc $x=3$ et $y=1$. Géométriquement, les droites se coupent en $(3,1)$.

\subsection{Invariance de l'ensemble des solutions sous les opérations élémentaires}

Les opérations suivantes ne changent pas l'ensemble des solutions :
\begin{itemize}
\item échanger deux équations;
\item multiplier une équation par un nombre non nul;
\item ajouter à une équation un multiple d'une autre.
\end{itemize}
Elles remplacent le système par une description équivalente des mêmes contraintes. Dans une matrice augmentée, elles deviennent les opérations élémentaires sur les lignes.

\subsection{Élimination de Gauss : calcul d'une forme échelonnée}

Résolvons
\[
\begin{cases}
x+2y-z=2,\\
2x-y+3z=9,\\
3x+y+2z=11.
\end{cases}
\]
La matrice augmentée est
\[
\left(
\begin{array}{ccc|c}
1&2&-1&2\\
2&-1&3&9\\
3&1&2&11
\end{array}
\right).
\]
On effectue
\[
L_2\leftarrow L_2-2L_1,
\qquad
L_3\leftarrow L_3-3L_1,
\]
ce qui donne
\[
\left(
\begin{array}{ccc|c}
1&2&-1&2\\
0&-5&5&5\\
0&-5&5&5
\end{array}
\right).
\]
Puis
\[
L_3\leftarrow L_3-L_2
\]
donne une ligne nulle. Le système possède une variable libre : les deux dernières équations sont identiques. On obtient
\[
-y+z=1,
\qquad
x+2y-z=2.
\]
En posant $z=t$,
\[
y=t-1,
\qquad
x=4-t.
\]
Il existe une infinité de solutions
\[
(x,y,z)=(4-t,t-1,t),
\qquad t\in\R.
\]

\subsection{Reconnaître les trois cas dans une forme échelonnée}

\begin{itemize}
\item Une ligne $[0\ 0\ \cdots\ 0\mid c]$ avec $c\neq0$ produit une contradiction : aucune solution.
\item Un pivot dans chaque colonne de variable donne une solution unique.
\item Une colonne sans pivot crée une variable libre : généralement une infinité de solutions.
\end{itemize}
La forme échelonnée révèle donc la structure du système, pas seulement les valeurs numériques.

\subsection{Déterminant et changement d'aire}

Pour
\[
A=\begin{pmatrix}a&b\\c&d\end{pmatrix},
\]
le déterminant est
\[
\det A=ad-bc.
\]
Sa valeur absolue est le facteur par lequel la transformation multiplie les aires. Son signe indique si l'orientation est conservée ou renversée.

Si $\det A=0$, le carré unité est écrasé sur une figure d'aire nulle, par exemple une droite. Les colonnes sont alors colinéaires et la transformation perd une direction; aucune inverse n'existe.

\subsection{Dérivation de l'inverse d'une matrice \texorpdfstring{$2\times2$}{2 par 2}}

Cherchons
\[
A^{-1}=\begin{pmatrix}p&q\\r&s\end{pmatrix}
\]
tel que $AA^{-1}=I$. Le produit donne un système sur $p,q,r,s$. La solution se rassemble sous la forme
\[
A^{-1}=\frac1{ad-bc}
\begin{pmatrix}d&-b\\-c&a\end{pmatrix},
\]
à condition que $ad-bc\neq0$.

On échange les éléments diagonaux, on change le signe des éléments hors diagonale, puis on divise par le déterminant. La division par $\det A$ explique pourquoi l'inverse n'existe pas lorsque le déterminant est nul.

\begin{examplebox}
Pour
\[
A=\begin{pmatrix}2&1\\5&3\end{pmatrix},
\]
on a $\det A=6-5=1$, donc
\[
A^{-1}=\begin{pmatrix}3&-1\\-5&2\end{pmatrix}.
\]
Le produit $AA^{-1}$ donne bien $I$.
\end{examplebox}

\subsection{Résolution par l'inverse}

Si
\[
A\mathbf x=\mathbf b
\]
et $A^{-1}$ existe, alors
\[
\mathbf x=A^{-1}\mathbf b.
\]
Cette formule est conceptuellement simple, mais l'élimination de Gauss est souvent plus efficace pour des systèmes plus grands. L'inverse est surtout utile lorsqu'une même matrice doit être appliquée à plusieurs vecteurs $\mathbf b$.

\subsection{Règle de Cramer}

Pour
\[
\begin{cases}
ax+by=e,\\
cx+dy=f,
\end{cases}
\]
avec
\[
\Delta=ad-bc\neq0,
\]
la règle de Cramer donne
\[
x=\frac{ed-bf}{\Delta},
\qquad
 y=\frac{af-ec}{\Delta}.
\]
Elle remplace successivement une colonne de la matrice des coefficients par le vecteur des constantes.

\begin{examplebox}
Pour
\[
\begin{cases}
2x+3y=7,\\
-x+4y=5,
\end{cases}
\]
\[
\Delta=\begin{vmatrix}2&3\\-1&4\end{vmatrix}=11.
\]
Alors
\[
\Delta_x=\begin{vmatrix}7&3\\5&4\end{vmatrix}=13,
\qquad
\Delta_y=\begin{vmatrix}2&7\\-1&5\end{vmatrix}=17,
\]
\[
x=\frac{13}{11},
\qquad y=\frac{17}{11}.
\]
La substitution vérifie les deux équations.
\end{examplebox}

\subsection{Modélisation par un système}

Une cafétéria vend deux menus. Le menu A utilise $2$ unités de pain et $1$ unité de fromage; le menu B utilise $1$ unité de pain et $3$ unités de fromage. Une journée utilise $70$ unités de pain et $90$ unités de fromage. Si $x$ et $y$ sont les nombres de menus,
\[
\begin{cases}
2x+y=70,\\
x+3y=90.
\end{cases}
\]
L'élimination donne $x=24$ et $y=22$. Les unités et la non-négativité donnent un contrôle contextuel immédiat.

\subsection{Choisir la méthode}

\begin{center}
\begin{tabularx}{\textwidth}{@{}lY@{}}
\toprule
Situation & Méthode utile\\\midrule
deux petites équations & substitution ou élimination\\
système plus grand & matrice augmentée et Gauss\\
même matrice, plusieurs seconds membres & inverse, si elle existe\\
petit système carré et formule explicite souhaitée & Cramer\\
classification des solutions & forme échelonnée et pivots\\
\bottomrule
\end{tabularx}
\end{center}

\begin{summarybox}
L'élimination transforme un système sans changer ses solutions. Les pivots révèlent l'unicité ou les variables libres. Le déterminant mesure la perte de dimension, contrôle l'inversibilité et permet la règle de Cramer.
\end{summarybox}


\chapter{Programmation linéaire à deux variables}

\begin{questionbox}
Comment maximiser ou minimiser une quantité lorsque plusieurs contraintes linéaires doivent être satisfaites simultanément?
\end{questionbox}

\section{Variables, contraintes et fonction objectif}

Un problème de programmation linéaire comprend :
\begin{itemize}
\item des variables de décision, par exemple $x$ et $y$;
\item des contraintes linéaires;
\item une fonction objectif linéaire à maximiser ou minimiser.
\end{itemize}

\begin{examplebox}
Une entreprise fabrique deux produits. Le produit A utilise $2$ heures de machine et $1$ heure de travail; le produit B utilise $1$ heure de machine et $2$ heures de travail. On dispose de $8$ heures de machine et de $8$ heures de travail. Si les profits unitaires sont $30$ et $40$ dollars, on écrit
\[
\begin{cases}
2x+y\le8,\\
x+2y\le8,\\
x\ge0,\ y\ge0,
\end{cases}
\qquad
P=30x+40y.
\]
\end{examplebox}

\section{Région réalisable}

Chaque inéquation linéaire détermine un demi-plan. L'intersection de tous les demi-plans est la région réalisable.

\begin{center}
\begin{tikzpicture}[scale=0.8]
\begin{axis}[
axis lines=middle,xmin=-0.5,xmax=5.2,ymin=-0.5,ymax=5.2,
xtick={0,1,2,3,4,5},ytick={0,1,2,3,4,5},
xlabel={$x$},ylabel={$y$},width=9cm,height=8cm,grid=both,grid style={gray!20}]
\addplot[domain=0:4,uqamblue,thick] {8-2*x};
\addplot[domain=0:5,teal,thick] {(8-x)/2};
\addplot[fill=softgreen,fill opacity=0.8] coordinates {(0,0) (4,0) (2.6667,2.6667) (0,4)} -- cycle;
\addplot[only marks,mark=*,mark size=2pt] coordinates {(0,0) (4,0) (2.6667,2.6667) (0,4)};
\node at (axis cs:1.25,1.4) {région réalisable};
\end{axis}
\end{tikzpicture}
\end{center}

Une région peut être bornée, non bornée ou vide. Une contrainte est redondante si elle ne modifie pas la région.

\section{Principe d'optimalité aux sommets}

Les lignes de niveau de l'objectif
\[
ax+by=k
\]
sont parallèles. En déplaçant cette ligne dans la direction où $k$ augmente, le dernier contact avec une région polygonale bornée se produit à un sommet ou sur tout un côté.

\begin{theorem}{Principe des sommets}{}
Sur une région polygonale bornée, non vide et possédant au moins un sommet, toute fonction objectif linéaire atteint un maximum et un minimum. Au moins une solution optimale est un sommet; si un côté entier est optimal, ses deux extrémités le sont aussi.
\end{theorem}

\begin{methodbox}
Pour résoudre graphiquement un problème à deux variables :
\begin{enumerate}
\item définir les variables et leurs unités;
\item traduire chaque contrainte en inéquation;
\item tracer la région réalisable;
\item déterminer les coordonnées de ses sommets;
\item évaluer la fonction objectif à chaque sommet;
\item interpréter la meilleure valeur dans le contexte.
\end{enumerate}
\end{methodbox}

\begin{examplebox}
Pour la région précédente, les sommets sont
\[
(0,0),\quad(4,0),\quad\left(\frac83,\frac83\right),\quad(0,4).
\]
On évalue $P=30x+40y$ :
\[
0,\quad120,\quad\frac{560}{3}\approx186{,}67,\quad160.
\]
Le maximum continu est atteint au point $(8/3,8/3)$. Si les produits doivent être fabriqués en nombres entiers, il faut comparer les points entiers réalisables voisins; le problème devient alors un problème de programmation entière.
\end{examplebox}

\section{Configurations particulières du domaine réalisable}

\begin{itemize}
\item \textbf{Plusieurs solutions optimales :} une ligne de niveau est parallèle à un côté optimal.
\item \textbf{Région vide :} les contraintes sont incompatibles.
\item \textbf{Objectif non borné :} la région permet d'augmenter indéfiniment la fonction objectif.
\end{itemize}

\begin{warningbox}
Un optimum mathématique doit respecter le sens des variables. Des quantités physiques peuvent exiger $x,y\ge0$, des valeurs entières ou des capacités minimales. Ces conditions doivent être écrites comme contraintes, et non ajoutées après le calcul.
\end{warningbox}


\section{Optimiser en faisant glisser une droite}

\begin{visualbox}
Les contraintes dessinent une région de choix possibles. La fonction objectif $P=ax+by$ ne sélectionne pas d'abord un point : chaque valeur $P=c$ dessine une droite. Maximiser $P$ revient à faire glisser cette droite parallèlement à elle-même dans la direction où $c$ augmente, jusqu'au dernier contact avec la région réalisable.
\end{visualbox}

\begin{center}
\begin{tikzpicture}[scale=0.88]
  \draw[-{Stealth},thick,slate] (-0.4,0)--(6.2,0) node[right] {$x$};
  \draw[-{Stealth},thick,slate] (0,-0.4)--(0,5.5) node[above] {$y$};
  \fill[mistteal] (0,0)--(4,0)--(3.2,2.4)--(1.2,4.2)--(0,3.6)--cycle;
  \draw[teal,very thick] (0,0)--(4,0)--(3.2,2.4)--(1.2,4.2)--(0,3.6)--cycle;
  \foreach \c/\sty in {4/slate!55,7/gold!65,10/terracotta}{
    \draw[\sty,thick,dashed] (-0.2,{\c/2})--({\c/3},0);
  }
  \fill[terracotta] (1.2,4.2) circle (3pt);
  \draw[terracotta,very thick,-{Stealth[length=3mm]}] (4.6,1.2)--(5.5,2.3);
  \node[terracotta,align=left] at (5.0,3.1) {valeurs de\\l'objectif croissantes};
  \node[teal] at (2.0,1.8) {région réalisable};
\end{tikzpicture}
\end{center}

\begin{connectionbox}
Pourquoi un sommet suffit-il? Sur un segment, une fonction linéaire varie sans courbure : elle augmente, diminue ou reste constante. Son maximum sur ce segment est donc à une extrémité, sauf si tout le segment est optimal. En répétant ce raisonnement sur les côtés d'un polygone, on obtient le principe d'optimalité aux sommets.
\end{connectionbox}

\subsection{Le modèle avant le graphique}
Les erreurs les plus graves viennent rarement du calcul des sommets. Elles viennent du choix des variables, des unités ou du sens d'une inégalité. Avant de tracer, traduire chaque phrase en une contrainte et tester un point évident, souvent $(0,0)$, pour vérifier le bon demi-plan.


\section{De la phrase à la région réalisable}

Considérons un atelier qui fabrique des tables et des étagères. Une table exige $3$ heures de découpe et $2$ heures d'assemblage; une étagère exige $1$ heure de découpe et $2$ heures d'assemblage. L'atelier dispose de $18$ heures de découpe et de $16$ heures d'assemblage.

On pose
\[
x=\text{nombre de tables},
\qquad
y=\text{nombre d'étagères}.
\]
Les contraintes deviennent
\[
3x+y\le18,
\qquad
2x+2y\le16,
\qquad
x\ge0,
\quad y\ge0.
\]
Chaque coefficient porte une unité : $3$ heures de découpe par table, par exemple. Les unités expliquent la forme de la contrainte et évitent d'intervertir les coefficients.

\subsection{Tester le bon demi-plan}
Pour $3x+y\le18$, le point $(0,0)$ satisfait l'inégalité; la région admissible se trouve donc du côté de la droite contenant l'origine. Ce test évite de mémoriser des règles de hachurage.

\begin{center}
\begin{tikzpicture}[scale=0.72]
 \begin{axis}[axis lines=middle,xmin=-0.5,xmax=7,ymin=-0.5,ymax=9,
 width=10cm,height=8cm,xlabel={$x$},ylabel={$y$},grid=both]
  \addplot[draw=none,fill=mistteal] coordinates {(0,0)(6,0)(4,4)(0,8)} -- cycle;
  \addplot[navy,very thick,domain=0:6] {18-3*x};
  \addplot[teal,very thick,domain=0:8] {8-x};
  \addplot[terracotta,thick,dashed,domain=0:7] {10-1.5*x};
  \addplot[only marks,mark=*,mark size=2pt,gold] coordinates {(0,0)(6,0)(4,4)(0,8)};
 \end{axis}
\end{tikzpicture}
\end{center}

\subsection{L'objectif comme famille de droites}
Si le profit est $P=70x+40y$, chaque valeur de $P$ donne une droite parallèle. Déplacer cette droite dans le sens des profits croissants montre visuellement quel sommet est atteint en dernier. Le calcul aux sommets confirme ensuite cette lecture.

\begin{connectionbox}
La programmation linéaire sépare trois questions :
\begin{enumerate}
\item qu'est-il \textbf{possible} de faire? - les contraintes;
\item qu'est-il \textbf{souhaitable} de faire? - la fonction objectif;
\item quelle solution respecte les \textbf{unités et le contexte}? - l'interprétation finale.
\end{enumerate}
\end{connectionbox}


\section{De la décision au modèle géométrique}

La programmation linéaire combine trois éléments : des variables de décision, des contraintes linéaires et une fonction objectif. La difficulté principale n'est pas le calcul aux sommets, mais la traduction fidèle du contexte en inégalités.

\subsection{Nommer les variables et leurs unités}

Avant toute formule, il faut écrire une phrase complète. Par exemple :
\[
x=\text{nombre de lots du produit A},
\qquad
 y=\text{nombre de lots du produit B}.
\]
Les contraintes de non-négativité
\[
x\geq0,
\qquad y\geq0
\]
proviennent du sens des variables, non d'une convention graphique.

\subsection{Construire les contraintes}

Supposons que A utilise $2$ h de machine et $1$ kg de matière, tandis que B utilise $1$ h de machine et $3$ kg de matière. On dispose de $40$ h et $45$ kg. Les contraintes sont
\[
2x+y\leq40,
\]
\[
x+3y\leq45.
\]
Chaque inégalité décrit un demi-plan. Leur intersection avec le premier quadrant est la région réalisable.

\subsection{Choisir le bon côté d'une droite}

Pour l'inégalité
\[
2x+y\leq40,
\]
on trace d'abord la frontière $2x+y=40$. Le point $(0,0)$ satisfait l'inégalité, donc le demi-plan contenant l'origine est retenu. Un point test qui n'appartient pas à la frontière suffit à choisir le côté.

\subsection{Fonction objectif et lignes de niveau}

Si les profits unitaires sont $30\,$\$ pour A et $40\,$\$ pour B,
\[
P=30x+40y.
\]
Les droites
\[
30x+40y=k
\]
sont parallèles. Augmenter $k$ déplace la droite dans la direction d'amélioration. Le dernier contact avec une région polygonale bornée se produit à un sommet ou le long d'un côté entier.

\subsection{Étude complète d'un problème}

Considérons
\[
\begin{cases}
2x+y\leq40,\\
x+3y\leq45,\\
x\geq0,\ y\geq0.
\end{cases}
\]
Les sommets sur les axes sont $(20,0)$ et $(0,15)$. L'intersection des deux frontières satisfait
\[
\begin{cases}
2x+y=40,\\
x+3y=45.
\end{cases}
\]
En éliminant $y$,
\[
6x+3y=120,
\]
\[
5x=75,
\qquad x=15,
\]

donc $y=10$. Les sommets sont
\[
(0,0),\ (20,0),\ (15,10),\ (0,15).
\]
Évaluons
\[
P=30x+40y:
\]
\[
P(0,0)=0,
\quad
P(20,0)=600,
\quad
P(15,10)=850,
\quad
P(0,15)=600.
\]
L'optimum est atteint en $(15,10)$ avec un profit de $850\,$\$.

La solution utilise entièrement les deux ressources, car les deux contraintes sont actives au sommet optimal.

\subsection{Principe d'optimalité aux sommets}

Sur un segment, une fonction linéaire varie sans courbure. Sa valeur maximale est donc atteinte à une extrémité, à moins qu'elle soit constante sur tout le segment. Une région polygonale est composée de segments; en déplaçant une ligne de niveau, on atteint donc un sommet ou un côté parallèle à la ligne de niveau.

Cette explication géométrique est plus importante que la règle « tester les sommets », car elle montre quand un optimum multiple peut apparaître.

\subsection{Contraintes redondantes}

Une contrainte est redondante si sa suppression ne change pas la région réalisable. Par exemple, si
\[
x+y\leq10,
\qquad x\geq0,
\qquad y\geq0,
\]
la contrainte $x+y\leq15$ n'ajoute aucune restriction. La repérer simplifie le graphique mais ne modifie pas le modèle.

\subsection{Région vide, non bornée et optimum multiple}

Une région vide signifie que les contraintes sont incompatibles. Une région non bornée ne signifie pas automatiquement que l'objectif est sans maximum : tout dépend de la direction de croissance de la fonction objectif. Un optimum multiple apparaît lorsque la ligne de niveau optimale est parallèle à un côté de la région.

\subsection{Procédure complète de modélisation}

\begin{enumerate}
\item nommer les variables et les unités;
\item traduire chaque ressource ou exigence en inégalité;
\item ajouter les contraintes naturelles, notamment la non-négativité;
\item écrire la fonction objectif;
\item tracer chaque frontière puis choisir le bon demi-plan;
\item calculer les sommets;
\item évaluer l'objectif;
\item interpréter la solution et vérifier l'intégralité si nécessaire.
\end{enumerate}

\begin{summarybox}
La programmation linéaire transforme une décision en région géométrique. Les contraintes définissent les choix admissibles; la fonction objectif déplace des lignes parallèles; les sommets concentrent les candidats à l'optimum.
\end{summarybox}


\part{Trigonométrie et phénomènes périodiques}
% ============================================================

\chapter{Angles et cercle trigonométrique}

\begin{questionbox}
Comment un angle devient-il un nombre, et comment le cercle unité permet-il de définir sinus et cosinus pour tous les angles?
\end{questionbox}

\section{Degrés et radians}

Un degré partage un tour complet en $360$ parties. Le radian relie directement un angle à une longueur d'arc.

\begin{definition}{Mesure en radians}{}
Pour un cercle de rayon $r$, un angle au centre qui intercepte un arc de longueur $s$ possède la mesure
\[
\theta=\frac sr
\]
en radians.
\end{definition}

Un tour complet a une longueur $2\pi r$, donc mesure $2\pi$ radians. Ainsi,
\[
180^\circ=\pi\text{ radians}.
\]
Les conversions sont
\[
\theta_{\mathrm{rad}}=\theta_{\mathrm{deg}}\frac{\pi}{180},
\qquad
\theta_{\mathrm{deg}}=\theta_{\mathrm{rad}}\frac{180}{\pi}.
\]

\begin{examplebox}
\[
60^\circ=60\frac{\pi}{180}=\frac\pi3,
\qquad
\frac{5\pi}{6}=\frac{5\pi}{6}\frac{180}{\pi}=150^\circ.
\]
\end{examplebox}

\section{Le cercle unité}

\begin{definition}{Cercle trigonométrique}{}
Le cercle trigonométrique est le cercle de centre $(0,0)$ et de rayon $1$. Le point associé à l'angle $\theta$, mesuré depuis l'axe positif des $x$ dans le sens antihoraire, est
\[
P(\cos\theta,\sin\theta).
\]
\end{definition}

\begin{center}
\begin{tikzpicture}[scale=2.3]
\draw[-{Stealth},thick] (-1.25,0)--(1.35,0) node[right] {$x$};
\draw[-{Stealth},thick] (0,-1.25)--(0,1.35) node[above] {$y$};
\draw[thick,uqamblue] (0,0) circle (1);
\def\ang{42}
\coordinate (P) at ({cos(\ang)},{sin(\ang)});
\draw[very thick,teal,-{Stealth[length=2mm]}] (0,0)--(P);
\draw[dashed] (P)--({cos(\ang)},0);
\draw[dashed] (P)--(0,{sin(\ang)});
\draw[gold,thick] (0.30,0) arc[start angle=0,end angle=42,radius=0.30];
\node[gold] at (0.42,0.16) {$\theta$};
\fill[deepblue] (P) circle (0.025) node[above right] {$P(\cos\theta,\sin\theta)$};
\node[below] at ({cos(\ang)/2},0) {$\cos\theta$};
\node[left] at (0,{sin(\ang)/2}) {$\sin\theta$};
\end{tikzpicture}
\end{center}

Comme $P$ est sur le cercle $x^2+y^2=1$, on obtient immédiatement l'identité fondamentale.

\begin{theorem}{Identité pythagoricienne}{}
Pour tout angle réel $\theta$,
\[
\cos^2\theta+\sin^2\theta=1.
\]
\end{theorem}

\section{Angles remarquables}

Les valeurs exactes proviennent des triangles $45^\circ$-$45^\circ$-$90^\circ$ et $30^\circ$-$60^\circ$-$90^\circ$.

\begin{center}
\begin{tabular}{c|ccccc}
\toprule
$\theta$ & $0$ & $\pi/6$ & $\pi/4$ & $\pi/3$ & $\pi/2$ \\
\midrule
$\sin\theta$ & $0$ & $1/2$ & $\sqrt2/2$ & $\sqrt3/2$ & $1$ \\
$\cos\theta$ & $1$ & $\sqrt3/2$ & $\sqrt2/2$ & $1/2$ & $0$ \\
$\tan\theta$ & $0$ & $1/\sqrt3$ & $1$ & $\sqrt3$ & non définie \\
\bottomrule
\end{tabular}
\end{center}

Dans les autres quadrants, les valeurs absolues se retrouvent avec l'angle de référence; les signes proviennent des coordonnées du point sur le cercle.

\begin{examplebox}
L'angle $150^\circ$ a pour angle de référence $30^\circ$ et se trouve dans le deuxième quadrant. Ainsi,
\[
\sin150^\circ=\frac12,
\qquad
\cos150^\circ=-\frac{\sqrt3}{2},
\qquad
\tan150^\circ=-\frac1{\sqrt3}.
\]
\end{examplebox}

\section{Périodicité et symétries}

Un tour complet ramène au même point :
\[
\sin(\theta+2\pi)=\sin\theta,
\qquad
\cos(\theta+2\pi)=\cos\theta.
\]
La tangente, définie par
\[
\tan\theta=\frac{\sin\theta}{\cos\theta},
\]
a une période $\pi$ et n'est pas définie lorsque $\cos\theta=0$.

Les symétries du cercle donnent
\[
\sin(-\theta)=-\sin\theta,
\qquad
\cos(-\theta)=\cos\theta.
\]
Le sinus est impair et le cosinus est pair.


\section{Le radian mesure naturellement une rotation}

\begin{visualbox}
Un degré dépend d'une convention de partage du cercle en $360$ parties. Un radian compare deux longueurs : l'arc parcouru et le rayon. Cette définition explique pourquoi les formules trigonométriques utilisent naturellement $\pi$ et pourquoi, sur le cercle unité, la mesure de l'angle est exactement la longueur de l'arc.
\end{visualbox}

\begin{center}
\begin{tikzpicture}[scale=2.0]
  \draw[slate,thick] (0,0) circle (1);
  \draw[navy,very thick] (0,0)--(1,0);
  \draw[navy,very thick] (0,0)--({cos(57.3)},{sin(57.3)});
  \draw[teal,very thick] (1,0) arc[start angle=0,end angle=57.3,radius=1];
  \node[teal] at (0.85,0.48) {$s=r$};
  \node[navy] at (0.5,-0.13) {$r$};
  \node[navy] at (0.28,0.52) {$1\text{ rad}$};
\end{tikzpicture}
\qquad
\begin{tikzpicture}[scale=2.0]
  \draw[slate,thick] (0,0) circle (1);
  \draw[-{Stealth},slate] (-1.25,0)--(1.3,0);
  \draw[-{Stealth},slate] (0,-1.25)--(0,1.3);
  \coordinate (P) at ({cos(140)},{sin(140)});
  \draw[gold,very thick] (0,0)--(P);
  \draw[dashed,teal] (P)--({cos(140)},0);
  \fill[terracotta] (P) circle (2pt);
  \node[anchor=south east] at (P) {$P(\cos\theta,\sin\theta)$};
\end{tikzpicture}
\end{center}

\begin{connectionbox}
Sur le cercle unité, le cosinus est la projection horizontale du rayon et le sinus sa projection verticale. L'identité
\[
\cos^2\theta+\sin^2\theta=1
\]
n'est donc pas une formule supplémentaire : c'est simplement l'équation $x^2+y^2=1$ du cercle appliquée au point $P$.
\end{connectionbox}

\subsection{Angles de référence et symétries}
Plutôt que mémoriser une longue table, on localise le quadrant, on ramène l'angle à un angle aigu de référence, puis on détermine les signes des coordonnées. Les symétries du cercle expliquent les relations comme $\cos(-\theta)=\cos\theta$ et $\sin(-\theta)=-\sin\theta$.


\section{Pourquoi les radians simplifient les formules}

Pour un angle $\theta$ en radians, la longueur d'arc est
\[
s=r\theta.
\]
Cette formule ne contient aucun facteur de conversion parce que le radian est précisément défini par le rapport $s/r$. En degrés, il faudrait écrire
\[
s=r\theta\frac{\pi}{180}.
\]
Le radian est donc l'unité qui rend la relation entre angle et longueur directement proportionnelle.

\subsection{Une petite rotation produit un petit déplacement}
Sur un cercle de rayon $r$, une petite variation angulaire $\Delta\theta$ produit un déplacement le long du cercle d'environ
\[
\Delta s=r\Delta\theta.
\]
Cette idée est la base de la vitesse angulaire et explique pourquoi les modèles de mouvement circulaire utilisent les radians.

\begin{center}
\begin{tikzpicture}[scale=1.5]
 \draw[slate,thick] (0,0) circle (2);
 \draw[navy,thick] (0,0)--(2,0);
 \draw[navy,thick] (0,0)--({2*cos(25)},{2*sin(25)});
 \draw[teal,very thick] (2,0) arc[start angle=0,end angle=25,radius=2];
 \draw[gold,thick,<->] (0.65,0) arc[start angle=0,end angle=25,radius=0.65];
 \node[gold!75!black] at (0.78,0.2) {$\Delta\theta$};
 \node[teal] at (1.65,0.45) {$\Delta s$};
 \node[navy] at (1,-0.15) {$r$};
\end{tikzpicture}
\end{center}

\subsection{Construire les valeurs remarquables}
Les valeurs de $30^\circ$, $45^\circ$ et $60^\circ$ ne doivent pas être apprises comme une liste sans origine. Elles viennent de deux triangles : un triangle équilatéral coupé en deux et un triangle rectangle isocèle.

\begin{center}
\begin{tikzpicture}[scale=0.95]
 \draw[navy,very thick] (0,0)--(4,0)--(2,{2*sqrt(3)})--cycle;
 \draw[teal,thick,dashed] (2,0)--(2,{2*sqrt(3)});
 \node[below] at (1,0) {$1$}; \node[below] at (3,0) {$1$};
 \node[left] at (1,{sqrt(3)}) {$2$};
 \node[right] at (2.1,{sqrt(3)}) {$\sqrt3$};
 \node[gold!75!black] at (0.65,0.28) {$60^\circ$};
 \node[gold!75!black] at (2.25,0.35) {$30^\circ$};
\end{tikzpicture}
\qquad
\begin{tikzpicture}[scale=0.95]
 \draw[navy,very thick] (0,0)--(3,0)--(3,3)--cycle;
 \draw (2.7,0) rectangle (3,0.3);
 \node[below] at (1.5,0) {$1$}; \node[right] at (3,1.5) {$1$};
 \node[above left] at (1.5,1.5) {$\sqrt2$};
 \node[gold!75!black] at (0.75,0.3) {$45^\circ$};
\end{tikzpicture}
\end{center}

Ces constructions rendent les coordonnées du cercle unité prévisibles et réduisent la mémorisation à deux figures fondamentales.


\section{Mesurer et lire une rotation}

Les degrés partagent un tour en $360$ parts; les radians mesurent directement une longueur d'arc rapportée au rayon. Les deux unités décrivent la même rotation, mais le radian relie naturellement angle, arc, vitesse angulaire et fonctions trigonométriques.

\subsection{Le radian comme rapport de longueurs}

Dans un cercle de rayon $r$, un angle central $\theta$ radians intercepte un arc de longueur
\[
s=r\theta.
\]
Ainsi
\[
\theta=\frac sr.
\]
Le rapport est sans unité, car le numérateur et le dénominateur sont deux longueurs. Sur le cercle unité, $r=1$, donc la mesure de l'angle est numériquement égale à la longueur de l'arc.

Un tour complet possède une longueur $2\pi r$, donc
\[
2\pi\ \text{radians}=360^\circ.
\]
On en déduit
\[
180^\circ=\pi,
\qquad
1^\circ=\frac\pi{180}.
\]

\subsection{Conversions raisonnées}

Pour convertir des degrés en radians, on multiplie par $\pi/180$ :
\[
150^\circ=150\cdot\frac\pi{180}=\frac{5\pi}{6}.
\]
Pour convertir des radians en degrés, on multiplie par $180/\pi$ :
\[
\frac{7\pi}{4}\cdot\frac{180}{\pi}=315^\circ.
\]
Les facteurs de conversion sont égaux à $1$ puisque $\pi$ radians et $180^\circ$ désignent le même angle.

\subsection{Le cercle unité comme système de coordonnées}

À l'angle $\theta$ correspond le point
\[
P(\theta)=(\cos\theta,\sin\theta).
\]
Le cosinus est l'abscisse et le sinus l'ordonnée. Comme $P$ appartient au cercle de rayon $1$,
\[
\cos^2\theta+\sin^2\theta=1.
\]
Cette identité fondamentale est simplement l'équation du cercle
\[
x^2+y^2=1
\]
appliquée aux coordonnées de $P$.

\subsection{Construire les valeurs remarquables}

Les valeurs de $45^\circ$ proviennent d'un triangle rectangle isocèle. Si les deux côtés de l'angle droit valent $1$, l'hypoténuse vaut $\sqrt2$. Après normalisation sur le cercle unité,
\[
\cos45^\circ=\sin45^\circ=\frac1{\sqrt2}=\frac{\sqrt2}{2}.
\]

Pour $30^\circ$ et $60^\circ$, on coupe un triangle équilatéral de côté $2$ en deux. On obtient un triangle rectangle de côtés
\[
1,\quad\sqrt3,\quad2.
\]
Ainsi
\[
\sin30^\circ=\frac12,
\qquad
\cos30^\circ=\frac{\sqrt3}{2},
\]
\[
\sin60^\circ=\frac{\sqrt3}{2},
\qquad
\cos60^\circ=\frac12.
\]
Les valeurs ne sont donc pas une table arbitraire à mémoriser; elles viennent de deux triangles fondamentaux.

\subsection{Angles coterminaux}

Deux angles qui diffèrent d'un nombre entier de tours aboutissent au même point :
\[
\theta+2k\pi,
\qquad k\in\Z.
\]
Ainsi
\[
\cos(\theta+2k\pi)=\cos\theta,
\qquad
\sin(\theta+2k\pi)=\sin\theta.
\]
Pour ramener un angle au tour principal, on ajoute ou soustrait des multiples de $2\pi$.

\begin{examplebox}
L'angle
\[
\frac{17\pi}{6}
\]
diffère de $2\pi=12\pi/6$ de l'angle
\[
\frac{5\pi}{6}.
\]
Les deux ont donc les mêmes valeurs trigonométriques.
\end{examplebox}

\subsection{Angle de référence et signes par quadrant}

L'angle de référence est l'angle aigu formé avec l'axe horizontal. Il permet de réutiliser les valeurs remarquables, puis de déterminer les signes à partir du quadrant :
\begin{center}
\begin{tabular}{c|ccc}
\toprule
Quadrant & $\cos$ & $\sin$ & $\tan$\\\midrule
I & $+$ & $+$ & $+$\\
II & $-$ & $+$ & $-$\\
III & $-$ & $-$ & $+$\\
IV & $+$ & $-$ & $-$\\
\bottomrule
\end{tabular}
\end{center}

Par exemple, $330^\circ$ a pour angle de référence $30^\circ$ et se trouve dans le quatrième quadrant :
\[
\cos330^\circ=\frac{\sqrt3}{2},
\qquad
\sin330^\circ=-\frac12,
\qquad
\tan330^\circ=-\frac1{\sqrt3}.
\]

\subsection{Symétries}

La réflexion par rapport à l'axe horizontal donne
\[
\cos(-\theta)=\cos\theta,
\qquad
\sin(-\theta)=-\sin\theta.
\]
Le cosinus est pair et le sinus impair. Une rotation d'un demi-tour donne
\[
\cos(\theta+\pi)=-\cos\theta,
\qquad
\sin(\theta+\pi)=-\sin\theta.
\]
Ces relations se voient directement sur les coordonnées du cercle.

\subsection{Longueur d'arc et aire d'un secteur}

Lorsque $\theta$ est en radians,
\[
s=r\theta
\]
pour la longueur d'arc, et
\[
A=\frac12r^2\theta
\]
pour l'aire du secteur. La seconde formule vient du rapport
\[
\frac{\theta}{2\pi}
\]
entre l'angle et le tour complet :
\[
A=\frac{\theta}{2\pi}\cdot\pi r^2=\frac12r^2\theta.
\]

\subsection{Degrés ou radians sur la calculatrice}

Une calculatrice en mode degré interprète $\sin(1)$ comme $\sin(1^\circ)$; en mode radian, elle interprète $\sin(1)$ comme $\sin(1\ \text{rad})$. Les valeurs sont très différentes. Avant tout calcul :
\begin{itemize}
\item repérer le symbole $^\circ$ ou les multiples de $\pi$;
\item choisir le mode correspondant;
\item tester une valeur connue, par exemple $\sin30^\circ=1/2$ ou $\sin(\pi/6)=1/2$.
\end{itemize}

\begin{summarybox}
Le radian mesure une rotation par un rapport de longueurs. Le cercle unité transforme l'angle en coordonnées, explique l'identité pythagoricienne, les signes, les symétries et la périodicité des fonctions trigonométriques.
\end{summarybox}


\chapter{Trigonométrie du triangle}

\begin{questionbox}
Comment déterminer une longueur ou un angle inaccessible à partir de quelques mesures disponibles?
\end{questionbox}

\section{Triangles rectangles}

Dans un triangle rectangle, relativement à un angle aigu $\theta$ :
\[
\sin\theta=\frac{\text{opposé}}{\text{hypoténuse}},
\qquad
\cos\theta=\frac{\text{adjacent}}{\text{hypoténuse}},
\qquad
\tan\theta=\frac{\text{opposé}}{\text{adjacent}}.
\]
Ces rapports ne dépendent que de l'angle, car tous les triangles rectangles possédant le même angle aigu sont semblables.

\begin{center}
\begin{tikzpicture}[scale=0.85]
\coordinate (A) at (0,0);
\coordinate (B) at (5,0);
\coordinate (C) at (5,3.2);
\draw[very thick] (A)--(B)--(C)--cycle;
\draw (4.7,0) rectangle (5,0.3);
\draw[gold,thick] (0.75,0) arc[start angle=0,end angle=32.62,radius=0.75];
\node[gold] at (0.95,0.25) {$\theta$};
\node[below] at (2.5,0) {adjacent};
\node[right] at (5,1.6) {opposé};
\node[above left] at (2.4,1.7) {hypoténuse};
\end{tikzpicture}
\end{center}

\begin{examplebox}
Une échelle de $8$ m forme un angle de $68^\circ$ avec le sol. La hauteur atteinte sur le mur est
\[
h=8\sin68^\circ\approx7{,}42\text{ m}.
\]
La distance du pied de l'échelle au mur est
\[
d=8\cos68^\circ\approx3{,}00\text{ m}.
\]
\end{examplebox}

\section{Vue d'ensemble des triangles quelconques}

Dans un triangle non rectangle, deux lois complètent les rapports trigonométriques. Elles ne sont pas deux recettes concurrentes : chacune correspond à une configuration de données différente.

\begin{center}
\begin{tabularx}{\textwidth}{@{}lYY@{}}
\toprule
Structure des données & Relation utile & Idée géométrique\\\midrule
une paire côté-angle opposé connue & $\dfrac a{\sin A}=\dfrac b{\sin B}=\dfrac c{\sin C}$ & une même hauteur vue dans deux triangles rectangles\\[0.7em]
deux côtés et l'angle compris, ou trois côtés & $c^2=a^2+b^2-2ab\cos C$ & correction apportée au théorème de Pythagore\\
\bottomrule
\end{tabularx}
\end{center}

\begin{center}
\begin{tikzpicture}[scale=0.82]
 \coordinate (A) at (0,0); \coordinate (B) at (5.6,0); \coordinate (C) at (1.9,3.2);
 \draw[navy,very thick] (A)--(B)--(C)--cycle;
 \draw[teal,thick,dashed] (C)--(1.9,0);
 \node[teal,right] at (1.9,1.6) {$h$};
 \node[gold!75!black] at (0.55,0.25) {$A$};
 \node[gold!75!black] at (5.05,0.25) {$B$};
 \node[left] at (0.9,1.65) {$b$}; \node[right] at (3.85,1.65) {$a$};
 \node[below] at (2.8,0) {$c$};
\end{tikzpicture}
\end{center}

\begin{warningbox}
Lorsque deux côtés et un angle non compris sont donnés, la loi des sinus peut conduire à zéro, un ou deux triangles. Il faut étudier la géométrie de la configuration et vérifier que la somme des angles reste inférieure à $180^\circ$ avant de retenir une seconde solution.
\end{warningbox}

\section{Choisir la formule à partir de la figure}

\begin{visualbox}
Les formules trigonométriques ne se choisissent pas par mots-clés, mais par la structure des données. Dans un triangle rectangle, les rapports sinus, cosinus et tangente suffisent. Dans un triangle quelconque, la loi des sinus exige une paire côté-angle opposé connue; la loi des cosinus convient lorsqu'on connaît deux côtés et l'angle compris, ou les trois côtés.
\end{visualbox}

\begin{center}
\resizebox{0.82\textwidth}{!}{%
\begin{tikzpicture}[node distance=0.75cm and 1.0cm,every node/.style={font=\small,rounded corners,inner sep=5pt,align=center}]
 \node[draw=navy,fill=mistblue] (start) {Le triangle est-il rectangle?};
 \node[draw=teal,fill=mistteal,below left=of start] (rect) {$\sin,\cos,\tan$};
 \node[draw=gold,fill=mistgold,below right=of start] (pair) {Une paire\\côté-angle opposé?};
 \node[draw=navy,fill=mistblue,below left=of pair] (sin) {loi des sinus};
 \node[draw=terracotta,fill=mistred,below right=of pair] (cos) {loi des cosinus};
 \draw[->,thick] (start)--node[left]{oui}(rect);
 \draw[->,thick] (start)--node[right]{non}(pair);
 \draw[->,thick] (pair)--node[left]{oui}(sin);
 \draw[->,thick] (pair)--node[right]{non}(cos);
\end{tikzpicture}}
\end{center}

\begin{center}
\begin{tikzpicture}[scale=0.8]
 \coordinate (A) at (0,0); \coordinate (B) at (5.5,0); \coordinate (C) at (1.8,3.2);
 \draw[navy,very thick] (A)--(B)--(C)--cycle;
 \draw[teal,thick,dashed] (C)--(1.8,0);
 \node[left] at (0.9,1.65) {$b$}; \node[right] at (3.75,1.7) {$a$}; \node[below] at (2.75,0) {$c$};
 \node[gold] at (0.55,0.25) {$A$}; \node[gold] at (4.95,0.25) {$B$};
 \node[teal] at (2.05,1.5) {$h$};
\end{tikzpicture}
\end{center}

\begin{connectionbox}
La hauteur $h=b\sin A$ rend visible la possibilité de zéro, un ou deux triangles lorsqu'on connaît deux côtés et un angle non compris. La seconde solution de $\sin B=k$ n'est pas une anomalie algébrique : elle correspond à une deuxième position géométrique possible du sommet.
\end{connectionbox}

\subsection{Toujours contrôler le triangle obtenu}
Les trois angles doivent totaliser $180^\circ$, le plus grand côté doit faire face au plus grand angle, et l'inégalité triangulaire doit être respectée. Ces contrôles géométriques sont souvent plus rapides qu'une reprise complète du calcul.


\section{Origine géométrique des lois du sinus et du cosinus}

\subsection{La loi des sinus vient d'une même hauteur}
Dans un triangle $ABC$, traçons la hauteur issue de $C$. Elle peut être exprimée à partir de deux triangles rectangles :
\[
h=b\sin A=a\sin B.
\]
On obtient alors
\[
\frac a{\sin A}=\frac b{\sin B}.
\]
La troisième égalité vient d'une autre hauteur. La loi des sinus compare donc chaque côté à la projection verticale du côté voisin.

\begin{center}
\begin{tikzpicture}[scale=0.9]
 \coordinate (A) at (0,0); \coordinate (B) at (6,0); \coordinate (C) at (2.1,3.5); \coordinate (H) at (2.1,0);
 \draw[navy,very thick] (A)--(B)--(C)--cycle;
 \draw[teal,thick,dashed] (C)--(H);
 \draw (1.85,0) rectangle (2.1,0.25);
 \node[left] at (1.05,1.8) {$b$}; \node[right] at (4.15,1.8) {$a$};
 \node[below] at (3,0) {$c$}; \node[teal,right] at (2.1,1.8) {$h$};
 \node[gold!75!black] at (0.55,0.25) {$A$}; \node[gold!75!black] at (5.45,0.25) {$B$};
\end{tikzpicture}
\end{center}

\subsection{La loi des cosinus mesure le défaut d'angle droit}
Si deux côtés de longueurs $b$ et $c$ forment un angle $A$, le troisième côté correspond à la différence de deux vecteurs. En décomposant le côté $b$ en une composante horizontale $b\cos A$ et une composante verticale $b\sin A$, on obtient
\[
a^2=(c-b\cos A)^2+(b\sin A)^2.
\]
Après développement et utilisation de $\sin^2A+\cos^2A=1$,
\[
a^2=b^2+c^2-2bc\cos A.
\]
Lorsque $A=90^\circ$, le terme correctif disparaît et on retrouve Pythagore.

\begin{center}
\begin{tikzpicture}[scale=0.95]
 \coordinate (A) at (0,0); \coordinate (B) at (5.5,0); \coordinate (C) at (2.2,3.0); \coordinate (H) at (2.2,0);
 \draw[navy,very thick] (A)--(B)--(C)--cycle;
 \draw[teal,thick,dashed] (C)--(H);
 \node[below] at (1.1,0) {$b\cos A$};
 \node[below] at (3.85,0) {$c-b\cos A$};
 \node[teal,right] at (2.2,1.5) {$b\sin A$};
 \node[left] at (1.0,1.5) {$b$}; \node[right] at (3.9,1.6) {$a$};
 \node[gold!75!black] at (0.6,0.25) {$A$};
\end{tikzpicture}
\end{center}

\begin{connectionbox}
La loi des sinus est particulièrement efficace lorsqu'une paire côté-angle opposé est connue. La loi des cosinus est particulièrement efficace lorsqu'on connaît deux côtés et leur angle compris, ou les trois côtés. Ce choix découle de ce que chaque formule rend directement accessible.
\end{connectionbox}



\subsection{Visualiser les deux triangles possibles}
Lorsque deux côtés et un angle non compris sont connus, un sommet peut parfois occuper deux positions. Fixons le côté $AB=c$ et l'angle $A$. Le sommet $C$ doit se trouver sur une demi-droite issue de $A$, mais aussi sur un cercle de centre $B$ et de rayon $a$. Le nombre d'intersections du cercle avec la demi-droite donne le nombre de triangles.

\begin{center}
\begin{tikzpicture}[scale=0.9]
 \coordinate (A) at (0,0); \coordinate (B) at (5.5,0);
 \draw[navy,very thick] (A)--(B);
 \draw[teal,very thick,-{Stealth[length=3mm]}] (A)--(6,3.7);
 \draw[terracotta,thick] (B) circle (3.1);
 \fill[gold] (2.75,1.695) circle (3pt) node[above left] {$C_1$};
 \fill[gold] (4.35,2.68) circle (3pt) node[above right] {$C_2$};
 \draw[gold,thick] (B)--(2.75,1.695);
 \draw[gold,thick] (B)--(4.35,2.68);
 \node[below] at (2.75,0) {$c$};
 \node[terracotta,right] at (7.2,0.8) {cercle : $BC=a$};
 \node[teal] at (1.2,0.45) {angle $A$ fixé};
\end{tikzpicture}
\end{center}

Si le cercle ne rencontre pas la demi-droite, aucun triangle n'existe. S'il est tangent, un seul triangle existe. S'il la coupe en deux points, deux triangles distincts satisfont les données.

\subsection{La fonction arcsinus ne voit qu'une branche principale}
Lorsque $\sin B=k$, la calculatrice retourne un angle principal $B_1=\arcsin k$. Sur $[0,180^\circ]$, l'autre angle ayant le même sinus est
\[
B_2=180^\circ-B_1.
\]
Il faut ensuite vérifier si $A+B_2<180^\circ$. La géométrie du cercle et la symétrie du sinus expliquent cette seconde possibilité; elle ne doit pas être ajoutée mécaniquement sans contrôle.


\section{Résoudre et valider un triangle}

Résoudre un triangle signifie déterminer ses côtés et ses angles à partir d'informations suffisantes. La première étape est toujours un schéma correctement étiqueté : le côté $a$ doit être opposé à l'angle $A$, le côté $b$ à $B$ et le côté $c$ à $C$.

\subsection{Triangles rectangles et similitude}

Les rapports
\[
\sin\theta=\frac{\text{opposé}}{\text{hypoténuse}},
\qquad
\cos\theta=\frac{\text{adjacent}}{\text{hypoténuse}},
\qquad
\tan\theta=\frac{\text{opposé}}{\text{adjacent}}
\]
ne dépendent que de l'angle, car tous les triangles rectangles partageant cet angle sont semblables.

\begin{examplebox}
Une échelle de $8$ m forme un angle de $68^\circ$ avec le sol. La hauteur atteinte est
\[
h=8\sin68^\circ\approx7{,}42\ \mathrm{m}.
\]
La distance du pied au mur est
\[
d=8\cos68^\circ\approx3{,}00\ \mathrm{m}.
\]
Le contrôle de Pythagore confirme que $h^2+d^2\approx8^2$.
\end{examplebox}

\subsection{Démonstration de la loi des sinus}

Dans un triangle $ABC$, traçons la hauteur $h$ issue de $C$ vers le côté $AB$. Dans les deux triangles rectangles obtenus,
\[
h=b\sin A
\]
et
\[
h=a\sin B.
\]
Donc
\[
b\sin A=a\sin B,
\]
d'où
\[
\frac a{\sin A}=\frac b{\sin B}.
\]
En répétant le raisonnement avec une autre hauteur,
\[
\frac a{\sin A}=\frac b{\sin B}=\frac c{\sin C}.
\]
La correspondance côté-angle opposé est la structure centrale de la formule.

\subsection{Utiliser la loi des sinus}

La loi est particulièrement efficace lorsqu'une paire opposée complète est connue.

\begin{examplebox}
Dans un triangle, $A=42^\circ$, $B=71^\circ$ et $a=9$ cm. Alors
\[
C=180^\circ-42^\circ-71^\circ=67^\circ.
\]
Puis
\[
\frac b{\sin71^\circ}=\frac9{\sin42^\circ},
\]
\[
b=9\frac{\sin71^\circ}{\sin42^\circ}\approx12{,}72\ \mathrm{cm}.
\]
De même,
\[
c=9\frac{\sin67^\circ}{\sin42^\circ}\approx12{,}38\ \mathrm{cm}.
\]
\end{examplebox}

\subsection{Analyse algébrique des solutions possibles}

Avec les mêmes notations, la loi des sinus conduit à
\[
\sin B=\frac{b\sin A}{a}=k.
\]
Trois situations doivent être distinguées :
\begin{itemize}
\item si $k>1$, aucune valeur réelle de $B$ n'existe;
\item si $k=1$, alors $B=90^\circ$ et un seul triangle est possible;
\item si $0<k<1$, deux candidats apparaissent :
\[
B_1=\arcsin k,
\qquad
B_2=180^\circ-B_1.
\]
\end{itemize}
Chaque candidat doit ensuite satisfaire $A+B<180^\circ$. La troisième mesure est $C=180^\circ-A-B$, puis le dernier côté se calcule par la loi des sinus.

\begin{examplebox}
Pour $A=30^\circ$, $a=7$ et $b=10$,
\[
k=\frac{10\sin30^\circ}{7}=\frac57.
\]
Ainsi,
\[
B_1\approx45{,}6^\circ,
\qquad
B_2\approx134{,}4^\circ.
\]
Les deux valeurs laissent un troisième angle positif :
\[
C_1\approx104{,}4^\circ,
\qquad
C_2\approx15{,}6^\circ.
\]
Les données déterminent donc deux triangles non congruents. Le contrôle géométrique par la hauteur et l'analyse algébrique aboutissent au même verdict.
\end{examplebox}

\subsection{Dérivation de la loi des cosinus}

Plaçons un triangle avec un côté $a$ sur l'axe horizontal et un côté $b$ formant l'angle $C$. Le troisième côté est la différence de deux vecteurs. Par le produit scalaire,
\begin{align*}
c^2
&=\norm{\mathbf a-\mathbf b}^2\\
&=(\mathbf a-\mathbf b)\cdot(\mathbf a-\mathbf b)\\
&=a^2+b^2-2\mathbf a\cdot\mathbf b\\
&=a^2+b^2-2ab\cos C.
\end{align*}
Lorsque $C=90^\circ$, $\cos C=0$ et on retrouve Pythagore.

\subsection{Utiliser la loi des cosinus}

La loi des cosinus convient à deux configurations :
\begin{itemize}
\item deux côtés et l'angle compris, pour trouver le troisième côté;
\item les trois côtés, pour trouver un angle.
\end{itemize}

\begin{examplebox}
Deux côtés mesurent $7$ cm et $11$ cm et l'angle compris vaut $54^\circ$. Le troisième côté $c$ satisfait
\[
c^2=7^2+11^2-2(7)(11)\cos54^\circ.
\]
Donc
\[
c\approx8{,}91\ \mathrm{cm}.
\]
La longueur est comprise entre $11-7=4$ et $11+7=18$, comme l'exige l'inégalité triangulaire.
\end{examplebox}

\subsection{Aire d'un triangle quelconque}

Si deux côtés $a,b$ et leur angle compris $C$ sont connus,
\[
\mathcal A=\frac12ab\sin C.
\]
En effet, la hauteur relative à la base $a$ vaut $b\sin C$. La formule généralise l'aire $\frac12\times\text{base}\times\text{hauteur}$.

\subsection{Triangulation d'une largeur inaccessible}

Deux points $A$ et $B$ sont séparés de $120$ m sur une rive. Un point $C$ de l'autre côté est observé sous les angles
\[
A=48^\circ,
\qquad B=63^\circ.
\]
Alors
\[
C=69^\circ.
\]
La loi des sinus donne
\[
AC=120\frac{\sin63^\circ}{\sin69^\circ}\approx114{,}54\ \mathrm{m}.
\]
La largeur perpendiculaire à la rive vaut ensuite
\[
AC\sin48^\circ\approx85{,}11\ \mathrm{m}.
\]
Le problème combine une mesure accessible, deux angles et une décomposition en triangle rectangle.

\subsection{Guide de choix}

\begin{center}
\begin{tabularx}{\textwidth}{@{}lY@{}}
\toprule
Données & Méthode naturelle\\\midrule
triangle rectangle & Pythagore et rapports trigonométriques\\
une paire côté-angle opposé & loi des sinus\\
deux côtés et angle compris & loi des cosinus\\
trois côtés & loi des cosinus pour un angle\\
deux côtés et angle non compris & loi des sinus, puis analyse d'existence et d'unicité\\
deux côtés et angle compris, aire cherchée & $\frac12ab\sin C$\\
\bottomrule
\end{tabularx}
\end{center}

\subsection{Contrôles}

La somme des angles doit être $180^\circ$. Le plus grand côté doit être opposé au plus grand angle. L'inégalité triangulaire doit être respectée. Une calculatrice doit être dans le bon mode. Enfin, les côtés et les angles doivent rester associés à leurs opposés dans toutes les formules.

\begin{summarybox}
La trigonométrie du triangle repose sur la similitude, la loi des sinus et la loi des cosinus. Un schéma étiqueté et l'identification de la configuration déterminent la méthode; les contrôles géométriques vérifient le résultat.
\end{summarybox}


\chapter[Fonctions trigonométriques et fonctions réciproques]{Fonctions trigonométriques\\et fonctions réciproques}

\begin{questionbox}
Comment passer du mouvement sur un cercle à des graphes périodiques, puis retrouver un angle à partir d'une valeur de sinus, cosinus ou tangente?
\end{questionbox}

\section{Graphes fondamentaux}

Lorsque le point parcourt le cercle unité, son ordonnée produit la fonction sinus et son abscisse produit la fonction cosinus.

\begin{center}
\begin{tikzpicture}
\begin{axis}[
axis lines=middle,xmin=-0.2,xmax=6.6,ymin=-1.4,ymax=1.4,
xtick={0,1.5708,3.1416,4.7124,6.2832},
xticklabels={$0$,$\pi/2$,$\pi$,$3\pi/2$,$2\pi$},
ytick={-1,0,1},width=12cm,height=6cm,grid=major,grid style={gray!20}]
\addplot[teal,very thick,domain=0:6.2832,samples=180] {sin(deg(x))};
\addplot[gold,very thick,densely dashed,domain=0:6.2832,samples=180] {cos(deg(x))};
\node[teal] at (axis cs:1.3,1.15) {$y=\sin x$};
\node[gold] at (axis cs:5.1,0.75) {$y=\cos x$};
\end{axis}
\end{tikzpicture}
\end{center}

Pour $\sin x$ et $\cos x$ :
\begin{itemize}
\item domaine : $\R$;
\item image : $[-1,1]$;
\item période : $2\pi$.
\end{itemize}
La tangente a pour domaine
\[
\R\setminus\left\{\frac\pi2+k\pi:k\in\Z\right\},
\]
image $\R$ et période $\pi$.

\section{Amplitude, période, déphasage et ligne médiane}

Une fonction sinusoïdale s'écrit
\[
y=A\sin(B(x-C))+D
\]
ou avec le cosinus. Ses paramètres ont un sens géométrique :
\begin{itemize}
\item amplitude : $\abs A$;
\item période : $T=\dfrac{2\pi}{\abs B}$;
\item déphasage horizontal : $C$;
\item ligne médiane : $y=D$;
\item maximum : $D+\abs A$;
\item minimum : $D-\abs A$.
\end{itemize}

\begin{examplebox}
Pour
\[
y=-3\cos\left(2\left(x-\frac\pi4\right)\right)+1,
\]
l'amplitude est $3$, la période est $\pi$, le déphasage est $\pi/4$ vers la droite et la ligne médiane est $y=1$. Les valeurs varient entre $-2$ et $4$.
\end{examplebox}

\section{Fonctions trigonométriques réciproques}

Le sinus et le cosinus n'admettent pas de fonction réciproque sur tout $\R$, car ils sont périodiques. On restreint donc leur domaine à des intervalles où ils sont injectifs.

\begin{definition}{Fonctions trigonométriques réciproques}{}
\begin{align*}
\arcsin&:[-1,1]\to\left[-\frac\pi2,\frac\pi2\right],\\
\arccos&:[-1,1]\to[0,\pi],\\
\arctan&:\R\to\left]-\frac\pi2,\frac\pi2\right[.
\end{align*}
Elles vérifient, sur ces intervalles,
\[
\sin(\arcsin x)=x,
\quad
\cos(\arccos x)=x,
\quad
\tan(\arctan x)=x.
\]
\end{definition}

\begin{warningbox}
La notation $\sin^{-1}x$ est ambiguë : elle désigne souvent $\arcsin x$, alors que $(\sin x)^{-1}$ signifie $1/\sin x$. Ce manuel utilise donc les notations $\arcsin$, $\arccos$ et $\arctan$.
\end{warningbox}

\begin{examplebox}
\[
\arcsin\left(\frac12\right)=\frac\pi6,
\qquad
\arccos\left(-\frac{\sqrt2}{2}\right)=\frac{3\pi}{4},
\qquad
\arctan(1)=\frac\pi4.
\]
Les réponses appartiennent toujours aux intervalles de sortie choisis.
\end{examplebox}

\section{Équations trigonométriques élémentaires}

Pour résoudre $\sin x=a$, on trouve d'abord un angle principal, puis on utilise les symétries et la périodicité.

\begin{examplebox}
Résolvons $\sin x=\frac12$ sur $[0,2\pi]$. L'angle de référence est $\pi/6$. Le sinus est positif dans les quadrants I et II, donc
\[
x=\frac\pi6
\quad\text{ou}\quad
x=\frac{5\pi}{6}.
\]
Sur tout $\R$,
\[
x=\frac\pi6+2k\pi
\quad\text{ou}\quad
x=\frac{5\pi}{6}+2k\pi,
\qquad k\in\Z.
\]
\end{examplebox}

\begin{examplebox}
Résolvons $\cos(2x)=0$ :
\[
2x=\frac\pi2+k\pi,
\]
donc
\[
x=\frac\pi4+\frac{k\pi}{2},
\qquad k\in\Z.
\]
\end{examplebox}


\section{La sinusoïde est l'ombre d'une rotation}

\begin{visualbox}
Lorsque le point $P$ tourne uniformément sur le cercle unité, sa hauteur monte et descend entre $-1$ et $1$. Reporter cette hauteur au fil de l'angle produit le graphe du sinus. La périodicité du graphe est donc la répétition d'un tour complet, et non une propriété arbitraire de la formule.
\end{visualbox}

\begin{center}
\begin{tikzpicture}[scale=0.75]
  \begin{scope}[xshift=-4.1cm]
    \draw[slate,thick] (0,0) circle (1.8);
    \draw[-{Stealth},slate] (-2.2,0)--(2.3,0);
    \draw[-{Stealth},slate] (0,-2.2)--(0,2.3);
    \coordinate (P) at ({1.8*cos(50)},{1.8*sin(50)});
    \draw[navy,very thick] (0,0)--(P);
    \draw[dashed,teal] (P)--({1.8*cos(50)},0);
    \fill[terracotta] (P) circle (2.5pt);
    \node[teal,right] at ({1.8*cos(50)},0.65) {$\sin\theta$};
  \end{scope}
  \begin{scope}[xshift=0cm]
    \draw[-{Stealth},thick,slate] (-0.3,0)--(7.4,0) node[right] {$\theta$};
    \draw[-{Stealth},thick,slate] (0,-2.2)--(0,2.3) node[above] {$y$};
    \draw[navy,very thick,domain=0:6.7,samples=120] plot(\x,{1.8*sin(57.2958*\x)});
    \draw[dashed,teal] (0.8727,{1.8*sin(50)})--(-2.25,{1.8*sin(50)});
    \fill[terracotta] (0.8727,{1.8*sin(50)}) circle (2.5pt);
    \node[below] at (3.14,-2.4) {$2\pi$ radians : un tour};
  \end{scope}
\end{tikzpicture}
\end{center}

\begin{connectionbox}
Dans
\[
y=A\sin(B(x-C))+D,
\]
$D$ déplace la ligne médiane, $|A|$ est la distance maximale à cette ligne, $2\pi/|B|$ est la durée d'un cycle, et $C$ indique où le cycle est placé horizontalement. Chaque paramètre correspond à un geste visible sur le graphe.
\end{connectionbox}

\subsection{Les fonctions inverses récupèrent un angle principal}
La valeur $\arcsin(y)$ ne donne pas tous les angles ayant le même sinus; elle donne l'angle principal dans un intervalle choisi pour rendre le sinus inversible. Les autres solutions viennent des symétries du cercle et de la périodicité.


\section{Effet des paramètres d'une sinusoïde}

Partons de $y=\sin x$. Chaque paramètre de
\[
y=A\sin(B(x-C))+D
\]
modifie une caractéristique précise.

\subsection{Amplitude et ligne médiane}
Multiplier par $A$ étire verticalement les hauteurs. Ajouter $D$ déplace la ligne médiane.

\begin{center}
\begin{tikzpicture}[scale=0.58]
 \draw[-{Stealth},thick,slate] (-0.3,0)--(7,0); \draw[-{Stealth},thick,slate] (0,-2.8)--(0,3.1);
 \draw[navy,very thick,domain=0:6.5,samples=100] plot(\x,{sin(57.3*\x)});
 \draw[teal,very thick,dashed,domain=0:6.5,samples=100] plot(\x,{2*sin(57.3*\x)});
 \node[navy] at (5.5,1.3) {$\sin x$}; \node[teal] at (5.5,2.35) {$2\sin x$};
\end{tikzpicture}
\qquad
\begin{tikzpicture}[scale=0.58]
 \draw[-{Stealth},thick,slate] (-0.3,0)--(7,0); \draw[-{Stealth},thick,slate] (0,-0.8)--(0,4.8);
 \draw[dashed,gold] (0,2)--(6.7,2);
 \draw[navy,very thick,domain=0:6.5,samples=100] plot(\x,{2+sin(57.3*\x)});
 \node[gold!75!black,right] at (6.6,2) {$D=2$};
\end{tikzpicture}
\end{center}

\subsection{Période et déphasage}
Remplacer $x$ par $Bx$ accélère la lecture de l'angle : lorsque $x$ parcourt $2\pi/|B|$, l'argument parcourt un tour complet. Remplacer $x$ par $x-C$ déplace le motif de $C$ unités vers la droite.

\begin{center}
\begin{tikzpicture}[scale=0.58]
 \draw[-{Stealth},thick,slate] (-0.3,0)--(7,0); \draw[-{Stealth},thick,slate] (0,-1.6)--(0,1.9);
 \draw[navy,very thick,domain=0:6.5,samples=120] plot(\x,{sin(114.6*\x)});
 \draw[terracotta,<->,thick] (0.05,1.45)--(3.14,1.45);
 \node[terracotta,fill=white] at (1.6,1.45) {période $\pi$};
\end{tikzpicture}
\qquad
\begin{tikzpicture}[scale=0.58]
 \draw[-{Stealth},thick,slate] (-0.3,0)--(7,0); \draw[-{Stealth},thick,slate] (0,-1.6)--(0,1.9);
 \draw[navy!40,thick,dashed,domain=0:6.5,samples=100] plot(\x,{sin(57.3*\x)});
 \draw[teal,very thick,domain=0:6.5,samples=100] plot(\x,{sin(57.3*(\x-1))});
 \draw[gold,->,thick] (0.1,1.4)--(1.0,1.4);
 \node[gold!75!black,above] at (0.55,1.4) {$C=1$};
\end{tikzpicture}
\end{center}

\subsection{Lire les paramètres à partir d'un graphe}
Si le maximum est $M$ et le minimum $m$, alors
\[
D=\frac{M+m}{2},
\qquad
|A|=\frac{M-m}{2}.
\]
La distance horizontale entre deux maxima successifs donne la période $T$, puis $|B|=2\pi/T$. La phase est déterminée en repérant un point caractéristique : maximum, minimum ou passage par la ligne médiane avec son sens de variation.

\begin{visualbox}
On ne devrait pas chercher simultanément $A,B,C,D$. Lire d'abord la géométrie verticale ($A,D$), puis la répétition horizontale ($B$), et enfin la position du motif ($C$).
\end{visualbox}



\subsection{Retrouver toutes les solutions à partir du cercle}
Pour résoudre $\sin x=k$, la calculatrice fournit l'angle principal $\alpha=\arcsin k$. Sur le cercle, une même hauteur coupe généralement le cercle en deux points. Les angles correspondants sont
\[
x=\alpha+2m\pi
\quad\text{ou}\quad
x=\pi-\alpha+2m\pi,
\qquad m\in\mathbb Z.
\]

\begin{center}
\begin{tikzpicture}[scale=1.8]
 \draw[slate,thick] (0,0) circle (1);
 \draw[-{Stealth},slate] (-1.25,0)--(1.3,0);
 \draw[-{Stealth},slate] (0,-1.25)--(0,1.3);
 \draw[teal,thick,dashed] (-1.1,0.62)--(1.1,0.62) node[right] {$y=k$};
 \fill[terracotta] ({sqrt(1-0.62^2)},0.62) circle (2.3pt);
 \fill[terracotta] ({-sqrt(1-0.62^2)},0.62) circle (2.3pt);
 \draw[gold,thick] (0,0)--({sqrt(1-0.62^2)},0.62);
 \draw[gold,thick] (0,0)--({-sqrt(1-0.62^2)},0.62);
 \node[gold!75!black] at (0.55,0.2) {$\alpha$};
 \node[gold!75!black] at (-0.55,0.2) {$\pi-\alpha$};
\end{tikzpicture}
\end{center}

Pour le cosinus, une même abscisse donne deux points symétriques par rapport à l'axe horizontal. Pour la tangente, les solutions se répètent après $\pi$ plutôt qu'après $2\pi$.

\subsection{Intervalle demandé ou solution générale}
Une réponse sur $[0,2\pi[$ est une liste finie. Une solution générale décrit toutes les répétitions. Il faut lire attentivement la question : une liste d'angles principaux ne répond pas à une demande de solution générale, et une formule avec $m\in\mathbb Z$ peut être inutile lorsqu'un intervalle précis est imposé.


\section{Construire une sinusoïde et retrouver un angle}

Les graphes de sinus et cosinus peuvent être lus comme les coordonnées d'un point qui tourne sur le cercle unité. Cette origine géométrique explique leur amplitude, leur période, leurs symétries et le fait qu'une même valeur soit atteinte plusieurs fois.

\subsection{Du cercle au graphe du sinus}

Lorsque l'angle $x$ augmente, le point
\[
(\cos x,\sin x)
\]
se déplace sur le cercle. Le graphe de $y=\sin x$ enregistre l'ordonnée du point en fonction de l'angle. Aux angles
\[
0,\ \frac\pi2,\ \pi,\ \frac{3\pi}2,\ 2\pi,
\]
les hauteurs sont
\[
0,\ 1,\ 0,\ -1,\ 0.
\]
Après $2\pi$, le point revient à sa position initiale; le motif se répète.

Le cosinus enregistre l'abscisse. Il commence à $1$ lorsque $x=0$, ce qui correspond à un décalage d'un quart de période par rapport au sinus.

\subsection{Lire les paramètres d'une sinusoïde}

Pour
\[
y=A\sin(B(x-C))+D,
\]
les paramètres ont les rôles suivants :
\[
\text{amplitude}=\abs A,
\qquad
\text{ligne médiane}=D,
\]
\[
\text{période}=\frac{2\pi}{\abs B},
\qquad
\text{déphasage}=C.
\]
Le signe de $A$ réfléchit le graphe autour de la ligne médiane.

Si le maximum est $M$ et le minimum $m$, alors
\[
D=\frac{M+m}{2},
\qquad
\abs A=\frac{M-m}{2}.
\]

\begin{examplebox}
Une fonction a un maximum $9$, un minimum $1$ et une période $6$. Alors
\[
D=5,
\qquad \abs A=4,
\qquad \abs B=\frac{2\pi}{6}=\frac\pi3.
\]
Si elle commence à son maximum au temps $0$, un modèle naturel est
\[
y=4\cos\left(\frac\pi3t\right)+5.
\]
\end{examplebox}

\subsection{Dérivation de la période \texorpdfstring{$2\pi/\abs B$}{2 pi sur valeur absolue de B}}

La fonction $\sin u$ se répète lorsque son argument augmente de $2\pi$. Pour
\[
f(x)=\sin(Bx),
\]
la période $T$ satisfait
\[
B(x+T)-Bx=2\pi.
\]
Donc
\[
BT=2\pi,
\qquad
T=\frac{2\pi}{\abs B}.
\]
La valeur absolue garantit une période positive.

\subsection{Choisir sinus ou cosinus}

Le choix n'est pas unique, mais certaines formes sont plus naturelles :
\begin{itemize}
\item départ sur la ligne médiane avec montée : sinus positif;
\item départ sur la ligne médiane avec descente : sinus négatif;
\item départ au maximum : cosinus positif;
\item départ au minimum : cosinus négatif.
\end{itemize}
Un déphasage permet de convertir toute forme sinus en forme cosinus et inversement.

\subsection{Fonctions trigonométriques inverses}

Le sinus n'est pas injectif sur $\R$, car il répète ses valeurs. Pour définir $\arcsin$, on restreint le sinus à
\[
\left[-\frac\pi2,\frac\pi2\right].
\]
Ainsi
\[
\arcsin:[-1,1]\to\left[-\frac\pi2,\frac\pi2\right].
\]
De même,
\[
\arccos:[-1,1]\to[0,\pi],
\]

a et
\[
\arctan:\R\to\left]-\frac\pi2,\frac\pi2\right[.
\]
La calculatrice retourne toujours la valeur principale dans ces intervalles.

\begin{warningbox}
L'égalité
\[
\arcsin(\sin x)=x
\]
n'est vraie que lorsque $x$ appartient à l'intervalle principal de $\arcsin$. Par exemple,
\[
\arcsin(\sin(5\pi/6))=\arcsin(1/2)=\pi/6,
\]
non $5\pi/6$.
\end{warningbox}

\subsection{Résoudre une équation sur un intervalle}

Résolvons
\[
\sin x=\frac12
\]
sur $[0,2\pi]$. L'angle de référence est $\pi/6$. Le sinus est positif dans les quadrants I et II, donc
\[
x=\frac\pi6
\quad\text{ou}\quad
x=\frac{5\pi}{6}.
\]
La calculatrice fournit seulement $\pi/6$; le cercle fournit la seconde solution.

Pour
\[
\cos x=-\frac{\sqrt2}{2},
\]
le cosinus est négatif dans les quadrants II et III :
\[
x=\frac{3\pi}{4}
\quad\text{ou}\quad
x=\frac{5\pi}{4}.
\]

\subsection{Solutions générales et périodicité}

Les solutions de
\[
\sin x=\sin\alpha
\]
s'écrivent
\[
x=\alpha+2k\pi
\quad\text{ou}\quad
x=\pi-\alpha+2k\pi,
\qquad k\in\Z.
\]
Pour le cosinus,
\[
\cos x=\cos\alpha
\]
donne
\[
x=\pm\alpha+2k\pi.
\]
Pour la tangente, dont la période est $\pi$,
\[
\tan x=\tan\alpha
\]
donne
\[
x=\alpha+k\pi.
\]

\subsection{Étude complète d'un graphe}

Considérons
\[
f(x)=-3\sin(2(x-\pi/4))+1.
\]
On lit :
\begin{itemize}
\item amplitude $3$;
\item ligne médiane $y=1$;
\item période $\pi$;
\item déphasage de $\pi/4$ vers la droite;
\item réflexion autour de la ligne médiane à cause du signe négatif;
\item maximum $4$ et minimum $-2$.
\end{itemize}
Pour tracer une période, on la divise en quatre quarts et on place les cinq points structurants.

\subsection{Erreurs fréquentes}

\begin{itemize}
\item Confondre amplitude et maximum : le maximum est $D+\abs A$.
\item Oublier que $B$ agit inversement sur la période.
\item Lire $C$ sans factoriser $B$ : dans $\sin(2x-\pi)$, on écrit $\sin(2(x-\pi/2))$.
\item Accepter seulement la valeur principale fournie par la calculatrice.
\item Mélanger degrés et radians.
\end{itemize}

\begin{summarybox}
Les graphes trigonométriques proviennent du cercle unité. Les paramètres contrôlent amplitude, ligne médiane, période et phase. Les fonctions inverses donnent un angle principal; les symétries et la périodicité fournissent les autres solutions.
\end{summarybox}


\chapter{Mouvement circulaire et modèles sinusoïdaux}

\begin{questionbox}
Comment construire une formule sinusoïdale à partir d'un phénomène périodique réel?
\end{questionbox}

\section{Vitesse angulaire}

Si un objet effectue un tour complet en une période $T$, sa vitesse angulaire est
\[
\omega=\frac{2\pi}{T}.
\]
L'angle après un temps $t$ est
\[
\theta(t)=\omega t+\theta_0,
\]
où $\theta_0$ est l'angle initial.

\begin{examplebox}
Une roue effectue un tour en $8$ secondes. Alors
\[
\omega=\frac{2\pi}{8}=\frac\pi4\text{ rad/s}.
\]
Après $3$ secondes, elle a tourné de $3\pi/4$ radians, plus son angle initial éventuel.
\end{examplebox}

\section{Coordonnées d'un mouvement circulaire}

Pour un cercle de centre $(h,k)$ et de rayon $R$,
\begin{align*}
x(t)&=h+R\cos(\omega t+\theta_0),\\
y(t)&=k+R\sin(\omega t+\theta_0).
\end{align*}
La coordonnée horizontale et la coordonnée verticale sont donc sinusoïdales.

\begin{examplebox}
Une grande roue de rayon $18$ m a son centre à $20$ m du sol et fait un tour en $90$ s. Si une cabine part du point le plus bas, sa hauteur peut être modélisée par
\[
h(t)=20-18\cos\left(\frac{2\pi}{90}t\right).
\]
Au départ, $h(0)=2$ m. Après $45$ s, la cabine est au sommet, à $38$ m.
\end{examplebox}

\section{Construire un modèle à partir de données}

\begin{methodbox}
Pour construire $y=A\sin(B(t-C))+D$ ou un modèle cosinus :
\begin{enumerate}
\item trouver le maximum $M$ et le minimum $m$;
\item calculer l'amplitude $\abs A=(M-m)/2$;
\item calculer la ligne médiane $D=(M+m)/2$;
\item mesurer la période $T$ et poser $B=2\pi/T$;
\item choisir sinus ou cosinus selon le point de départ;
\item ajuster le déphasage $C$ et le signe de $A$.
\end{enumerate}
\end{methodbox}

\begin{examplebox}
Une température varie entre $12^\circ$C et $24^\circ$C avec une période de $24$ heures. Le maximum est atteint à $15$ h. L'amplitude vaut $6$, la ligne médiane vaut $18$ et $B=2\pi/24=\pi/12$. Un modèle naturel est
\[
T(t)=18+6\cos\left(\frac\pi{12}(t-15)\right).
\]
\end{examplebox}

\section{Interpréter un modèle}

\begin{examplebox}
Pour
\[
y(t)=4\sin\left(\frac\pi3(t-1)\right)+7,
\]
on lit :
\begin{itemize}
\item amplitude $4$;
\item période $T=2\pi/(\pi/3)=6$;
\item ligne médiane $y=7$;
\item déphasage de $1$ unité vers la droite;
\item valeurs comprises entre $3$ et $11$.
\end{itemize}
\end{examplebox}

\section{Interprétation ondulatoire}

Une onde sinusoïdale peut s'écrire
\[
y(t)=A\sin(2\pi ft+\varphi),
\]
où $f$ est la fréquence en cycles par unité de temps et $\varphi$ est la phase initiale. La période est $T=1/f$.

\begin{intuitionbox}
L'amplitude décrit la taille de l'oscillation; la fréquence décrit sa rapidité. Doubler la fréquence ne double pas l'amplitude : ces paramètres contrôlent des aspects indépendants du signal.
\end{intuitionbox}


\section{Du mouvement circulaire au modèle temporel}

\begin{visualbox}
Une grande roue ne « suit » pas une courbe sinusoïdale dans l'espace : la cabine se déplace sur un cercle. C'est sa hauteur, observée en fonction du temps, qui devient sinusoïdale. Le modèle relie donc trois objets : une rotation, une projection verticale et une horloge.
\end{visualbox}

\begin{center}
\begin{tikzpicture}[scale=0.8]
  \draw[slate,thick] (0,0) circle (2.2);
  \draw[slate,thick] (-2.2,-2.5)--(2.2,-2.5);
  \coordinate (P) at ({2.2*cos(135)},{2.2*sin(135)});
  \draw[navy,very thick] (0,0)--(P);
  \draw[dashed,teal] (P)--({2.2*cos(135)},-2.5);
  \fill[terracotta] (P) circle (3pt);
  \draw[teal,very thick,<->] (2.7,-2.5)--(2.7,{2.2*sin(135)});
  \node[teal,right] at (2.7,-0.45) {hauteur $h(t)$};
  \node[navy] at (-0.65,1.0) {$R$};
  \node[gold!75!black] at (0,-2.85) {sol};
\end{tikzpicture}
\qquad
\begin{tikzpicture}[scale=0.63]
  \draw[-{Stealth},thick,slate] (-0.3,0)--(8.2,0) node[right] {$t$};
  \draw[-{Stealth},thick,slate] (0,-0.5)--(0,5.2) node[above] {$h$};
  \draw[dashed,gold] (0,2.4)--(8,2.4) node[right] {$D$};
  \draw[navy,very thick,domain=0:7.8,samples=120] plot(\x,{2.4+1.7*cos(46.15*\x+180)});
  \draw[teal,<->,thick] (1.95,0.7)--(1.95,2.4);
  \node[teal,left] at (1.95,1.55) {$R$};
  \draw[terracotta,<->,thick] (0.05,4.55)--(7.85,4.55);
  \node[terracotta,fill=white] at (3.95,4.55) {une période $T$};
\end{tikzpicture}
\end{center}

\begin{connectionbox}
Le modèle général
\[
h(t)=D+A\cos(\omega t+\varphi)
\]
se lit physiquement : $D$ est la hauteur du centre, $|A|$ le rayon, $\omega=2\pi/T$ la vitesse angulaire et $\varphi$ la position initiale. Le signe de $A$ et la phase ne sont pas deux détails séparés : ils codent ensemble le point de départ et le sens dans lequel on lit le cycle.
\end{connectionbox}

\subsection{Construire un modèle avec peu de données}
Le maximum et le minimum donnent d'abord $D=(M+m)/2$ et $|A|=(M-m)/2$. La période donne ensuite $\omega$. On choisit enfin sinus ou cosinus, puis la phase, pour faire coïncider la position initiale et le sens de variation. Cet ordre évite de deviner quatre paramètres simultanément.


\section{Construire un modèle périodique à partir d'observations}

Supposons qu'une hauteur varie entre $4$ m et $28$ m, avec une période de $60$ s. Elle atteint son minimum à $t=0$ et commence ensuite à augmenter.

\subsection{Étape 1 : centre et amplitude}
Le centre du mouvement est
\[
D=\frac{28+4}{2}=16,
\]
et l'amplitude est
\[
|A|=\frac{28-4}{2}=12.
\]

\subsection{Étape 2 : vitesse angulaire}
Un cycle complet en $60$ s donne
\[
\omega=\frac{2\pi}{60}=\frac\pi{30}.
\]

\subsection{Étape 3 : position initiale}
Le mouvement part du minimum. Le cosinus vaut $1$ à l'origine; on choisit donc un coefficient négatif :
\[
h(t)=16-12\cos\left(\frac\pi{30}t\right).
\]
On vérifie
\[
h(0)=4,
\qquad
h(30)=28,
\qquad
h(60)=4.
\]

\begin{center}
\begin{tikzpicture}[scale=0.095]
 \draw[-{Stealth},thick,slate] (-3,0)--(130,0) node[right] {$t$};
 \draw[-{Stealth},thick,slate] (0,-2)--(0,34) node[above] {$h$};
 \draw[dashed,gold] (0,16)--(125,16);
 \draw[navy,very thick,domain=0:120,samples=180] plot(\x,{16-12*cos(6*\x)});
 \foreach \x/\y/\lab in {0/4/$4$,30/28/$28$,60/4/$4$,90/28/$28$,120/4/$4$}{\fill[terracotta] (\x,\y) circle (14pt) node[above right] {\lab};}
 \draw[teal,<->,very thick] (0,31)--(60,31);
 \node[teal,fill=white] at (30,31) {$T=60$};
\end{tikzpicture}
\end{center}

\subsection{Validation et limites}
Un modèle périodique doit reproduire les extrêmes, la période et la position initiale. Il doit aussi être utilisé dans un contexte où le phénomène se répète de manière suffisamment stable. Des variations de vitesse, du vent ou une déformation mécanique peuvent rendre le modèle seulement approximatif.

\begin{connectionbox}
La formule est la dernière étape. Le raisonnement commence par une image du mouvement, puis identifie successivement la ligne médiane, l'amplitude, la période et la phase. Cette démarche est plus sûre que d'essayer de reconnaître directement une formule complète.
\end{connectionbox}


\section{Construire et valider un modèle périodique}

Un mouvement circulaire uniforme produit des coordonnées sinusoïdales. Cette relation permet de modéliser une grande roue, une marée, une vibration ou toute quantité qui oscille régulièrement autour d'une valeur moyenne.

\subsection{Vitesse angulaire}

Si un tour complet prend une période $T$, le point parcourt $2\pi$ radians en $T$ unités de temps. Sa vitesse angulaire est
\[
\omega=\frac{2\pi}{T}.
\]
Avec une position initiale $\theta_0$,
\[
\theta(t)=\theta_0+\omega t.
\]
Le signe de $\omega$ indique le sens de rotation selon la convention choisie.

\subsection{Coordonnées d'un point en rotation}

Pour un cercle de rayon $R$ centré en $(h,k)$,
\[
x(t)=h+R\cos(\theta_0+\omega t),
\]
\[
y(t)=k+R\sin(\theta_0+\omega t).
\]
Le rayon devient l'amplitude, les coordonnées du centre deviennent les lignes médianes et la vitesse angulaire contrôle la période.

\subsection{Grande roue : modèle complet}

Une grande roue a un rayon de $18$ m, son centre est à $20$ m du sol et elle effectue un tour en $90$ s. Une cabine part du point le plus bas.

La hauteur minimale est $2$ m et la hauteur maximale $38$ m. La vitesse angulaire est
\[
\omega=\frac{2\pi}{90}=\frac\pi{45}.
\]
Comme la cabine part au minimum, un modèle naturel est
\[
H(t)=20-18\cos\left(\frac\pi{45}t\right).
\]
On vérifie :
\[
H(0)=2,
\qquad
H(45)=38,
\qquad
H(90)=2.
\]
La moitié d'une période conduit au sommet, et une période complète ramène au point initial.

Pour savoir quand la cabine atteint $30$ m,
\[
20-18\cos\left(\frac\pi{45}t\right)=30,
\]
\[
\cos\left(\frac\pi{45}t\right)=-\frac59.
\]
Sur un tour, il existe deux instants : un pendant la montée et un pendant la descente.

\subsection{Construire un modèle à partir d'un maximum et d'un minimum}

Si une quantité varie entre $M$ et $m$,
\[
D=\frac{M+m}{2},
\qquad
A=\frac{M-m}{2}.
\]
La période observée fournit
\[
B=\frac{2\pi}{T}.
\]
Il reste à choisir la phase à partir de la valeur et du sens de variation au temps initial.

\begin{examplebox}
Une marée varie entre $1{,}2$ m et $5{,}8$ m avec une période de $12{,}4$ h. Une marée haute se produit à $t=0$. On a
\[
D=\frac{5{,}8+1{,}2}{2}=3{,}5,
\qquad
A=\frac{5{,}8-1{,}2}{2}=2{,}3,
\]
\[
B=\frac{2\pi}{12{,}4}.
\]
Un modèle est
\[
H(t)=3{,}5+2{,}3\cos\left(\frac{2\pi}{12{,}4}t\right).
\]
\end{examplebox}

\subsection{Interpréter chaque paramètre dans le contexte}

Pour
\[
y(t)=A\sin(B(t-C))+D,
\]
\begin{itemize}
\item $D$ est la valeur d'équilibre ou moyenne géométrique du modèle;
\item $\abs A$ est l'écart maximal à cette valeur;
\item $T=2\pi/\abs B$ est le temps d'un cycle;
\item $C$ localise un événement de référence dans le temps;
\item le signe de $A$ indique l'orientation initiale par rapport au modèle de base.
\end{itemize}
Les unités doivent être indiquées : $B$ s'exprime en radians par unité de temps.

\subsection{Fréquence et période}

La fréquence $f$ est le nombre de cycles par unité de temps :
\[
f=\frac1T.
\]
Si l'argument est écrit sous la forme $2\pi ft$, le lien avec la période est immédiat :
\[
\sin(2\pi ft).
\]
Une fréquence plus grande signifie davantage d'oscillations dans le même intervalle et donc une période plus courte.

\subsection{Modèles ondulatoires élémentaires}

Une vibration idéale peut être modélisée par
\[
y(t)=A\sin(2\pi ft+\varphi).
\]
L'amplitude décrit l'intensité géométrique de l'oscillation, la fréquence le nombre de cycles et $\varphi$ la phase initiale. Deux ondes de même fréquence mais de phases différentes atteignent leurs maxima à des moments différents.

\subsection{Lire un graphique périodique}

Pour reconstruire une formule à partir d'un graphe :
\begin{enumerate}
\item lire le maximum et le minimum;
\item calculer la ligne médiane et l'amplitude;
\item mesurer la distance entre deux maxima consécutifs pour obtenir la période;
\item choisir un point caractéristique : maximum, minimum ou passage de la ligne médiane;
\item déterminer si la courbe monte ou descend à ce point;
\item écrire le modèle puis le vérifier sur plusieurs points.
\end{enumerate}

\subsection{Prédiction et répétition}

La périodicité permet de transporter une observation d'un cycle à un autre. Si $T$ est la période,
\[
y(t+T)=y(t).
\]
Ainsi, connaître la valeur à $t=2{,}5$ h donne automatiquement la valeur à $t=2{,}5+kT$ pour tout entier $k$, tant que le modèle reste valable.

\subsection{Limites du modèle sinusoïdal}

Un phénomène réel n'est jamais parfaitement périodique. Une marée subit plusieurs influences; une vibration s'amortit; une grande roue peut accélérer ou s'arrêter. Le modèle sinusoïdal est pertinent lorsque l'amplitude, la période et la ligne médiane restent approximativement stables sur l'intervalle étudié.

Un bon modèle doit donc annoncer son domaine temporel et la précision attendue. Une formule élégante n'est pas une garantie de validité hors du contexte observé.

\subsection{Contrôles}

\begin{itemize}
\item La valeur du modèle reste entre $D-\abs A$ et $D+\abs A$.
\item La période calculée doit correspondre à la répétition du phénomène.
\item La valeur initiale et le sens initial de variation doivent être corrects.
\item Les unités de $t$, $T$, $f$ et $B$ doivent être compatibles.
\item Une équation de hauteur peut avoir deux solutions dans un cycle.
\end{itemize}

\begin{summarybox}
Le mouvement circulaire uniforme produit des coordonnées sinusoïdales. Le rayon devient l'amplitude, le centre devient la ligne médiane, la vitesse angulaire fixe la période et la position initiale fixe la phase. Le modèle doit toujours être vérifié et limité au contexte où la périodicité est plausible.
\end{summarybox}



\backmatter
\chapter*{Synthèse des formules essentielles}
\addcontentsline{toc}{chapter}{Synthèse des formules essentielles}

\section*{Algèbre et fonctions}
\begin{align*}
(a+b)^2&=a^2+2ab+b^2,
&a^2-b^2&=(a-b)(a+b),\\
x_{1,2}&=\frac{-b\pm\sqrt{b^2-4ac}}{2a},
&\log_bx&=\frac{\ln x}{\ln b}.
\end{align*}

\section*{Dénombrement et probabilités}
\begin{align*}
P_n&=n!,&
A_n^k&=\frac{n!}{(n-k)!},\\
\binom nk&=\frac{n!}{k!(n-k)!},&
P(A^c)&=1-P(A),\\
P(A\mid B)&=\frac{P(A\cap B)}{P(B)},&
P(A\cup B)&=P(A)+P(B)-P(A\cap B).
\end{align*}

\section*{Vecteurs et matrices}
\begin{align*}
\norm{\mathbf v}&=\sqrt{v_1^2+\cdots+v_n^2},
&\mathbf u\cdot\mathbf v&=\norm{\mathbf u}\norm{\mathbf v}\cos\theta,\\
\det\begin{pmatrix}a&b\\c&d\end{pmatrix}&=ad-bc,
&\begin{pmatrix}a&b\\c&d\end{pmatrix}^{-1}
&=\frac1{ad-bc}\begin{pmatrix}d&-b\\-c&a\end{pmatrix}.
\end{align*}

\section*{Trigonométrie}
\begin{align*}
\sin^2\theta+\cos^2\theta&=1,
&\tan\theta&=\frac{\sin\theta}{\cos\theta},\\
\frac{a}{\sin A}&=\frac{b}{\sin B}=\frac{c}{\sin C},
&c^2&=a^2+b^2-2ab\cos C,\\
T&=\frac{2\pi}{\abs B},
&\omega&=\frac{2\pi}{T}.
\end{align*}

\chapter*{Conclusion}
\addcontentsline{toc}{chapter}{Conclusion}
Les chapitres de MAT0339 sont reliés par une même idée : choisir une représentation adaptée. Les nombres donnent les objets de base; l'algèbre transforme les expressions; les fonctions décrivent la dépendance; les probabilités modélisent l'incertitude; les vecteurs et matrices organisent l'espace et les systèmes; la trigonométrie décrit les angles, les triangles et les phénomènes périodiques.

Le manuel a volontairement approfondi ces contenus à l'intérieur de leurs chapitres. Les exemples détaillés, les démonstrations, les erreurs fréquentes et les méthodes de vérification ne forment pas un programme supplémentaire : ils rendent le programme principal plus compréhensible et plus autonome.

\end{document}
